\documentclass[11pt,oneside]{book}
\usepackage[margin=1.2in]{geometry}
\usepackage[toc,page]{appendix}
\usepackage{graphicx}
\usepackage{natbib}
\usepackage{lipsum}
\usepackage{caption}
\usepackage[T1]{fontenc}
\usepackage{titlesec, blindtext, color}
\usepackage{xcolor,tikz}
\usetikzlibrary{patterns}
\usetikzlibrary{decorations.markings}
\usepackage{amsmath,amssymb,amsthm,mathrsfs,amsfonts,xfrac,pifont,bbold,physics}
\usepackage[utf8]{inputenc}
\usepackage{amsthm}
\usepackage[breakable, theorems, skins]{tcolorbox}
\usepackage[colorlinks = true,
            linkcolor = red,
            urlcolor  = blue,
            citecolor = red,
            anchorcolor = red]{hyperref}
\usepackage{cleveref}

\usepackage{soul}
\usepackage{frcursive}
\usepackage{booktabs}
\usepackage{enumitem}


\usepackage{array}
\newcolumntype{C}[1]{>{\centering\arraybackslash}m{#1}}

% -------------------------------------------------------------------
% Theorem Styles
% -------------------------------------------------------------------

\theoremstyle{definition} % Define theorem styles here based on the definition style (used for definitions and examples)
\newtheorem*{definition}{Definition}

\theoremstyle{plain} % Define theorem styles here based on the plain style (used for theorems, lemmas, propositions)
\newtheorem{theorem}{Theorem}[section]
\newtheorem{axiom}{Axiom}
\newtheorem{corollary}[theorem]{Corollary}
\newtheorem{lemma}[theorem]{Lemma}
\newtheorem{proposition}[theorem]{Proposition}

\theoremstyle{remark} % Define theorem styles here based on the remark style (used for remarks and notes)
\newtheorem*{notation}{Notation}
\newtheorem*{solution}{Solution}

\newtheoremstyle{underline}% name
{}        % Space above, empty = `usual value'
{}              % Space below
{}              % Body font
{}    % Indent amount (empty = no indent, \parindent = para indent)
{}              % Thm head font
{.}             % Punctuation after thm head
{1.5mm}         % Space after thm head: \newline = linebreak
{{\underline{\textit{\thmname{#1}\thmnumber{ #2}}~\thmnote{(#3)}\unskip}}}  % Thm head spec

\theoremstyle{underline}

\newtheorem{remark}[theorem]{Remark}
\newtheorem{adremark}[theorem]{Additional Remark}
\newtheorem{example}[theorem]{Example}
\newtheorem{claim}[theorem]{Claim}
% -------------------------------------------------------------------
% Chapter Headings
% -------------------------------------------------------------------

\setcounter{chapter}{-1}

\makeatletter
%\renewcommand{\@chapapp}{Lecture}
\makeatother
\definecolor{lightergray}{rgb}{0.9,0.9,0.9}

\usepackage{titlesec}
\titleformat{\section}{\large\bfseries\raggedright}{}{0em}{\colorsection}[\titlerule]
\titleformat{name=\section,numberless}{\large\scshape\bfseries\raggedright}{}{0em}{\colorsectionnonumber}[\titlerule]

\titleformat{\subsection}{\bfseries\raggedright}{}{0em}{\colorsubsection}
\titleformat{name=\subsection,numberless}{\bfseries\raggedright}{}{0em}{\colorsubsectionnonumber}

\newcommand{\colorsection}[1]{%
    \colorbox{lightergray}{\parbox{\dimexpr\textwidth-2\fboxsep}{\thesection\ \ #1}}}
\newcommand{\colorsectionnonumber}[1]{%
    \colorbox{lightergray}{\parbox{\dimexpr\textwidth-2\fboxsep}{#1}}}
    
\newcommand{\colorsubsection}[1]{%
    \colorbox{lightergray}{\parbox{\dimexpr\textwidth-2\fboxsep}{\thesubsection\ #1}}}
\newcommand{\colorsubsectionnonumber}[1]{%
    \colorbox{lightergray}{\parbox{\dimexpr\textwidth-2\fboxsep}{#1}}}
    
\definecolor{gray75}{gray}{0.75}
\newcommand{\hsp}{\hspace{20pt}}
\titleformat{\chapter}[hang]{\Huge\bfseries}{\thechapter\hsp\textcolor{gray75}{|}\hsp}{0pt}{\Huge\bfseries}

% Begin/end Shortcuts
\def\bc{\begin{corollary}}
\def\ec{\end{corollary}}
\def\bcl{\begin{claim}}
\def\ecl{\end{claim}}
\def\bd{\begin{definition}}
\def\ed{\end{definition}}
\def\ben{\begin{enumerate}}
\def\een{\end{enumerate}}
\def\be{\begin{equation}}
\def\ee{\end{equation}}
\def\bse{\begin{equation*}}
\def\ese{\end{equation*}}
\def\bex{\begin{example}}
\def\eex{\end{example}}
\def\bit{\begin{itemize}}
\def\eit{\end{itemize}}
\def\bl{\begin{lemma}}
\def\el{\end{lemma}}
\def\bnn{\begin{notation}}
\def\enn{\end{notation}}
\def\bn{\begin{note}}
\def\en{\end{note}}
\def\bp{\begin{proposition}}
\def\ep{\end{proposition}}
\def\bq{\begin{proof}}
\def\eq{\end{proof}}
\def\br{\begin{remark}}
\def\er{\end{remark}}
\def\badr{\begin{adremark}}
\def\eadr{\end{adremark}}
\def\bs{\begin{solution}}
\def\es{\end{solution}}
\def\btab{\begin{table}}
\def\etab{\end{table}}
\def\btb{\begin{tabular}}
\def\etb{\end{tabular}}
\def\bt{\begin{theorem}}
\def\et{\end{theorem}}
\def\btik{\begin{tikzpicture}}
\def\etik{\end{tikzpicture}}

% Letter Shortcuts
% Greek letters
\def\a{\alpha}
\def\b{\beta}
\def\del{\delta}
\def\g{\gamma}
\def\l{\lambda}
\def\n{\nabla}
\def\sig{\sigma}


% Mathbb
\def\b1{\mathbb{1}} % identity
\def\C{\mathbb{C}} % complex numbers
\def\F{\mathbb{F}}
\def\N{\mathbb{N}} % Natural
\def\Q{\mathbb{Q}} % Rational
\def\R{\mathbb{R}} % Reals
\def\Z{\mathbb{Z}} % Integers

% Mathcal
\def\cA{\mathcal{A}}
\def\cB{\mathcal{B}}
\def\cD{\mathcal{D}}
\def\cF{\mathcal{F}}
\def\cH{\mathcal{H}}
\def\cJ{\mathcal{J}}
\def\cK{\mathcal{K}}
\def\cL{\mathcal{L}}
\def\cM{\mathcal{M}}
\def\cN{\mathcal{N}}
\def\cO{\mathcal{O}}
\def\cP{\mathcal{P}}
\def\cT{\mathcal{T}}

% Mathfrak 
\DeclareMathOperator{\GL}{GL}
\def\gl{\mathfrak{gl}}
\DeclareMathOperator{\Ort}{O}
\def\ort{\mathfrak{o}}
\DeclareMathOperator{\SL}{SL}
\def\sl{\mathfrak{sl}}
\DeclareMathOperator{\SO}{SO}
\def\so{\mathfrak{so}}
\DeclareMathOperator{\SU}{SU}
\def\su{\mathfrak{su}}
\def\fF{\mathfrak{F}}

\def\md{\mathrm{d}}
% Misc. 
\def\diag{\text{diag}}
\def\p{\partial}
\def\ss{\subset}
\def\se{\subseteq}
\def\sm{\setminus}
\def\ve{\varepsilon}
\def\vn{\varnothing}
\def\cl{\colon}
\def\la{\langle}
\def\ra{\rangle}

% Boxes 
\DeclareRobustCommand{\mybox}[2][colframe=blue!10!black,before skip=10pt]{%
\begin{tcolorbox}[
        breakable,
        width=\dimexpr\textwidth\relax, 
        boxsep=5pt,
        arc=5pt,outer arc=5pt,
        ]
        #2
\end{tcolorbox}
}
\def\bbox{\begin{tcolorbox}[colback=white!5,colframe=blue!75!black,title=Exercise]}
\def\ebox{\end{tcolorbox}}

% Feynman Diagram Stuff 
\def\midarrow{\draw[thick, decoration={markings, mark=at position 0.5 with {\arrow{>}}}, postaction={decorate}]}
\def\midarrowblue{\draw[thick, blue, decoration={markings, mark=at position 0.5 with {\arrow{>}}}, postaction={decorate}]}
\def\midarrowred{\draw[thick, red, decoration={markings, mark=at position 0.5 with {\arrow{>}}}, postaction={decorate}]}
\def\aftermidarrow{\draw[thick, decoration={markings, mark=at position 0.7 with {\arrow{>}}}, postaction={decorate}]}
\def\aftermidarrowblue{\draw[thick, blue, decoration={markings, mark=at position 0.7 with {\arrow{>}}}, postaction={decorate}]}
\def\aftermidarrowred{\draw[thick, red, decoration={markings, mark=at position 0.7 with {\arrow{>}}}, postaction={decorate}]}
\def\beforemidarrow{\draw[thick, decoration={markings, mark=at position 0.4 with {\arrow{>}}}, postaction={decorate}]}
\def\wavey{\draw[thick, decorate, decoration={snake, segment length=2.5mm, amplitude=0.75mm}]}
\newcommand{\AxisRotatorL}[1][rotate=0]{%
    \tikz [x=0.25cm,y=0.60cm,line width=.2ex,-stealth,#1] \draw (0,0) arc (-150:150:1 and 1);%
}
\newcommand{\AxisRotatorR}[1][rotate=0]{%
    \tikz [x=0.25cm,y=0.60cm,line width=.2ex,-stealth,#1] \draw (0,0) arc (150:-150:1 and 1);%
}

\begin{document}

\captionsetup[figure]{margin=1.5cm,font=small,labelfont={bf},name={Figure},labelsep=colon,textfont={it}}
\captionsetup[table]{margin=1.5cm,font=small,labelfont={bf},name={Table},labelsep=colon,textfont={it}}
\setlipsumdefault{1}

\frontmatter

\begin{titlepage}
	\centering
	    \scshape % Use small caps for all text on the title page
        \vspace*{\baselineskip} % White space at the top of the page
        
	    \rule{\textwidth}{1.6pt}\vspace*{-\baselineskip}\vspace*{2pt} % Thick horizontal rule
	    \rule{\textwidth}{0.4pt} % Thin horizontal rule
	    
	    \vspace{0.75\baselineskip} % Whitespace above the title
	    
	    {\LARGE Conformal Field Theory} % Title
	    
	    \vspace{0.75\baselineskip} % Whitespace below the title
	    
	    \rule{\textwidth}{0.4pt}\vspace*{-\baselineskip}\vspace{3.2pt} % Thin horizontal rule
	    \rule{\textwidth}{1.6pt} % Thick horizontal rule
	    
        \vspace{5\baselineskip} % Whitespace after the title block

	    Course delivered in 2020 by 
	
	    \vspace{0.5\baselineskip} % Whitespace before the editors
	
	    {\scshape\Large Dr. Paul Heslop} % Lecturer Name
	
	    \vspace{0.5\baselineskip} % Whitespace below the editor list
	
	    \textit{Durham University} % Lecturer Institution
	    
	    \vspace{2\baselineskip} % Whitespace after the title block
	    
	    \includegraphics[width=8cm]{images/DurhamLogo.png}\\[1cm] % Logo 

	    Notes taken by 
	
	    \vspace{0.5\baselineskip} % Whitespace before my name
	
	    {\scshape\Large Richie Dadhley} % Lecturer Name
	   
	    \vspace{0.5\baselineskip} % Whitespace below my name
	    \textit{richie@dadhley.com} % Email
	
	    \vfill % Whitespace between editor names and publisher logo
\end{titlepage}



% -------------------------------------------------------------------
% Acknowledgements
% -------------------------------------------------------------------

\newpage
\section*{Acknowledgements}

These are my notes on the 2020 lecture course "Conformal Field Theory" taught by Dr. Paul Heslop at Durham University as part of the Particles, Strings and Cosmology Msc. For reference, the course lasted 16 hours and was taught over 4 weeks. \\

I have tried to correct any typos and/or mistakes I think I have noticed over the course. I have also tried to include additional information that I think supports the taught material well, which sometimes has resulted in modifying the order the material was taught. Obviously, any mistakes made because of either of these points are entirely mine and should not reflect on the taught material in any way. \\

% For any other notable sources used:
% One of the main sources I have used to obtain the additional information is Dr. David Tong’s String Theory notes, available via the following link  

%\begin{center}
%    \textcolor{red}{Link}
%\end{center}

I would like to extend a message of thanks to Dr. Heslop for teaching this course. I would also like to thank Thimo Preis for helping getting these notes started.\\

If you have any comments and/or questions please feel free to contact me via the email provided on the title page. \\

These notes are currently a work in progress, so for updated notes, as well as a list of other notes/works I have available, visit my blog site

\begin{center}
    \href{https://richie291.wixsite.com/theoreticalphysics}{https://richie291.wixsite.com/theoreticalphysics}
\end{center}

These notes are not endorsed by Dr. Heslop or Durham University.

\vspace{1cm}

\begin{flushright}
    \Huge{{\cursive\setul{0.1ex}{}\ul{~~Richie Dadhley~~}}}
\end{flushright}

% -------------------------------------------------------------------
% Contents
% -------------------------------------------------------------------

\tableofcontents

% -------------------------------------------------------------------
% Main sections (as required)
% -------------------------------------------------------------------

\mainmatter

\chapter{Preliminary Comments On Notes}

These notes were taken during lectures themselves, but have undergone significant editing afterwards by both myself and Paul. In some places this has resulted in me (Richie) adding completely new material that was not discussed during the course. This is mainly done for my own understanding on the background of some of the results used. Sections that contain a lot of unlectured material will be marked with an asterisk.

\mybox{
    \begin{center}
        The important results that appear in these section will be put in boxes like this.
    \end{center}
}

I will try not to rely too heavily on the details of these additional sections, and any calculations used later in the course I shall present again. This is done so that, in principle, the non-asterisked sections can be read in a self contained way (up to accepting the boxed results). The disclaimer is that I will use the notation from these sections (e.g. putting tildes on the charges associated to the conformal differential operators), so if any notation seems strange, the reader is requested to glance these sections for clarity.  

\badr 
    Other, non-asterisked, sections will also include small additions here and there and these will normally be included as additional remarks like this.\footnote{Some will also appear as footnotes.} 
\eadr

\br
    Within the asterisked sections I will just use `normal' remarks (i.e. not "Additional Remark"), like this.
\er     

Of course some of these additions will naturally be part of the text itself. 

All this additional material should be understood with caution as I may have misinterpreted results I've seen from places. Of course any errors etc in these parts are entirely my fault and should not reflect on the course itself. 
\chapter{Introduction \& Motivations}

We start this course with a introduction and motivation of the work that is to follow. It is beneficial to do this as a some of the work that follows can be rather abstract and it is always useful to have some kind of grounding to remind us why we care about what we're doing. However it is worth saying that the things said in the introduction aren't meant to be fully understood on a first read, but will slowly `fall into place' as the course goes on.\footnote{I simply say this because I have felt overwhelmed at introductions before.} So without further ado, let's get introducing and motivating. 

\section{What is CFT ?}

Of course the most immediate question to ask is "what is conformal field theory (CFT)?" This is something that will become more precise as we progress, but as a guiding principle, a CFT is a \textit{scale-invariant} quantum field theory (QFT).\footnote{We should clarify and say that there are indeed ways to get CFTs without making reference to a Lagrangian, something you usually have in a QFT.} By scale-invariant we mean that you can expand/shrink the underlying spacetime without changing the theory. In other (even more rough) words you can't tell how `zoomed in' you are to the theory, and you don't really care. Not only is a CFT invariant under \textit{global} scalings (i.e. scale the whole spacetime by the same amount, so relative lengths are preserved), but it is also invariant under \textit{local} scalings. Essentially what this means is that a CFT is invariant under any transformation that preserves the \textit{angles}\footnote{In the Euclidean sense.} between things, but it need not preserve lengths. On top of this scaling symmetry, a CFT also possesses Poincar\'{e} symmetry. This follows simply from the fact that a CFT is a QFT (which we construct with Poincar\'{e} invariance in mind!)

\begin{center}
    \btik 
        \draw[thick, ->, rotate around={70:(0,0)}] (0,0) -- (2,0) node [right] {$v_1$};
        \draw[thick, ->] (0,0) -- (3,0)  node [right] {$v_2$};
        \draw[->] (1,0) arc (0:70:1cm);
        \node at (0.5,0.4) {$\theta$};
        %%
        \draw[->] (4.5,1) -- (7.5,1) node [midway, above] {Conformal Transformation};
        %%
        \draw[thick, ->, rotate around={150:(12,0)}] (12,0) -- (15,0) node [left] {$v_1^{\prime}$};
        \draw[thick, ->, rotate around={80:(12,0)}] (12,0) -- (13.5,0)  node [right] {$v_2^{\prime}$};
        \draw[->, rotate around={80:(12,0)}] (13,0) arc (0:70:1cm);
        \node at (11.75,0.5) {$\theta$};
    \etik  
\end{center}

Combining these symmetries can lead to further symmetries. It turns out that in $2$-dimensions we have infinitely many symmetries, while in higher dimensions there is a finite number. We can collect all these symmetries into a group, known as the \textit{conformal group}, and thus we can define a CFT to be a QFT that is invariant under the conformal group.

\br 
    It is important to note that the real world is \textit{not} scale invariant! For example QCD is not scale invariant as 
    \ben[label=(\roman*)]
        \item The gluon masses: this just comes from the idea that mass and length are related (i.e. the higher the energy the smaller the scale). 
        \item Even in massless QCD, we introduce a scale via the renormalisation process, i.e. the cut-off, and so we still break scale invariance. 
    \een 
    In general for a theory to be CFT, the classical Lagrangian has to be scale-invariant. When you quantize it, it either retains scale-invariance or it does not, and as we have just explained QCD does not. 
\er 

\section{Why CFT?}

Ok so armed with our introductory understanding of what CFT is, we can ask the next important question "Why do we care about CFTs?" Well there's always the standard answer of "things are interesting to study," but really we would like a more motivating answer. 

It turns out that there exist many theories similar to QCD which \textit{are} scale invariant, e.g. $\cN=4$ super-Yang-Mills theory.\footnote{This is the most symmetry QFT you can have in 4-dimensions. On top the conformal symmetry, we have 4 supersymmetries. Obviously more details on supersymmetry in my SUSY notes. Also "Yang-Mills" is just the fancy name for non-Abelian gauge theory --- see my QFT II notes for more details.} This is worth studying as it plays a role in the study of the AdS/CFT correspondance.

Furthermore RG flow fixed points of \textit{any} QFT are conformally invariant (and thus CFTs).

More recently, \textit{numerical conformal bootstrap} procedures are ongoing research. We will discuss these in a bit more detail later, but essentially this is a way to constrain a system by imposing self-consistency checks. This allows us to find stuff out about the theory without having to do any explicit calculations (e.g. Feynman diagrams).

The above motivational examples are, almost certainly, completely new names and so might not provide a huge motivation. The next example will perhaps provide more motivation: string theory is a $2$-dimensional CFT! Indeed almost any string theory course you will read will most likely contain a reasonable amount of information presented in this course. It is really important to remember, though, that this is only a specific application of CFT and so it is worth studying outside that. Note also that string theory happens to fall into our `special' class of CFTs, namely 2-dimensional ones. For this reason we should be careful about using string theory has our only grounding when it comes to studying CFTs.

\section{Parts}

Part I of this course will focus on studying $D>2$ CFTs (with some $D=2$ comments made along the way) then Part II will focus on $D=2$ CFTs.
\part{CFT in $D > 2$}
\chapter{Conformal transformations}

Now that we are (hopefully) motivated, we can begin a more technical study of conformal symmetry. We start by being a bit more precise as to what each of the symmetries we have mentioned mean. 

\section{Symmetries \& Conformal Killing Equation}

\bd
    The \textit{Poincar\'{e} group} is the full set of symmetries of special relativity, i.e. it is the union of the Lorentz symmetries (spatial rotations and boosts) with translations. Poincar\'{e} transformations can therefore be defined as the set of transformations that map flat space to flat space. 
\ed

We should be familiar enough with Poincar\'{e} transformations to not need any examples.

We can put a QFT on an arbitrary background, coupling the QFT to gravity, resulting in a theory which is diffeomorphism\footnote{A diffeomorphism is the structure preserving map between two manifolds. Diffeomorphisms play a \textit{huge} role in GR, where they (roughly speaking) correspond to coordinate changes. Diffeomorphisms are  gauge symmetry in GR and essentially give rise to the connection coefficients, $\Gamma$, in exactly the same way that the $SU(3)$ gauge symmetry of QCD gives gluons (although this comparison is not always made!).} invariant. Fixing the background to flat space we recover the original QFT and the  Poincar\'{e} transformations are simply the diffeomorphisms which leave the flat metric invariant.

Some QFTs when coupled to gravity in this way possess an additional symmetry known as a Weyl transformation, simply a local rescaling of the metric.

\bd
    A \textit{Weyl transformation} is a local rescaling of the metric by a factor, i.e. it maps 
    \bse 
        g \mapsto \Omega^2(x) g,
    \ese 
    where the squared is purely convention. We call $\Omega^2(x)$ a \textit{Weyl factor}.
\ed 

A Weyl invariant theory when restricted to a fixed flat background metric is a called a CFT (at least classically).

\bd 
    A \textit{conformal transformation} is thus a spacetime transformation that leaves the \textit{flat} metric invariant up to a Weyl factor.
\ed 

\badr 
    There is a subtle difference between a Weyl transformation and a conformal transformation. A Weyl transformation is a scaling of the metric itself and it does not effect the spacetime itself, while a conformal transformation leaves the metric alone but changes the underlying spacetime. We can see the difference between the two by transformation in terms of coordinates:\footnote{If we want to be a bit more fancy, we can define the difference using the notion of a pull-back. Let $f:\cM\to\cM$ be a smooth map with $\cM$ being our manifold. Then a conformal transformation is defined via $f^* g = \Omega^2(x) g$, where $\Omega^2(x) \in C^{\infty}(\cM)$. On the other hand a Weyl transformation has $f=\b1_{\cM}$ and we simply have $h = \Omega^2(x)g$.} 
    \begin{itemize}
        \item A Weyl transformation is defined via 
        \bse 
            h_{\mu\nu} = \Omega^2(x) g_{\mu\nu}.
        \ese
        \item A conformal transformation is defined via 
        \bse
            \frac{\p x^{\prime\a}}{\p x^{\mu}} \frac{\p x^{\prime\beta}}{\p x^{\nu}} \eta^{\prime}_{\a\beta} = \Omega^2(x) \eta_{\mu\nu}, 
        \ese 
        where the left-hand side of this expression is just the general formula for how the components of the metric change under a coordinate transformation.
    \end{itemize}
    These two things are \textit{almost} essentially the same, and it is likely that I will slip up here and there and use them interchangeably, but it's important to note they are indeed different things. 
\eadr 

Let's try break down conformal transformations a little bit further. Recall that the metric essentially tells us the lengths and angles between things. The length comes from taking $g(V,V)$, while the angle comes from considering $g(U,V)$ along with the lengths of $U$ and $V$. If we take a conformal transformation then the lengths are clearly going to change, as we have the $\Omega^2(x)$ factor. However when looking at angles really all we are interested in is the projections of lengths, and here the $\Omega^2(x)$ factors cancel. That is the lengths of \textit{both} $U$ and $V$ change by the same amount\footnote{This follows from the fact that we can only take the angle between them if they are defined at the same $x$, so the scaling is the same for both.} and so the ratio to the projection remains invariant. We can therefore intuitively think of conformal transformations as spacetime transformations that locally vary lengths but preserve angles. This is exactly the introductory definition we gave above. 

The above result is actually very important, and so we write it again in a nice box. 
\mybox{
    \be 
    \label{eqn:ConformalTransformation}
        \frac{\p x^{\prime\a}}{\p x^{\mu}} \frac{\p x^{\prime\beta}}{\p x^{\nu}} \eta^{\prime}_{\a\beta} = \Omega^2(x) \eta_{\mu\nu}. 
    \ee 
}

\bnn
    From now on we will assume we are working in a flat spacetime, i.e. we replace $g_{\mu\nu}$ with $\eta_{\mu\nu}$. 
\enn 

The above result holds for a general \textit{finite} conformal transformation. We can use it to find out how an infinitesimal conformal transformation acts. Start by writing
\bse
    x^{\prime \mu} = x^{\mu} + \xi^{\mu} (x), \qand \Omega = 1 +\kappa(x),
\ese
where $\xi$ and $\kappa$ are infinitesimal (i.e. ignore $\cO(\xi^2),\cO(\kappa^2),\cO(\xi \kappa)$). Then inserting this into \Cref{eqn:ConformalTransformation} gives us 
\mybox{
    \be
    \label{eqn:ConformalKilling}
        2 \kappa(x) \eta_{\mu \nu}= \p_{\mu} \xi_{\nu} + \partial_{\nu} \xi_{\mu},
    \ee
}
\noindent which is known as the \textit{conformal Killing equation}. It is the defining equation for infinitesimal conformal transformations.

\br 
    Note we call it the conformal \textit{Killing equation} as it is of the same form as the Killing equation from general relativty, which reads\footnote{In flat space, otherwise we replace $\p_{\mu}$ with $\nabla_{\mu}$.}
    \be 
    \label{eqn:KillingEquation}
        \p_{\mu}\xi_{\nu} + \p_{\nu}\xi_{\mu} = 0.
    \ee 
    Note it is the $\kappa$ that distinguishes the two, and this is where the "conformal" part of the name comes from; $\kappa$ comes from the $\Omega^2(x)$ factor. 
\er 

\section{Conformal Killing Equation For $d>2$}

The idea is to solve \Cref{eqn:ConformalTransformation} to get the infinitesimal symmetries and then we can get the full picture back using the exponential map. The general solution (for $d>2$, with $d$ being the spacetime dimension) is
\be
\label{eqn:ConformalKillingSolution}
    \xi^{\mu}(x) = a^{\mu} + {\omega^{\mu}}_{\nu} x^{\nu} + \sigma x^{\mu} + b^{\mu} x^2- 2 b_{\nu} x^{\nu} x^{\mu}, 
\ee 
with (contracting \Cref{eqn:ConformalKilling} with $\eta^{\mu \nu}$)
\bse 
    \kappa(x)=\frac{1}{d} \p_{\mu} \xi^{\mu} = \sigma-2 b_{\nu} x^{\nu},
\ese  
where 
\ben[label=(\roman*)]
    \item $\omega$ is the Lorentz part, so it must obey $\omega_{\mu\nu}=-\omega_{\nu\mu}$,
    \item $\sigma$ is the scaling part (known as a \textit{dilatation}), and
    \item $b$ is something new we call the \textit{special conformal transformation}.
\een 

\bbox
    \ben
        \item Check that the above solves \Cref{eqn:ConformalKilling}.
        \item Derive \Cref{eqn:ConformalKillingSolution}. \textit{Hint: Start by recalling (or deriving) the general solution to Killings equation, \Cref{eqn:KillingEquation}, and then include the contributions from $\kappa$. You want to show that $\kappa$ is linear in $x$, i.e. two derivatives on it are zero. Then go from there by picking suitable linear combinations to find the solution.}
    \een
\ebox 

\br 
    As we said above, \Cref{eqn:ConformalKillingSolution} only holds for dimensions $d \neq 2$. As we have mentioned a couple times, for the specific case of $d=2$ things change and we get a much larger expression. 
\er 

\subsection{The Conformal Algebra}

\bcl 
    The conformal group is a Lie group.\footnote{Note that the conformal group is \textit{not} a compact Lie group, though. This is easily seen from the fact that the Lorentz group is non-compact. See Remark 6.1.1 of my Group Theory notes for a little more detail.}
\ecl 

We can pick a basis of the Lie algebra so that our conformal Killing vector $\xi^{\mu}\p_{\mu}$ is a generator of the Lie group. We can do this by decomposing it in terms of the generators of the subgroups. We have 
\ben[label=(\roman*)]
    \item Momentum:
    \be 
    \label{eqn:MomentumGenerator}
        P_{\mu} := \p_{\mu},
    \ee
    which generates spacetime translations, $a^{\mu}$.
    \item Lorentz: 
    \be 
    \label{eqn:LorentzGenerator}
        L_{\mu\nu} := x_{\mu}\p_{\nu} - x_{\nu}\p_{\mu},
    \ee 
    which generate our boosts and spatial rotations, $\omega^{\mu\nu}$.
    \item Dilatations: 
    \be 
    \label{eqn:DilatationGenerator}
        D := x^{\mu}\p_{\mu},
    \ee 
    which generates our scale transformations, $\sig$.
    \item Special conformal generator:
    \be 
    \label{eqn:SpecialConformalGenerator}
        K_{\mu} := x^2 \p_{\mu} - 2x_{\mu}x^{\nu} \p_{\nu},
    \ee 
    which generates our new $b^{\mu}$ parameter. 
\een 
From \Cref{eqn:ConformalKillingSolution} we then read off 
\bse
    \xi^{\mu} \p_{\mu} = a^{\mu} P_{\mu} + \frac{1}{2} {\omega^{\mu}}_{\nu} {L^{\nu}}_{\mu} + \sigma D + b^{\mu} K_{\mu},
\ese 

Now we have claimed that our conformal group is a Lie group and that \Crefrange{eqn:MomentumGenerator}{eqn:SpecialConformalGenerator} are our generators. If this is the case they must be elements of the associated Lie algebra, and so we really should check that they are closed under the commutator (which is the Lie bracket here). 

\bp 
    The generators \Crefrange{eqn:MomentumGenerator}{eqn:SpecialConformalGenerator} satisfy the following commutation relations. 
    \mybox{
        \be 
        \label{eqn:GeneratorCommutators}
            \begin{split}
                [P_{\mu} , L_{\nu\rho}] & = \eta_{\mu\nu}P_{\rho} - \eta_{\mu\rho}P_{\nu}, \\
                [P_\mu, K_\nu] & = -2 \eta_{\mu \nu} D+2 L_{\mu \nu}, \\
                [K_\mu, L_{\nu \rho}] & = \eta_{\mu \nu} K_\rho - \eta_{\mu \rho} K_\nu, \\
                [K_\mu, D] & = - K_\mu \\
                [P_\mu, D] &= P_\mu, \\
                [L_{\mu \nu}, L_{\rho \sigma}] & = \eta_{\nu \rho} L_{\mu \sigma} + \eta_{\mu \sigma}  L_{\nu \rho} -\eta_{\mu \rho} L_{\nu \sigma} - \eta_{\nu \sigma} L_{\mu \rho},
            \end{split}
        \ee 
    }
    \noindent with all others vanishing.
\ep 


\bbox 
    Derive the commutation relations \Cref{eqn:GeneratorCommutators}. \textit{Hint: Recall that the commutator of differential operators only really make sense when acting on a function, $f$. Then remember (or show) that the $\p \p f$ terms always cancel.}
\ebox 

Now these commutation relations are not the prettiest things to remember, so it would be nice if we could repackage the information in a nicer way. We can do this by introducing two new `dimensions', which we call "$-1$" and "$d$". In other words we let $M, N \in \{ -1 , 0, ... , d-1 , d\}$, where we have used new indices to distinguish them from our Lorentz ones, $\mu,\nu$. We then make the following definition.
\bd
    Using our new index range, we define $J_{MN}$ via
    \be
    J_{\mu \nu} = L_{\mu \nu}, \quad J_{-1 \mu} = \frac{1}{2} (P_\mu + K_\mu) \quad J_{d \mu}  = \frac{1}{2} (P_\mu -K_\mu) \quad J_{d -1} = D, \quad J_{MN} = -J_{NM}.
    \ee
\ed 

We can write this definition as a matrix\footnote{The convention here is that the first index, i.e. $M$, tells us the row, and the second index, $N$, tells us the column.}
\bse
    J_{MN} =
    \begin{pmatrix}
    0 & \dots & \frac{1}{2}\big(P_{\nu} + K_{\nu}\big) & \dots & -D \\
    \vdots && \dots && \vdots  \\
    -\frac{1}{2}\big(P_{\mu} + K_{\mu}\big) && L_{\mu\nu} && \frac{1}{2}\big(K_{\mu} - P_{\mu}\big) \\
    \vdots && \dots && \vdots \\
    D & \dots & \frac{1}{2}\big(P_{\nu} - K_{\nu}\big) & \dots & 0 
    \end{pmatrix}.
\ese 
It follows from this that $J_{MN}$ satisfies 
\bse
    [J_{mn},J_{PQ}] = \eta_{n P} J_{m Q} + \eta_{m Q} J_{n P} - \eta_{m P} J_{n Q} - \eta_{n Q} J_{m P},
\ese 
where $\eta_{MN}= \text{diag}(-1,-1,+1,...,+1)$, but this is a representation of a Lorentz-type Lie algebra (compare it to the last line of  \Cref{eqn:GeneratorCommutators}). In other words this is a representation of $\so(2,\mathrm{d})$.

\subsection{Finite Conformal Transformations}

Now that we have the conformal Lie algebra, we can use the exponential map to recover a \textit{finite} conformal transformation:
\be 
\label{eqn:FiniteConformalTransformation}
    x^{\prime \nu} = e^{\xi^{\mu} \p_{\mu}} x^{\nu} = x^{\nu} + \xi^{\mu} \p_{\mu} x^{\nu} + \frac{1}{2} \xi^{\mu}\p_{\mu} \big(\xi^{\rho}\p_{\rho} x^{\nu}\big) + \cO(\xi^3).
\ee  

Let's look at some examples. 

\bex 
    The easiest example is to consider just a translation, i.e. $\xi^{\mu} = a^{\mu}$. Plugging this into \Cref{eqn:FiniteConformalTransformation} we have 
    \bse
        \begin{split}
            x^{\prime \nu} & = x^{\nu} + a^{\mu} \delta^{\nu}_{\mu} + \frac{1}{2}a^{\mu} \p_{\mu}\big( a^\rho \delta^{\nu}_{\rho} \big) \\
            & = x^{\nu} + a^{\nu},
        \end{split}
    \ese
    where the second line follows from the fact that $a^{\rho}$ is constant.
\eex 

\bex 
    Next let's consider just a dilatation, i.e. $\xi^{\mu} = \sig x^{\mu}$. We then have 
    \bse 
        \begin{split}
            x^{\prime \nu} & = x^{\nu} + \sigma x^{\mu} \delta^{\nu}_{\mu} + \frac{1}{2} \sigma x^{\mu} \p_{\mu} \big( \sigma x^{\rho} \delta^{\nu}_{\rho} \big) + \dots \\
            & = x^{\nu} +\sigma x^{\nu} + \frac{1}{2} \sigma^2 x^{\nu} + \dots \\
            & = e^{\sigma} x^\nu.
        \end{split}
    \ese 
    Now, as $\sig$ is a constant, this just tells us to scale our coordinates about $x=0$: 
    \begin{center}
        \btik 
            \draw[thick, dashed, <->] (-1,-1) -- (1,1);
            \draw[thick, dashed, <->] (-1,1) -- (1,-1);
            \draw[thick, ->] (1,1) -- (2,2);
            \draw[thick, ->] (1,-1) -- (2,-2);
            \draw[thick, ->] (-1,1) -- (-2,2);
            \draw[thick, ->] (-1,-1) -- (-2,-2);
            \draw[fill=black] (1,1) circle [radius=0.05cm];
            \draw[fill=black] (1,-1) circle [radius=0.05cm];
            \draw[fill=black] (-1,1) circle [radius=0.05cm];
            \draw[fill=black] (-1,-1) circle [radius=0.05cm];
            \node at (1,0.7) {$x$};
            \node at (1,-0.7) {$x$};
            \node at (-1,0.7) {$x$};
            \node at (-1,-0.7) {$x$};
            \node at (2,1.7) {$x^{\prime}$};
            \node at (2,-1.6) {$x^{\prime}$};
            \node at (-2,1.7) {$x^{\prime}$};
            \node at (-2,-1.6) {$x^{\prime}$};
        \etik 
    \end{center}
\eex 

We have seen a translation and dilatation, but what about our special conformation transformations? That is what does
\bse 
    \xi^\mu = b^\mu x^2 - 2 b_\nu x^\nu x^\mu
\ese 
give us. 
\bcl 
    The special conformal transformations correspond to the finite transformation 
    \be
    \label{eqn:FiniteSpecialConformalTransformation}
        x^{\prime \nu} = e^{\xi^{\mu} \p_{\mu} } x^{\nu} = \frac{x^{\nu} +x^2 b^{\nu}}{1+2 b_{\nu} x^{\nu} +b^2 x^2}
    \ee 
\ecl 

\bbox
    Show \Cref{eqn:FiniteSpecialConformalTransformation} holds up to $\cO(b^2)$, i.e. expand both sides. 
\ebox 

We can now check that these finite transformations agree with the original definition, i.e. check that \Cref{eqn:ConformalTransformation} holds for the above $x^{\prime \nu}$s with some $\Omega(x)$.
\bbox
    Given that 
    \be 
    \label{eqn:OmegaSpecialConformal}
        \Omega^{-1}(x) = 1+2 b_\nu x^\nu +b^2 x^2
    \ee
    for special conformal transformations, check that \Cref{eqn:ConformalTransformation} is satisfied for special conformal transformations. That is use the given formula along with \Cref{eqn:FiniteSpecialConformalTransformation} to verify our definition holds.
\ebox 

\subsection{Inversions}

Although \Cref{eqn:FiniteSpecialConformalTransformation} looks unapealing at first, it actually allows us to see a really nice result. Consider 
\be
\label{eqn:xPrimedSquared}
    (x^\prime)^2 = \frac{(x+ x^2 b)^2}{(1+2b \cdot x+b^2 x^2)^2}=\frac{x^2+2 x^2 x\cdot b +x^4 b^2}{(1+2b\cdot x +x^2 b^2)^2}= \frac{x^{\prime} \cdot (x + x^2b) }{1+2 b\cdot x +x^2 b^2},
\ee 
from which we see that 
\be
\label{eqn:xByxSquared}
    \frac{x^{\prime \mu}}{x^{\prime 2 }} = \frac{x^\mu + x^2 b^\mu}{x^2} = \frac{x^\mu}{x^2}+b^\mu. 
\ee 
In order to understand what this tells us, let's introduce another definition. 

\bd[Inversion Map] 
    We define the \textit{inversion map} via 
    \be 
    \label{eqn:InversionMap}
        I: x^\mu \rightarrow \frac{x^\mu}{x^2}.
    \ee 
\ed 

As the name suggests, the inversion map inverts our coordinates, e.g. it sends $0 \to \infty$. Now we see that \Cref{eqn:xByxSquared} tells us that if we invert our coordinates then a special conformal transformation corresponds simply to a translation. That is 
\bse 
    I x^{\prime\mu} = I x^{\mu} + b^{\mu},
\ese
which is our finite translation formula. Infinitesimally we can write this as
\be
\label{eqn:KEqualIPI}
    K = I P I.
\ee 
\chapter{Conformal Transformation Of Classical Fields}

The fundamental objects in a CFT are \textit{local}\footnote{The reason for emphasis on the word "local" will become clearer when we look at the state operator correspondance.} fields, which  become local operators on quantisation. We can then use these operators to work out the correlation functions of the theory. When we know all the operators\footnote{That is, knowing what representation of the conformal group they transform under.} and their correlation functions we say that we can \textit{solve} the theory. Said another way, if we know all the operators and all the different correlation functions between them we can work out anything we would want to know about the theory. 

\bex 
    In a theory that has a Lagrangian, the fields will be simply linear combinations of products of (derivatives of) the fundamental fields appearing in $\cL$. For example: 
    \ben[label=(\roman*)]
        \item If you have a scalar field theory the fundamental field is $\phi(x)$. Our Lagrangian is then built out of products, $\phi^n(x)$, and derivative terms $\p_{\mu}\phi(x)$. We also have products of both kinds, e.g. $\phi(x)\p_{\mu}\phi(x)$.
        \item If you are considering a gauge theory, you must ensure that the fields are gauge invariant. So $A_{\mu}(x)$ is not considered by itself, as it is not gauge invariant. However terms such as $\tr[F_{\mu\nu}F_{\rho\sig}]$, with $F_{\mu\nu} = \p_{\mu}A_{\nu}-\p_{\nu}A_{\mu} + ig[A_{\mu},A_{\nu}]$, will appear.\footnote{If you are not familiar with such terms, see, for example, my QFT II notes.}
    \een 
\eex 

\br 
\label{rem:FieldsCanHaveLorentzIndices}
    Note that our fields can transform outside the trivial representation of the Lorentz group, as the Lorentz group is a subgroup of the conformal group. However gauge transformations are non-physical and do not form a subgroup, so we have to ensure gauge invariance if we want to calculate physical things. By this we just mean that our operators, $\cO$, will be allowed to carry Lorentz indices $\mu,\nu$ etc., however they \textit{cannot} carry Lie algebra indices $a,b$ etc.
\er 

\section{Transformation Of Fields}

So far we have seen how conformal transformations act on the spacetime coordinates $x^{\mu}$, however we are yet to see how the conformal transformations act on the fields themselves. This kind of question should sound familiar from a introductory field theory course where you discuss active vs. passive transformations of the Lorentz group. We shall recap the general idea here.\footnote{For a slightly more intuitive explanation, see Section 2.1.1 of my IFT notes.}

\subsection{Recap Of Lorentz Transformations Of Fields}

Under a Lorentz transformation the space transforms as (indices left implicit for notational convenience) 
\bse 
    x \to x^{\prime} = \Lambda x
\ese 
while the scalar field transforms as 
\bse 
    \phi \to \phi^{\prime}.
\ese 
Now if $\phi$ is a scalar it is, essentially by definition, invariant under the transformation. Mathematically that is
\bse 
    \phi^{\prime}(x^{\prime}) = \phi(x).
\ese 
If we put this all together, we can conclude
\mybox{
    \be 
    \label{eqn:ScalarFieldTransformation}
        \phi^{\prime}(x) = \phi(\Lambda^{-1}x).
    \ee 
}

Ok great, we know how scalar fields transform, however we just said in \Cref{rem:FieldsCanHaveLorentzIndices} that our fields can have Lorentz indices, the next question is how do these transform? Well of course the argument still transforms as above, and we just need to take into account the transformation of the Lorentz index, e.g.
\bse 
    V^{\mu} \to V^{\prime\mu} 
\ese 
where 
\be 
\label{eqn:VmuTransform}
    V^{\prime\mu}(x^{\prime}) = {\Lambda^{\mu}}_{\nu} V^{\nu}(x).
\ee

\badr 
    Technically speaking we should write 
    \bse 
        V^{\prime \mu} (x^{\prime}) = {D[\Lambda]^{\mu}}_{\nu} V^{\nu}(x),
    \ese
    where $D$ is the representation of the Lorentz group that $V^{\mu}(x)$ transforms in. In the above we have just used the fundamental representation, as this is the representation for $4$-vectors. 
\eadr 

\subsection{Adding Dilatations}

Ok so we know how the fields transform under the Lorentz subgroup of our conformal group. We now need to ask how our dilatations act. Recall that these are "scalings" and so anything that has dimension will be effected by these. That is, even the fundamental scalar fields will transform non-trivially. To be more clear, recall that 
\bse 
    S = \int d^dx  \, (\p_{\mu}\phi)^2 \qquad \implies \qquad [\phi] = \frac{d-2}{2}.
\ese
and so our scalar fields are effected by dilatations. The question we want to ask is "how are they effected?" Well recall that natural units ($c=1=\hbar$) are defined by mass dimensions, and it follows from Einstein's little equation $E=mc^2$, that in natural units $[E]=[m]$. Then from Heisenberg uncertainty that $[E] = -[t]$, and so spacetime lengths and mass will scale oppositely. That is things with positive mass dimensions "shrink" when we use a dilatation to "zoom in". We summarise this with the following definition. 

\bd[Dilatation Weight]
    The \textit{dilatation weight} (or just \textit{weight}), $\Delta$, is defined to equal to the mass dimension. 
\ed 

For clarity, the dilatation weight of our scalar field above is 
\bse 
    \Delta_{\phi} = \bigg(\frac{d-2}{2}\bigg).
\ese 
So if we use a dilatation of the form 
\be 
\label{eqn:xDilatation}
    x^{\mu} \to x^{\prime\mu} = \l x^{\mu}, \qquad \l\in \R 
\ee 
we have 
\be
\label{eqn:phiDilatation}
    \phi^{\prime}(x^{\prime}) = \l^{-\big(\frac{d-2}{2}\big)} \phi(x),
\ee
where the prefactor is coming from the fact that $\phi$ has mass dimensions and $[m] = [x^{-1}] = -[x]$ and so scales with the negative power.

More generally, all local fields have a dilatation weight, and scale under the dilatation \Cref{eqn:xDilatation} as
\mybox{
    \be 
    \label{eqn:ScalarFildDilatation}
        \phi^{\prime}(x^{\prime}) = \l^{-\Delta} \phi(x).
    \ee
}

\badr 
    Note the above formula makes sense; it is just the dilatation equivalent of including a term for non-trivial Lorentz transformations, as in \Cref{eqn:VmuTransform}. The equivalent of the trivial representation (i.e. a Lorentz scalar) is simply $\Delta=0$, which gives us precisely $\phi^{\prime}(x^{\prime})=\phi(x)$. % Remember to ask Paul about the $f(x) = x^n$ example. Wiki gives it as an example, are they just wrong?
\eadr 

\section{Primary Fields}

All we have left to add are the special conformal transformations, however here we just summarise how a scalar field transforms under a general conformal transformation. The first type of field we consider are so-called scalar primary fields.

\bd[Primary Field]
    A scalar \textit{primary field}, $\phi(x)$, with weight $\Delta$, transforms under a general conformal transformation as 
    \be 
    \label{eqn:PrimaryFieldTransformation}
        \phi^{\prime}(x^{\prime}) = \Omega^{-\Delta}(x) \phi(x),
    \ee 
    where $\Omega(x)$ is a general conformal transformation. 
\ed 

Infinitesimally we can write \Cref{eqn:PrimaryFieldTransformation} as
\bse 
    \del\phi(x) =  -\kappa(x) \Delta \phi(x) - \xi^{\mu}\p_{\mu}\phi(x),
\ese 
where the second term comes from the fact that we have shifted the coordinates. 

\bbox 
    Use this definition to show that a scalar primary field transforms under special conformal transformations as 
    \be 
    \label{eqn:SpecialConformalTransformationOfPrimaryField}
        \phi(x) = \frac{\phi(x^{\prime})}{(1+2b^{\mu}x_{\mu} +b^2x^2)^{\Delta}}. 
    \ee 
    % In response to Pual's question about if this was from him, no but it's in Simon Ross' notes, equation (4.4).
\ebox 

Ok that's scalar primaries, but what about non-scalar primaries? Well non-scalar primaries carry some form of Lorentz index (e.g. $\mu$ if it is in the vector representation or $\a$ if it is in the spinor representation). In order to account for all of these we shall denote a general index by $I$,\footnote{Not to be confused with the $I$ for inversions.} $J$ etc. All we have to do, then, is account for the transformation of the Lorentz indices, so our non-scalar primary field transforms (infinitesimally) as 
\bse 
    \del\phi_{\Delta,I}(x) = -\kappa(x) \Delta \phi_{\Delta,I}(x) - \xi^{\mu}(x)\p_{\mu}\phi_{\Delta,I}(x) + {\rho^{\mu}}_{\nu} {({S^{\nu}}_{\mu})_I}^J \phi_{\Delta, J}(x),
\ese 
where we have used
\bse 
    \begin{split}
        \kappa(x) & = \sig - 2b\cdot x \\
        {\rho^{\mu}}_{\nu} & = {\omega^{\mu}}_{\nu} + 2(b^{\mu}x_{\nu} - x^{\mu}b_{\nu}) \\
        \xi^{\mu}(x) & = a^{\mu} + {\omega^{\mu}}_{\nu}x^{\nu} + \sig x^{\mu} + b^{\mu} x^2 - 2b\cdot x x^{\mu}
    \end{split}
\ese 
and $S^{\mu\nu}$ is the appropriate matrix for the Lorentz representation in question\footnote{For example for spinors $S^{\mu\nu} = \frac{1}{4}[\g^{\mu},\g^{\nu}]$.}

Now of course, by taking coefficients of $a_{\mu}, b_{\mu}, \sig, \omega^{\mu\nu}$, we get the action of $P_{\mu}, K_{\mu}, D, L_{\mu\nu}$ (respectively) on conformal primary fields, the results are set as an exercise now. 

\bbox 
    Use the above to show that 
    \be 
    \label{eqn:VariationGenerators}
        \begin{split}
            \del_{P_{\mu}} & = -\p_{\mu} \\
            \del_{K_{\mu}} & = 2x_{\mu}\Delta - (x^2\p_{\mu} - 2x_{\mu}x^{\nu} \p_{\nu}) + 4(x_{\nu} {S^{\nu}}_{\mu}) \\
            \del_D & = -(\Delta + x^{\mu}\p_{\mu}\big) \\
            \del_{L_{\mu\nu}} & = x_{\nu}\p_{\mu} - x_{\mu}\p_{\nu} + S_{\nu\mu}
        \end{split}
    \ee 
\ebox 


\section{Descendants}

So far we have only discussed what we called primary fields. Of course we do not expect these to be the only types of fields in a CFT, and indeed they are not. There is another important type of field known as \textit{descendants}, however in order to show where they come from we need to introduce some subtle points.\footnote{The arguments made here are based off the ones made in Simmons-Duffin, [1602.07982]. So see there for more details.}

\subsection{Charges*}

Recall from a canonical QFT course that for every conserved Noether current we can construct a conserved charge. For example, spacetime translations, $P_{\mu} = \p_{\mu}$, give rise to the stress\footnote{It is likely I will change between saying stress/energy-momentum/stress-energy/other-similar-combinations tensor, so keep on your toes!} tensor $T^{\mu\nu}$. We can extend this to a more general statement. 

\bp 
    For a QFT (need not be a CFT) with well-defined stress tesnor, any Killing vector field --- that is a $\xi = \xi^{\mu}(x)\p_{\mu}$ obeying \Cref{eqn:KillingEquation} --- has a conserved charge given by 
    \be 
    \label{eqn:Charges}
        Q_{\xi}(\Sigma) = - \oint_{\Sigma} dS_{\mu} \, \xi_{\nu}(x) T^{\mu\nu}(x),
    \ee 
    where $\Sigma$ is the surface we integrate over with boundary $S$. 
\ep 

\bq 
    Simply take the derivative:\footnote{We drop all the integral etc to reduce notation.} 
    \bse 
        \begin{split}
            \p_{\mu} Q_{\xi} & \sim \p_{\mu} \big( \xi_{\nu}(x) T^{\mu\nu}(x) \big) \\
            & = \big(\p_{\mu}\xi_{\nu}(x)\big) T^{\mu\nu}(x) + \xi_{\nu}(x) \p_{\mu}T^{\mu\nu}(x) \\
            & = \frac{1}{2}\big(\p_{\mu}\xi_{\nu} + \p_{\nu}\xi_{\mu}\big) T^{\mu\nu} \\
            & = 0
        \end{split}
    \ese 
    where we have used $\p_{\mu}T^{\mu\nu}=0$ and the symmetry condition $T^{\mu\nu}=T^{\nu\mu}$. 
\eq 

As we said, this is true for QFTs outside CFTs. As we will show shortly, it turns out that in a CFT the stress tensor is traceless.\footnote{Foreshadowing to a remark that is to come, there are caveats to this.} It follows from this and the proof above that we can relax our condition to give the following Lemma. 

\bl 
    For a CFT a conformal Killing vector $\xi$ has an associated conserved charge given by \Cref{eqn:Charges}. 
\el 

\bbox 
    Prove this Lemma. \textit{Hint: Just take the above proof and recall \Cref{eqn:ConformalKilling}.}
\ebox 

\br 
    Note that \Cref{eqn:Charges} reduces to what we're familiar from in canonical QFT. There we take our slices to be equal time slices, so our boundary $dS_{\mu}$ points `in time' and our $\Sigma$ is a spatial slice. So our formula reduces to 
    \bse 
        Q_{\xi} = \int d^3 x \, \xi_{\nu}(x) T^{0\nu}(x) ,
    \ese 
    where the minus sign goes because we work in the signature $(-,+,+,+)$ in QFT.
\er 

Why are we talking about all this? Well it allows us to define conserved charges associated to our conformal Killing vectors $P_{\mu}, K_{\mu}, D$ and $L_{\mu\nu}$. We shall adopt the notation of using a tilde to indicate the corresponding conserved charge,\footnote{Note this is different to Simmons-Duffin's notation where a lower case letter indicates the vector field while a capital letter is used to denote the charge. We have already been using capital letters for the vector fields, so I've decided to adopt a tilde notation. Of course this is just notation and so does not mean anything itself, but this footnote is just for cross comparisons.} that is
\bse 
    \widetilde{P}_{\mu} := Q_{P_{\mu}}
\ese 
etc. However if we want to remain general (i.e. consider any of the conformal Killing vectors) we sill stick to the notation $Q_{\xi_i}$ where $i$ is meant to label whether we have $P_{\mu}$ or $K_{\mu}$ etc.

Now there is a highly non-intuitive result that we can show 

\bp 
    The commutators between our conformal charges satisfy 
    \be 
    \label{eqn:CommutatorOfChargesRelation}
        [Q_{\xi_i}, Q_{\xi_j}] = Q_{-[\xi_1,\xi_2]}.
    \ee 
\ep 

\bq 
    It is the minus sign that is highly non-trivial, so let's outline how you get it here. The first step is magically pull a relation out of the air 
    \bse 
        [Q_{\xi_i}, T^{\mu\nu}] = \xi_i^{\rho}\p_{\rho} T^{\mu\nu} + \p_{\rho}\xi_i^{\rho} T^{\mu\nu} - \p_{\rho}\xi_i^{\mu} T^{\rho\nu} + \p^{\nu} \xi_{i\rho} T^{\rho\mu}. 
    \ese 
    We do not prove this here but simply refer readers to Exercise 3.3. of Simmons-Duffin. Ok taking this formula as given let's try and prove \Cref{eqn:CommutatorOfChargesRelation}. The key point is to remember that the charges come as integrals and so the derivatives that will appear in the magic formula above will be with respect to different things. That is we have, for example,
    \bse 
        Q_{\xi_i} = -\oint dS_{\mu} \, \xi_{i\nu}(x) T^{\mu\nu}(x)
    \ese 
    and 
    \bse 
        Q_{\xi_j} = -\oint dS_{\mu} \, \xi_{j\nu}(y) T^{\mu\nu}(y). 
    \ese 
    Now we use the Jacobi identity (which our charges inherit from the Lie algebra) to write 
    \bse 
        \big[Q_{\xi_i}, [Q_{\xi_j}, T^{\mu\nu}]\big] - \big[Q_{\xi_j}, [Q_{\xi_i}, T^{\mu\nu}]\big] = \big[ [Q_{\xi_i},Q_{\xi_j}], T^{\mu\nu} \big]. 
    \ese 
    The idea is to expand the left-hand side out and then compare the result to our magic formula to deduce the commutation relation. Let's consider just the first term: we do the inner commutator first to give us 
    \bse 
        \begin{split}
            \big[Q_{\xi_i}, [Q_{\xi_j}, T^{\mu\nu}]\big] & = \big[ Q_{\xi_i} , \xi_j^{\rho}\p_{\rho} T^{\mu\nu} + \p_{\rho}\xi_j^{\rho} T^{\mu\nu} - \p_{\rho}\xi_j^{\mu} T^{\rho\nu} + \p^{\nu} \xi_{j\rho} T^{\rho\mu} \big] \\
            & = \xi_j^{\rho}\p_{\rho} [ Q_{\xi_i} , T^{\mu\nu}] + \p_{\rho}\xi_j^{\rho} [ Q_{\xi_i} , T^{\mu\nu} ] - \p_{\rho}\xi_j^{\mu} [ Q_{\xi_i} , T^{\rho\nu}] + \p^{\nu} \xi_{j\rho} [ Q_{\xi_i} , T^{\rho\mu}],
        \end{split}
    \ese 
    where the second line follows from the linearity of the commutator bracket and the comment we made above about the derivatives being w.r.t. different variables. To avoid any confusion, what we mean is that the derivatives appearing in the above expression are w.r.t $y$ whereas $Q_{\xi_i}$ is a function of $x$ and so we can freely take the derivatives outside. 
    
    Hopefully it is clear from here\footnote{As otherwise it's a lot of writing for me...} how you continue the calculation and how it results in a left-hand side which implies \Cref{eqn:CommutatorOfChargesRelation}, so we finish the proof here.
\eq 

\br 
    It turns out that the magic formula we quoted above is only true in $d\geq 3$, so again our $2$-dimensional CFT is special. However for what we're going to use these results for here this doesn't matter as we have another way to deal with it in $2$-dimensions.
\er 

We now introduce a new notation 
\be
\label{eqn:CommutatorOfChargesRelationNew}
    [Q_{\xi_i}, Q_{\xi_j}] = - \widetilde{[\xi_i, \xi_j]}
\ee 
so that, for example, 
\be 
\label{eqn:CommutatorKDCharges}
    [\widetilde{K}_{\mu}, \widetilde{D}] = - \widetilde{[K_{\mu}, D]} = \widetilde{K}_{\mu}. 
\ee 
Again this is a highly non-trivial result and essentially corresponds to saying that we have to swap the signs everywhere in our previous generator commutation relations, \Cref{eqn:GeneratorCommutators}. 

\mybox{
    \begin{center}
        Taking the commutators of charges (tilded letters) differs from the commutators of vector fields (no tilde) by a minus sign, as in \Cref{eqn:CommutatorKDCharges}. We therefore have to put minus signs on the right-hand side of \Cref{eqn:GeneratorCommutators} when considering the charges.
    \end{center}
}

\subsection{Descendant Fields*}
\label{sec:DescendantFields}

Ok why are we bothering to do all of this? Well it allows us to get a different definition of a primary field, which in turn will allow us to define descendant fields. The first thing we note is that if we take our local operator to be at $x=0$, then our dilatation operator acts as 
\bse 
    D \phi(0) = -\Delta \phi(0),
\ese 
which is easily seen from \Cref{eqn:VariationGenerators}. We are ultimately interested in the quantum theory where the local fields $\phi$ become local \textit{operators} $\cO_{\phi}$. The action of a charge on the local operator is given by the commutator, e.g.
\bse 
    [\widetilde{D}, \cO_{\phi}(0)] = \Delta \cO_{\phi}(0).
\ese 
We now adopt a new notation for the action of the charges on the fields by simply dropping the commutator brackets, i.e. we define
\bse 
    Q_{\xi_i}\cO(x) := [Q_{\xi_i}, \cO(x)].
\ese 
We extend this to nested commutators, 
\bse 
    Q_{\xi_j}Q_{\xi_i}\cO(x) := \big[Q_{\xi_j} , [Q_{\xi_i}, \cO(x)]\big].
\ese 

Now using this notation we can finally obtain the result we've been driving at. Consider the following (the subscript $\Delta$ is to label the weight of our operator)
\bse 
    \begin{split}
        \widetilde{D}\widetilde{K}_{\mu}\cO_{\Delta}(0) & = \big( \widetilde{K}_{\mu}\widetilde{D} + [\widetilde{D},\widetilde{K}_{\mu}]\big) \cO_{\Delta}(0) \\
        & = \big(\Delta - 1\big) \widetilde{K}_{\mu}\cO_{\Delta}(0),
    \end{split}
\ese 
where we have made use of \Cref{eqn:CommutatorKDCharges}. This tells us that $\widetilde{K}_{\mu}\cO_{\Delta}$ has dilatation weight $(\Delta-1)$. In other words we can view $\widetilde{K}_{\mu}$ as a lowering operator for the weight. Now any physically reasonable theory will have a lower bound on the dilatation weight of a field/operator, and so it follows that there \textit{must} exist an operator such that 
\bse 
    \widetilde{D}\cO(0) = 0.
\ese 
This is what we can take as the definition of a primary operator, which we can relate to a definition of a primary field which we state now. 

\bd[Primary Operator]
    A primary operator of a given representation is an operator with the lowest dilatation weight in that representation of the conformal algebra. 
\ed 

Armed with this definition we can (finally!) explain descendant fields. Recall that the dilation weight is equal to the mass dimension. What we therefore want is something that raises the mass dimension of a field. Well the partial derivative has $[\p]=+1$ and so it follows that the dilatation weight of $\p\phi$ is $(\Delta+1)$. We generate derivatives using the momentum operator, and so we can define our descendant operators accordingly. 

\bd[Descendant Operator]
    Let $\cO_{\Delta}$ be a local primary operator of weight $\Delta$. Then the set 
    \bse 
        D_{\cO} := \{ \widetilde{P}_{\mu_1} \widetilde{P}_{\mu_2} ... \widetilde{P}_{\mu_n} \cO_{\Delta} \, | \, n \in \N \}
    \ese 
    are the \textit{descendant} local operators associated to $\cO$. A sum of such terms is also a descendant operator.
\ed 

We can of course translate this into a definition of descendant fields via the action of the un-tilded $P_{\mu}$s on $\phi(x)$. 

\mybox{
    \begin{center}
        A descendant field is given by a linear combination of primary fields and their derivatives.
    \end{center}
}

\subsection{Examples Of Fields}

Let's give some examples of primary and descendant fields now. We shall use a free, complex, massless scalar field theory, 
\bse 
    S = \int d^d x \,  \p_{\mu}\phi \p^{\mu}\bar{\phi}
\ese 
as our starting point. Our discussions above tell us that:

\begin{center}
	\begin{tabular}{@{} C{5cm} C{5cm} C{4cm} @{}}
		\toprule
		Field & Primary Or Descendant & Weight \\
		\midrule 
		$\phi(x)$ & Primary & $\Delta$ \\ \\
		$\phi^n(x)$ &  Primary & $n\Delta$ \\ \\
		$\p_{\mu}\phi(x)$ & Descendant & $\Delta+1$ \\ \\
		$\phi(x)\p_{\mu}\phi(x)$ & Descendant & $2\Delta+1$ \\ \\
		$\phi(x)\p_{\mu}\bar{\phi}(x) + \bar{\phi}(x)\p_{\mu}\phi(x)$ & Descendant & $2\Delta+1$ \\ \\
		$\phi(x)\p_{\mu}\bar{\phi}(x) - \bar{\phi}(x)\p_{\mu}\phi(x)$ & Descendant & $2\Delta+1$ \\ 
		\bottomrule
	\end{tabular}
\end{center}

\br 
    Note that the 4th and 5th examples are descendants as they are, respectively, 
    \bse 
        \p_{\mu}\phi^2(x) \qand \p_{\mu}\big(\bar{\phi}(x)\phi(x)\big)
    \ese 
\er 

\bbox 
    Prove the above primaries are indeed primary and check their weights. That is check that $\del_{K_{\nu}}$ and $\del_D$ are correct for primary. For the last one you will need (for this case) that.
    \bse 
        {({S^{\nu}}_{\mu})_{\rho}}^{\sig} = \frac{1}{2}\big( - \del^{\nu}_{\rho} \del^{\sig}_{\mu} + \del^{\nu}_{\mu}\del_{\rho}^{\sig}\big)
    \ese
    \textit{Hint: You start by assuming that $\phi(x)$ and $\overline{\phi}(x)$ is a primary. To be clear, we show $\del_D$ for $\phi^n(x)$ here:}
    \bse 
        \del_D\big(\phi^n(x)\big)  = n \phi^{n-1}\del_D \phi(x) = -n\phi^{n-1}\big(\Delta + x^{\mu}\p_{\mu}\big)\phi(x) = -\big(n\Delta + x^{\mu}\p_{\mu}\big)\phi^n(x),
    \ese 
    \textit{which is an operator of dimension $n\Delta$.}
\ebox 

\section{Stress-Energy Tensor}

Now we have already made serious use of the stress-energy tensor above in deriving descendant operators, but we haven't actually talked about the existence of $T^{\mu\nu}$ in a CFT. It turns out that \textit{every} CFT has a well defined stress-energy tensor. Indeed we have seen that the stress-energy tensor actually generates the conformal transformations themselves via \Cref{eqn:Charges}.

So how do we define the stress-energy tensor in a CFT? Well if we have a theory with a Lagrangian, we start by coupling the theory to gravity, which is accomplished via 
\bse 
    S = \int d^4x (\p\phi)^2 \to \int d^4x \sqrt{-g} \nabla_{\mu}\phi \nabla_{\nu}\phi g^{\mu\nu},
\ese
which should be familiar from a GR course. We then proceed exactly as in GR: we take the variation w.r.t. the metric $g_{\mu\nu}$ to obtain
\bse 
    \del S = -\frac{1}{2} \int d^Dx \, \sqrt{-g} T^{\mu\nu} \del g_{\mu\nu},  
\ese 
or 
\mybox{
    \be 
    \label{eqn:StressEnergyTensor}
        T^{\mu\nu} := -\frac{2}{\sqrt{-g}} \frac{\del S}{\del g_{\mu\nu}}
    \ee 
}

What properties does the stress-energy tensor have? 
\ben 
    \item Because it is coupled to gravity, we have diffeomorphism invariance\footnote{Namely $\del S = 0$ under $\del g_{\mu\nu} = \p_{(\mu}\epsilon_{\nu)}$. We can see this is a translation by plugging $x^{\prime\mu} = x^{\mu} + \epsilon^{\mu}$ into \Cref{eqn:ConformalTransformation}.} which implies $\p_{\mu}T^{\mu\nu} = 0$, which is a conservation equation. Note this is a specialisation to flat space, as otherwise we need a covariant derivative. However it is true for any theory in flat space.
    \item We also have Weyl invariance, which tells us that $\del S = 0$ under $\del g_{\mu\nu} = \kappa \eta_{\mu\nu}$ which implies $T^{\mu\nu}\eta_{\mu\nu} = {T^{\mu}}_{\mu} = 0$, so it is traceless. This is \textit{not} true for a non-CFT, so we can sort of see this as a defining property of a CFT.
    \item We can use the stress energy tensor to construct \textit{all} the Noether currents associated with conformal symmetries. In this sense we can say that the stress-energy tensor generates our conformal field theories.
\een 

\badr
\label{rem:WeylAnom}
    Condition 2 is only true classically. It turns out that in even dimensional CFTs, upon quantisation the trace of the stress-energy tensor is some constant times terms that depend on the curvature, e.g. the Ricci scalar. These are known as \textit{Weyl anomalies}. We are working in flat space in this course, though, so these Weyl anomolies will not bother us. 
\eadr 
\chapter{Ward Identity \& Correlation Functions}

% Note to Paul: To be honest I find the Ward identity very confusing. I have seen multiple different statements of it which do not seem to be trivially compatible to me. I have written here what I think is true and have tried to work off Simmons-Duffin, but it's likely I've connected dots that shouldn't be connected. 

So far we have been able to categorise all the fields (and their associated operators) in our $d>2$ CFT. However we know that in QFT it is important to also know about the \textit{correlation functions} for the fields. As we said earlier, if we know all the operators and all the corresponding correlators then we say was have `solved' the theory, in the sense that we can use that information to obtain anything we want to know about the CFT. Before looking at the correlators, let's just recap what we have seen as it will be useful going forward.

\section{Ward Identities*}

We have seen that in the quantum theory, where our local fields become local operators acting on some Hilbert space, and we can implement the conformal transformations via commutators of the charges and operators. We recall that when we do this we pick up an additional minus sign compared to the differential operator case. In other words, we have
\bse 
    [Q_{\xi_i},\cO(x)] = - \del_{\xi_i} \cO(x)
\ese 
where the $\del_{\xi_i}$s are given by \Cref{eqn:VariationGenerators}. For example we have 
\bse 
    [\widetilde{D},\cO(x)] = (\Delta + x^{\mu}\p_{\mu}) \cO(x) \qand [\widetilde{P}_{\mu}, \cO(x)] = \p_{\mu}\cO(x),
\ese 
etc. 

Why are we repeating this? Well now let's consider a correlation function with our charges inserted. For concreteness we shall use $\widetilde{P}^{\mu}$
\bse 
    \la \widetilde{P}_{\mu}(\Sigma) \cO(x) \ra, 
\ese 
where we have explicitly reinserted the boundary $\Sigma$. Now correlation functions come time ordered, which we need to account for. Now our $\Sigma$ is the boundary of some surface in our spacetime, and so we really need to account for the fact that it `spans some time period'. Let's take $\Sigma_1, \Sigma_2$ to be spatial slices at times $t_1 < t < t_2$, where $t=x^0$ is the time of our local operator. Now with \Cref{eqn:Charges} in mind we take the orientation of these spatial surfaces to be opposite; that is our charges are the boundary of some surface and the $dS^{\mu}$ appearing in \Cref{eqn:Charges} always points out of this surface. Another way to think of this is to imagine picking $\Sigma$ so that it stretches really far in the spatial direction, then in the local picture it gives spatial slices with different orientations, as illustrated below. 

\begin{center}
    \btik 
        \begin{scope}[xshift = -4cm]
            \draw[dashed, thick] (0,0) ellipse (3cm and 1cm);
            \draw[thick, ->] (0,-1) -- (0,-1.5);
            \draw[thick, ->] (1,-0.9) -- (1,-1.4);
            \draw[thick, ->] (-1,-0.9) -- (-1,-1.4);
            \draw[thick, ->] (2,-0.75) -- (2,-1.25);
            \draw[thick, ->] (-2,-0.75) -- (-2,-1.25);
            % 
            \draw[thick, ->] (0,1) -- (0,1.5);
            \draw[thick, ->] (1,0.9) -- (1,1.4);
            \draw[thick, ->] (-1,0.9) -- (-1,1.4);
            \draw[thick, ->] (2,0.75) -- (2,1.25);
            \draw[thick, ->] (-2,0.75) -- (-2,1.25);
            % 
            \node at (0,0) {$\cO(x)$};
        \end{scope}
        \draw[ultra thick, ->] (0,0) -- (2,0) node [above,midway] {Local};
        \begin{scope}[xshift=4cm]
            \draw[dashed, thick] (-1,1) -- (1,1);
            \draw[dashed, thick] (-1,-1) -- (1,-1);
            \draw[thick, ->] (-0.5,1) -- (-0.5,1.5);
            \draw[thick, ->] (0.5,1) -- (0.5,1.5);
            \draw[thick, ->] (-0.5,-1) -- (-0.5,-1.5);
            \draw[thick, ->] (0.5,-1) -- (0.5,-1.5);
            % 
            \node at (0,0) {$\cO(x)$};
        \end{scope}
    \etik 
\end{center}

We therefore have two terms 
\bse 
    \la \cT\big\{ \big(\widetilde{P}_{\mu}(\Sigma_2) - \widetilde{P}_{\mu}(\Sigma_1)\big)\cO(x)\big\} \ra 
\ese 
which using the time ordering gives us a commutator
\bse 
    \la [\widetilde{P}_{\mu}, \cO(x)] \ra = \la \p_{\mu} \cO(x) \ra = \p_{\mu} \la \cO(x) \ra.
\ese 
It is important to notice that this result comes from the fact that $t_1<t<t_2$, or equivalently that $\cO(x)$ lies within our $\Sigma$. It is hopefully clear that we will get this result regardless of how we choose our spatial slices; that is $\cO(x)$ is always contained within the boundary and so we will always get a commutator. 

We can now extend this argument to the case when we have more then one operator inserted within our $\Sigma$. Here we will get the sum of commutators. That is, let's assume that the first $i$ operators, $\{\cO_1(x_1) , ... , \cO_i(x_i)\}$, lie within $\Sigma$, then we get
\bse 
    \begin{split}
        \la \cT\{ \widetilde{P}^{\mu}(\Sigma) \cO_1(x_1) ... \cO_i(x_i) ... \cO_n(x_n) \} \ra & = \la \big[\widetilde{P}^{\mu}(\Sigma), \cO_1(x_1) ... \cO_i(x_i)\big] \cO_{i+1} ... \cO_n(x_n) \} \ra \\
        & = \la \Big( \big[\widetilde{P}^{\mu}(\Sigma), \cO_1(x_1)\big] \cO_2(x_2) ... \cO_i(x_i) \\
        & \quad + \cO_1(x_1) \big[\widetilde{P}^{\mu}(\Sigma), \cO_2(x_2)\big] \cO_3(x_3) ... \cO_i(x_i) + ... \\
        & \quad + \cO_1(x_1)... \cO_{i-1}(x_{i-1}) \big[\widetilde{P}^{\mu}(\Sigma), \cO_i(x_i)\big]\Big) \cO_{i+1} ... \cO_n(x_n) \} \ra \\ 
        & = \big(\p_1^{\mu} + ... + \p^{\mu}_i \big) \la \cO_1(x_1) ... \cO_n(x_n) \ra,  
    \end{split}
\ese 
where we have used the commutator result 
\bse 
    [A,BC] = [A,B]C + B[A,C]
\ese 
repeatedly to get the sums on the second line. The derivatives are meant to be understood with the mantra that they only act on the relevant fields, i.e. $\p^{\nu}_i$ only acts on $\cO_i(x_i)$ and none of the others. For even more clarity, we could imagine labelling the different $x_i$s as $\{x,y,z,...\}$ and then our derivatives would be $\frac{\p}{\p x}$, $\frac{\p}{\p y}$ etc.  

Now if we recall that our charges are given by integrals over the stress-energy tensor we can reverse engineer\footnote{I say reverse engineer because Simmons-Duffin obtains the above result from the Ward identity. } this result to arrive at the Ward identity. 

\be 
\label{eqn:WardIdentity}
    \p_{\mu}\la T^{\mu\nu}(x) \cO_1(x_1) ... \cO_n(x_n) \ra = - \sum_i \del^{(4)}(x-x_i) \p_i^{\nu} \la \cO_1(x_1) ... \cO_n(x_n) \ra.
\ee 

Let's just check this makes sense. Integrating both sides over the region $\Sigma$ is the boundary of: the derivative on the left-hand side can be removed using Stoke's theorem, leaving us with the integral over $\Sigma$ as needed for the charge; on the right-hand side the delta functions give us a sum of partial actions. The minus sign on the right-hand side is just included to account for the minus in \Cref{eqn:Charges}.

The above was derived using the momentum charge, but we can extend it to the others by simply including the $\xi_{\nu}$ contractions, i.e. 
\bse 
    \p_{\mu}\la \xi_{\nu}(x)T^{\mu\nu}(x) \cO_1(x_1) ... \cO_n(x_n) \ra = - \sum_i \del^{(4)}(x-x_i) \p_i^{\nu} \la \xi_{\nu}(x) \cO_1(x_1) ... \cO_n(x_n) \ra,
\ese 
which we can rewrite as 
\be 
\label{eqn:WardIdentityCharge}
    \p_{\mu}\la \xi_{\nu}(x) T^{\mu\nu}(x) \cO_1(x_1) ... \cO_n(x_n) \ra = \sum_i \del^{(4)}(x-x_i) \la \cO_1(x_1) ... \del_{\xi}\cO_i(x_i) ... \cO_n(x_n) \ra.
\ee 

We did all this so that we could draw the following conclusion: in terms of transformations of correlation functions, all that matters is whether the operator is contained within the $\Sigma$ region or not. This corresponds to the current (i.e. $T^{\mu\nu}$) having common support to the local operators, when it does the delta function in \Cref{eqn:WardIdentity} is hit and we pick up a term. Such terms are known as \textit{contact terms}, for intuitive reasons. 

There's a nice result we can get from this. Let's consider the case when $\Sigma$ encloses \textit{all} the operators. Then our time ordering gives us the commutator as 
\bse 
    \begin{split}
        \la \cT\{ Q_{\xi_i}(\Sigma) \cO_1(x_1) ... \cO_n(x_n) \} \ra & = \la [Q_{\xi_i}(\Sigma), \cO_1(x_1)...\cO_n(x_n)] \ra \\
        & = \la Q_{\xi_i}(\Sigma_2)\cO_1(x_1)...\cO_n(x_n) \ra - \la \cO_1(x_1)...\cO_n(x_n)Q_{\xi_i}(\Sigma_1) \ra. 
    \end{split}
\ese 
Now it turns out to be true\footnote{For a reason that will become clearer later.} that all the conformal charges annihilate the vacuum both as a left and right action, i.e. $\bra{0}Q_{\xi_i} = 0 = Q_{\xi_i}\ket{0}$, and so the above result just vanishes. However we could equally expand the commutator out as 
\bse 
    \begin{split}
        [Q_{\xi_i}(\Sigma), \cO_1(x_1) ... \cO_n(x_n)] & = [Q_{\xi_i}(\Sigma),\cO_1(x_1)]\cO_2(x_2) ... \cO_n(x_n) \\
        & \qquad + \cO_1[Q_{\xi_i}(\Sigma),\cO_2(x_2)]\cO_3(x_3) ... \cO_n(x_n) + ...\\
        & \qquad \qquad  + \cO_1(x_1) ... \cO_{n-1}(x_{n-1})[Q_{\xi_i}(\Sigma),\cO_n(x_n)],
    \end{split}
\ese 
and then using $[Q_{\xi_i}(\Sigma), \cO_j(x_j)] = -\del_{\xi_i} \cO(x_j)$, we get 
\bse 
    \la \del\cO_1(x_1) \cO_2(x_2) ... \cO_n(x_n)\ra + ... + \la \cO_1(x_1) \cO_2(x_2) ... \del\cO_n(x_n)\ra = 0.
\ese 
We can also write this result as the following.
\mybox{

    If $x\mapsto x'$ is a quantum symmetry with $\cO \mapsto \cO'$, then the correlators of the transformed operators is equal to the correlator of the untransformed ones. That is 
    \be
    \label{eqn:WardIdentityCourse}
        \la \cO_1(x_1) ... \cO_n(x_n) \ra = \la \cO_1^{\prime}(x_1) ... \cO_n^{\prime}(x_n) \ra,
    \ee 
    where we note that the $x$s on the right-hand side are \textit{not} primed, as otherwise we have a trivial statement from $\cO(x)=\cO^{\prime}(x^{\prime})$.
}
\noindent where we get the above result by considering $\cO^{\prime} = \cO + \del\cO$. 

\section{Correlation Functions}

As we've said a few times, the key objects in a CFT are the local operators, and the key data defining the CFT are the correlation functions of these local operators. For example
\bse 
    \la \cO(x_1) \cO(x_2) ... \cO(x_n)\ra = f(x_1,...,x_n)
\ese
where $f(x_1,...,x_n)\in \C$ for a given $(x_1,...,x_n)$. The important thing to note is that the result is a function --- i.e. there is no information left about the operators themselves on the right-hand side. This is why we say that once we know the correlation functions we have solved the theory --- we have `removed' all the information about the operators themselves and are just left with theory specific results. 

In principle, if you have a Lagrangian you can compute all the correlation functions essentially via Feynman diagrams. However, in practice this is exceedingly difficult!\footnote{Besides that, what about theories that don't have a Lagrangian?} We therefore want to try find some other way to get the correlation function results. Given that we just worked out how conformal transformations act on the correlation functions, \Cref{eqn:WardIdentityCourse}, the obvious question to ask is "does our conformal symmetry tells us anything about them?" The answer is "yes", and we shall now see what they tell us. 

\subsection{Conformal Constraints}

% Commenting out for now as doesn't use language of course, but don't want to loose the information.

% Now it follows from our derivation of \Cref{eqn:WardIdentityCharge} that if the charge has no common support with any operator (i.e. all the charges are outside $\Sigma$) that the right-hand side vanishes. Now let's consider the case when the charge has common support with \textit{every} operator (i.e. they're all within $\Sigma$). Then every delta function is hit, and so we pick up every $\del_{\xi}\cO(x_i)$. However we know that in a conformal symmetry we have inversion symmetry, which essentially corresponds to turning the inside of the $\Sigma$ to the outside and visa versa. This has the effect of resulting in a situation where all the operators lie \textit{outside} $\Sigma$ and so the result must vanish. Equating these two ideas allows us to conclude 
%\be
%\label{eqn:CorrelationFunctionVariationOperatorsXi}
%    \la \del_{\xi_i}\cO_1(x_1) \cO_2(x_2) ... \cO_n(x_n)\ra + \la \cO_1(x_1) \cO_2(x_2) ... \del_{\xi_i}\cO_n(x_n)\ra = 0.
%\ee 

Let's consider a infinitesimal conformal transformation. We then have 
\bse 
    \cO^{\prime}(x) = \cO(x) + \del\cO.
\ese 
Plugging this into \Cref{eqn:WardIdentityCourse} we see that the left-hand side will be cancelled by the $\cO(x)$ on the right-hand side of the above formula. We are then just left with 
\be 
\label{eqn:CorrelationFunctionVariationOperators}
    \la \del\cO_1(x_1) \cO_2(x_2) ... \cO_n(x_n)\ra + ... + \la \cO_1(x_1) \cO_2(x_2) ... \del\cO_n(x_n)\ra = 0.
\ee 
We therefore have\footnote{Note the signs flip compared to \Cref{eqn:VariationGenerators}, due to the fact we are considering the charges. Of course we're setting it equal to 0 so an overall sign makes no difference.} 
\ben[label=(\roman*)]
    \item Translations, $\del=\del_{\widetilde{P}_{\mu}} = \p_{\mu}$:
    \bse 
        \bigg(\frac{\p}{\p x_1^{\mu}} + \frac{\p}{\p x_2^{\mu}} + ... + \frac{\p}{\p x_n^{\mu}}\bigg) \la \cO_1(x_1) ... \cO_n(x_n)\ra = 0.
    \ese
    \item Dilatations, $\del = \del_{\widetilde{D}} = (x^{\mu}\p_{\mu} + \Delta)$:
    \bse 
        \big(x_1^{\mu}\p_{1\mu} + x_2^{\mu}\p_{2\mu} + ... + x_3\p_{3\mu} + \Delta_1 + \Delta_2 + ... \Delta_n\big) \la \cO_1(x_1) ... \cO_n(x_n)\ra = 0.
    \ese
    \item Special conformal, for primary scalar operators, $\del = \del_{\widetilde{K}_{\mu}} = -2x_{\mu}\Delta +(x^2\p_{\mu} - 2x_{\mu} x^{\nu}\p_{\nu}) $:
    \bse 
        \bigg[\sum_{i=1}^n \big(2x_{i\mu} \Delta_i -(x_i^2 \p_{i\mu} - 2x_{i\mu} x_i^{\nu}\p_{i\nu})\big)\bigg] \la \cO_1(x_1) ... \cO_n(x_n)\ra = 0.
    \ese
    We can simplify this by defining\footnote{Note, no tilde here. This is the vector field $K_{\mu}$.} 
    \bse 
        K_{\mu} = x_i^2 \p_{i\mu} - 2x_{i\mu} x_i^{\nu}\p_{i\nu},
    \ese 
    which is the generator of special conformal transformations at $x_i$.
    \item Lorentz, for primary scalar operators, $\del = \del_{\widetilde{L}_{\mu\nu}} = -x_{\nu}\p_{\mu} + x_{\mu}\p_{\nu}$:
    \bse 
        \bigg[\sum_{i=1}^n \big(x_{i\mu}\p_{i\nu} - x_{i\nu}\p_{i\mu}\big)\bigg]\la \cO_1(x_1) ... \cO_n(x_n) \ra = 0.
    \ese 
\een

\br 
    We can, of course, adapt the special conformal and Lorentz ones to non-scalars by including the $\rho$ terms, which results in including $S_{\mu\nu}$ factors. 
\er 

We can use these relations to constrain the correlators, which we now do. 

\subsection{2-Point Functions}

As always we start with the simplest case, the $2$-point function:
\bse 
    \la \cO_1(x_1)\cO_2(x_2)\ra = f(x_1,x_2).
\ese 
Here we will assume scalar operators, but not a great deal changes for other types; you just need to consider the Lorentz representation too.  

\subsubsection{Translations}

First we impose translation symmetry,
\bse 
    (\p_1+\p_2)f(x_1^{\mu},x_2^{\mu}) = 0 \qquad \implies \qquad f(x_1^{\mu},x_2^{\mu}) = f\big((x_1-x_2)^{\mu}\big),
\ese 
where the implication arrow can be seen by taking a change of variables as follows: let $y_1=x_1-x_2$ and $y_2=x_2$, then by chain rule we have (suppressing indices)
\bse 
    \p_1 f(y_1,y_2) = \frac{\p f}{\p y_1}\p_1 y_1 + \frac{\p f}{\p y_2}\p_1 y_2 = \frac{\p f}{\p y_1}
\ese 
and 
\bse 
    \p_2 f(y_1,y_2) = \frac{\p f}{\p y_1}\p_2 y_1 + \frac{\p f}{\p y_2}\p_2 y_2 = -\frac{\p f}{\p y_1} + \frac{\p f}{\p y_2},
\ese 
adding these together gives us
\bse 
    \frac{\p }{\p y_2} f(y_1,y_2) = 0 \qquad \implies \qquad f(y_1,y_2) = f(y_1) = f(x_1-x_2).
\ese 

\badr 
    Note this result makes perfect sense if we think about the problem in terms of our $\Sigma$ picture. In general the operators will be inserted at $x_1 \neq 0 \neq x_2$, however we can use a translation to move everything such that $x_1=0$, say. It is clear from this that the answer could then only depend on $x_2$. We could equally have shifted everything so that $x_2=0$, and so the result then only depends on $x_1$. It's not a huge leap to go from here to seeing that the result only depends on the difference between the two points. 
\eadr 

\bnn 
    From now one we shall use the notation $x_{12} := x_1-x_2$. So we write the above as $f(x_{12}^{\mu})$.
\enn 

\subsubsection{Lorentz}

We could proceed similarly to above to see what the Lorentz transformations tells us. However we can save a bit of time by using the idea of "Lorentz symmetry corresponds to not having indices left over". We have seen that $f$ is a function of $x_{12}^{\mu}$ and so if we want to remove the indices we have to make it a function of $x_{12}^2 := x_{12}^{\mu}(x_{12})_{\mu}$. 

\badr 
    As with the remark above, we can make a diagrammatic argument here: Lorentz symmetries contains spatial rotations and boosts. The rotations tell us that our $f(x_{12}^{\mu})$ can't depend on how $\Vec{x}_{12}$ is orientated relative to the coordinate axes. It follows from this that it must only depend on the absolute value $x_{12}^2$.
\eadr 

So we have 
\bse 
    f(x_1,x_2) = f(x_{12}^2) \equiv f(|x_1-x_2|^2).
\ese    

\subsubsection{Dilatations}

Let's now look at something that is CFT specific.\footnote{The Ward identity formula holds for a generic QFT and so the above results also hold there.} 
\bse 
    \big(x_1\cdot \p_1 + x_2\cdot \p_2 + \Delta_1 + \Delta_2\big) f(x_{12}^2) = 0.
\ese 
Using the chain rule 
\bse 
    \p_{1\mu} f(x_{12}^2) = 2 x_{12\mu} f'(x_{12}^2) \qand \p_{2\mu} f(x_{12}^2) = -2 x_{12\mu} f'(x_{12}^2),
\ese 
we have 
\bse 
    0 = \big(2x_1\cdot x_{12} - 2x_2\cdot x_{12}^2\big) f'(x_{12}^2) + \big(\Delta_1+\Delta_2)f(x_{12}^2) = 2x_{12}^2 f'(x_{12}^2) + (\Delta_1+\Delta_2)f(x_{12}^2).
\ese 
We can solve this using separation of variables: let $x_{12}^2=x$ for notational reasons, then we have 
\bse 
    \ln f = A + \frac{\Delta_1+\Delta_2}{2} \ln x \qquad \implies \qquad f = \frac{A^{\prime}}{x^{(\Delta_1+\Delta_2)/2}}
\ese
which we can write as 
\bse 
    \la \cO_1(x_1)\cO_2(x_2) \ra = \frac{C_{12}}{|x_1-x_2|^{\Delta_1+\Delta_2}}.
\ese

\subsubsection{Special Conformal}

Finally we have special conformal transformations. These give us 
\bse 
    \big(-2x_{1\mu}\Delta_1 - 2x_{2\mu}\Delta_2 + K_{1\mu}+K_{2\mu}\big)f(x_{12}^2) = 0.
\ese 

\bcl 
    The following equation holds 
    \bse 
        (K_{1\mu}+K_{2\mu}) |x_{12}| = -(x_{1\mu}+x_{2\mu}) |x_{12}|.
    \ese 
\ecl 

\bq 
    \bse 
        \begin{split}
            (K_{1\mu}+K_{2\mu}) x_{12}^2 & = \big( x_1^2\p_{1\mu} - 2x_{1\mu}x_1^{\nu}\p_{1\nu} + x_2^2\p_{2\mu} - 2x_{2\mu}x_2^{\nu}\p_{2\nu}\big) x_{12}^2 \\
            & = \big(2x_1^2(x_1-x_2)_{\mu} - 4x_{1\mu}x_1^{\nu}(x_1-x_2)_{\nu} - 2x_2^2(x_1-x_2)_{\mu} + 4x_{2\mu}x_2^\nu(x_1-x_2)_{\nu} \\
            & = -2x_1^2 x_{1\mu} - 2 x_1^2x_{2\mu} + 4x_{1\mu}x_1\cdot x_2 - 2x_2^2x_{2\mu} - 2x_2^2x_{1\mu} + 4x_{2\mu}x_1\cdot x_2 \\
            & = -2(x_{1\mu}+x_{2\mu})(x_1^2 - 2x_1\cdot x_2 + x_2^2) \\
            & = -2(x_{1\mu}+x_{2\mu})(x_1-x_2)^2.
        \end{split}
    \ese
    Then using the fact that $K_{\mu}$ is a differential operator, if we act on the square root, we just pull down a factor of a $1/2$, which gives us 
    \bse 
        (K_{1\mu}+K_{2\mu}) (x_{12}^2)^{1/2} = \frac{1}{2(x_{12}^2)^{1/2}} (K_{1\mu}+K_{2\mu}) x_{12}^2
    \ese 
    which is the result we want. 
\eq 

\bbox 
Using this result show that our special conformal transformations give 
\bse 
    \big(-2x_{1\mu}\Delta_1 - 2x_{2\mu}\Delta_2 + K_{1\mu}+K_{2\mu}\big)\frac{C_{12}}{|x_{12}|^{\Delta_1+\Delta_2}} = (\Delta_2-\Delta_1)(x_{1\mu} - x_{2\mu}) \frac{C_{12}}{|x_{12}|^{\Delta_1+\Delta_2}} = 0
\ese 
\ebox 

The conclusion is that we either have $C_{12}=0$ or $\Delta_1=\Delta_2$.\footnote{We don't want to have $x_1=x_2$ here as this would correspond to putting two operators on top of each other, which is not a nice idea.} To conclude we have 
\mybox{
    \be 
    \label{eqn:TwoPointCorrelator}
        \la \cO_1(x_1) \cO_2(x_2)\ra = \begin{cases} 
            \frac{C_{12}}{|x_1-x_2|^{2\Delta}} & \text{ if } \Delta_1=\Delta_2 \\
            0 & \text{ otherwise}.
        \end{cases}
    \ee 
}

So we have seen that conformal symmetries fix the 2-point correlators up to a factor $C_{12}$. We can even fix this constant by renormalising our fields (and therefore operators) via 
\bse 
    \cO \to \frac{1}{\sqrt{C_{12}}} \cO,
\ese 
which puts a $1$ in the numerator of \Cref{eqn:TwoPointCorrelator}. This tells us that for a CFT, defining the space of operators and their dimensions is equivalent to the space of 2 point functions.

% Note for Paul: You made some comment about getting anomalous dimension running with the coupling. I didn't understand this as I am not overly familiar with RG, so haven't made any remark here. 

\subsection{3-Point Functions Of Scalar Operators}

That's two-point functions done, next on the list is three-point functions
\bse 
    \la \cO_1(x_1) \cO_2(x_2) \cO_3(x_3) \ra = f(x_1,x_2,x_3).
\ese 
We proceed similarly to find
\mybox{
    \be 
    \label{eqn:ThreePointCorrelator}
        \la \cO_1(x_1)\cO_2(x_2)\cO_3(x_3)\ra = \frac{C_{123}}{|x_{12}|^{\Delta_1+\Delta_2-\Delta_3} |x_{23}|^{\Delta_2 + \Delta_3 - \Delta_1} |x_{13}|^{\Delta_3 + \Delta_1 - \Delta_2} }
    \ee
}

\br 
    Note that once we have fixed our $C_{12}$ for the two point functions we cannot change the $C_{123}$. That is the three-point functions can only be defined up to $C_{123}$.
\er 

\bbox 
    Derive \Cref{eqn:ThreePointCorrelator}.
\ebox 

\subsection{4-Point Functions}

Next, 4-point functions. Unfortunately, these are no longer fixed by conformal symmetry. This is not surprising as we expect there to be some point at which our conformal transformations stop helping us --- that is we have a limited number of constraints we can put on the system and so there has to be a point at which we cannot fully constrain the system anymore. 

The result of imposing conformal symmetries is
\bse
    \la \cO_1(x_1) ... \cO_4(x_4) \ra = f(u,v) \prod_{i<j} |x_{ij}|^{\frac{\Delta_1 + ... + \Delta_4}{3} - \Delta_i - \Delta_j}
\ese
where $u,v$ are the so-called \textit{conformal cross-ratios}
\be 
\label{eqn:CrossRatios}
    u := \frac{|x_{12}||x_{34}|}{|x_{13}||x_{24}|} \qand v := \frac{|x_{14}||x_{23}|}{|x_{13}||x_{24}|}.
\ee 
The big problem is that $f$ is an arbitrary function of $u,v$. The claim is that the product part
\bse
    \prod_{i<j} |x_{ij}|^{\frac{\Delta_1 + ... + \Delta_4}{3} - \Delta_i - \Delta_j}
\ese 
satisfies \Cref{eqn:CorrelationFunctionVariationOperators}, from which it follows from this that the cross-ratios are conformally invariant themselves. In particular they satisfy 
\be 
\label{eqn:CrossRatiosSpecialConformal}
    (K_{1\mu}+K_{2\mu}+K_{3\mu}+K_{4\mu})u = 0 = (K_{1\mu}+K_{2\mu}+K_{3\mu}+K_{4\mu})v.
\ee 
\bbox 
    Prove that the conformal cross-ratios are invariant under special conformal transformations. That is prove \Cref{eqn:CrossRatiosSpecialConformal}.
\ebox 

\br 
    Note that even though we haven't fully solved this problem we have taken it from being a function of $4D$ variables (the $4$ $x^{\mu}$s) and reduced it to the dependence on $2$ variables. 
\er 

\subsection{Higher Point Functions}

When we consider higher point functions we get more and more invariant cross-ratios, and so \textit{any} function solves the CFT. This is an obvious problem and we have to resort to other methods to fix these. In order to do this we need to derive the (several times mentioned) operator-state correspondence and so-called \textit{crossing symmetry}. 
\chapter{Radial Quantisation \& The State-Operator Correspondence}

\section{Radial Quantisation}

Recall from a canonical QFT course that we tend to quantise the fields on equal time slices, i.e. we take our (anti)commutators to be equal time (anti)commutators. However this is not the only slicing we can take. This might sound strange at first but recall that how we choose to slice the spacetime is essentially arbitrary so there is no reason to assume that equal time is the only way.\footnote{In fact if it was then this would be non-Lorentz invariant, as we are giving time `more priority' to space.} It turns out for CFTs it is useful to use "equal radial slices" in order to do our quantisation, and this process is known, unsurprisingly, as \textit{radial quantisation}. 

We can motivate this by first considering the theory on a cylinder. That is we consider a Euclidean\footnote{The $\R$ here is the Euclidean time.} CFT on $\R\times S^{D-1}$ (the $D$-dimensional cylinder) instead of on the more common $\R^{1,D-1}$.  We coordinatise our cylinder by using $\tau$ for the Euclidean time, $\R$, and $\underline{n}$ for the sphere, $S^{D-1}$. We also choose the coordinates such that we have a unit sphere, i.e. $\underline{n}\cdot \underline{n} =1$. We illustrate this below for $D=2$.\footnote{As it's hard to draw $4D$ and higher objects...}

\begin{center}
    \btik 
        \draw[thick] (0,3) ellipse (0.75 and 0.2);
        \draw[thick] (-0.75,0) arc (180:360:0.75 and 0.2);
        \draw[thick, dashed] (-0.75,0) arc (180:360:0.75 and -0.2);
        \draw[thick] (-0.75,0) -- (-0.75,3);
        \draw[thick] (0.75,0) -- (0.75,3);
        \draw[thick, dashed] (0,1.5) ellipse (0.75 and 0.2);
        \draw[thick, ->] (-1,1) -- (-1,2) node [midway, left] {$\tau$};
        \draw[thick, ->] (-0.5,1.1) arc (200:340: 0.5 and 0.133) node [midway, below] {$\underline{n}$};
    \etik 
\end{center}

\noindent The metric\footnote{Technically speaking this is a line element, however we can extract the metric components from it. We wont be so picky to distinguish between the two in this course (apart from this footnote...)} on our cylinder is just
\bse 
    ds^2 = d\tau^2 + d\underline{n}^2
\ese 
where we have a plus sign for both as we are considering the Euclidean cylinder. For example, if we had $D=3$ then we would have the 2-sphere $S^2$, and our metric would be
\bse 
    ds^2 = d\tau^2 + d\theta^2 + \sin^2\theta d\varphi^2.
\ese 

Now let's consider a coordinate transformation
\bse 
    \tau \to r = e^{\tau} \qquad \implies \qquad \frac{dr}{d\tau} = e^{\tau} = r,
\ese 
so our metric becomes 
\bse 
    ds^2 = \frac{dr^2}{r^2} + d\underline{n}^2 = \frac{1}{r^2}\big( dr^2 + r^2d\underline{n}^2)
\ese
where we notice that the term in brackets is flat space, or  Euclidean,\footnote{Not to be confused with Euclidean as in Wick rotate.} metric. 

So all our coordinate transformation has done is scale the metric by an overall factor. This is just a Weyl transformation, and so we have shown that there exists a conformal transformation from
\bse 
    \R\times S^{D-1} \to \R^D
\ese 
with 
\be 
\label{eqn:ConformalFactorRadial}
    \Omega = r = |x|.
\ee 
Let's look at the properties of this map.
\ben[label=(\roman*)] 
    \item The infinite past, $\tau=-\infty$, is mapped to the origin $r=0$. 
    \item Increasing time corresponds to increasing radius. More specifically, 
    \bse 
        \frac{\p}{\p\tau} \to \frac{\p}{\p \ln r} = r\frac{\p}{\p r}.
    \ese 
    This tells us that equal time quantisation corresponds to equal radius quantisation. We also see that time ordering corresponds to radial ordering. 
    \item From the above point we see that the energy on the cylinder corresponds to the dilatation weight on the plane. This is because the energy on the cylinder is given by the Hamiltonian which is given by $\p_{\tau}$, whereas a dilatation corresponds to scaling the system. We have just seen, though, that propagating in time on the cylinder corresponds to scaling on the plane. This further supports our idea that $\Delta$ weight is given by the mass dimension (as $[E]=[m]$). 
\een 
We summarise the idea in the following diagram.
\begin{center}
    \btik
        \draw[thick] (0,0) -- (0,5);
        \draw[thick] (2,0) -- (2,5);
        \draw[thick] (1,5) ellipse (1 and 0.18);
        \draw[thick] (0,0) arc (180:360: 1 and 0.18);
        \draw[thick, dashed] (0,0) arc (180:360: 1 and -0.18); 
        \draw[thick, dashed, blue] (1,1.5) ellipse (1 and 0.18);
        \draw[thick, dashed, red] (1,3) ellipse (1 and 0.18);
        \node at (-0.5, 1.5) {$\tau_1$};
        \node at (-0.5, 3) {$\tau_2$};
        %
        \draw[ultra thick, ->] (2.5,2.5) -- (5.5,2.5) node [midway, above] {$\tau \to r=e^{\tau}$};
        %
        \draw[thick] (6,0) -- (11,0) -- (11,5) -- (6,5) -- (6,0);
        \draw (8.5,0) -- (8.5,5);
        \draw (6,2.5) -- (11,2.5);
        \draw[thick, dashed, blue] (8.5,2.5) circle (1cm);
        \draw[thick, dashed, red] (8.5,2.5) circle (2cm);
        \node at (9.35,3.35) {$r_1$};
        \node at (10.1,4.1) {$r_2$};
    \etik 
\end{center}

\badr 
    Note that, just as spacelike separated operators commute in a causal Lorentz invariant QFT, we know that operators at the same \textit{radius} on our flat space will commute. These two statements are indeed exactly the same, as per the explanations above.
\eadr 

\badr 
    There is another way to understand the mapping from the cylinder to the plane pictorially, using what is known as \textit{stereographic projection}. Imagine putting the cylinder on top of the plane, and then drawing a line from the top of the cylinder\footnote{Which is infinitely far away, but you get the idea.} to the plane. The point on the cylinder that the line passes through corresponds to point of intersection on the plane. We can see from this that the infinite past (i.e. bottom of the cylinder) corresponds to the origin of the plane and increasing time corresponds to increasing radius on the plane.\footnote{To save myself messing with Tikz to get the angles correct, I request you google stereographic projection for a illustration.}
\eadr 

\subsection{Mapping Of Operators}

Ok great, so we have seen how to map into the radial picture, the next thing we need to ask is "how are the local operators mapped?" Well we have already said that our mapping is just a conformal transformation with the conformal factor given by $\Omega = r$, \Cref{eqn:ConformalFactorRadial}. So we just plug this into the transformation behaviour of our fields (and therefore operators), \Cref{eqn:PrimaryFieldTransformation}. For a scalar primary operator we therefore just have 
\bse 
    \cO_{\text{cyl}}(\tau,\underline{n}) \to \cO_{\text{flat}}(r,\underline{n}) = \frac{1}{r^{\Delta}} \cO_{\text{cyl}}(\tau,\underline{n}).
\ese 

Now we also know from QM that the Hermitian conjugate plays an important role, so let's see how that acts. On the cylinder we have 
\bse 
    \cO_{\text{cyl}}(\tau,\underline{n})^{\dagger} = \cO_{\text{cyl}}(-\tau,\underline{n})
\ese
where the minus sign comes from our Wick rotation, $\tau=-it$. This concept is called \textit{reflection positivity}, it is the Euclidean equivalent to unitarity in the Lorentzian picture.\footnote{In general reflection positivity tells us whether our Euclidean picture is just a Wick-rotated version of a unitary Lorentzian theory or not. That is, given a arbitrary Euclidean theory, there is no reason a priori to assume the Lorentzian theory you get by `counter' Wick rotating will be unitary. If the Euclidean picture is reflection positive, then we know the Lorentzian theory is unitary. This is apparently the content of the \textit{Osterwalder-Schrader reconstruction theorem}. Note that of course this worked out for us because we got our Euclidean picture by Wick rotating a unitary Lorentzian theory.} We can then use our radial mapping to see what happens to the local operators on our flat space:
\bse 
    \cO_{\text{flat}} (r,\underline{n})^{\dagger} = \bigg( \frac{1}{r^{\Delta}} \cO_{\text{cyl}}(\tau,\underline{n})\bigg)^{\dagger} = \frac{1}{r^{\Delta}} \cO_{\text{cyl}}(-\tau,\underline{n}) = \frac{1}{r^{\Delta}} \bigg(\frac{1}{r}\bigg)^{\Delta} \cO_{\text{flat}}\bigg(\frac{1}{r}, \underline{n}\bigg)
\ese
Therefore Hermitian conjugation we have
\be 
\label{eqn:HermitianConjugateOperatorFlat}
    \cO_{\text{flat}}(r,\underline{n})^{\dagger} = \frac{1}{r^{2\Delta}} \cO_{\text{flat}}\bigg(\frac{1}{r}, \underline{n}\bigg).
\ee

We now note that this is just the action of an inversion on the operator, i.e. if we act with 
\bse 
    I : x^{\mu} \mapsto \frac{x^{\mu}}{x^2}
\ese 
on the operator we will get \Cref{eqn:HermitianConjugateOperatorFlat}. 

\bcl 
    We can use this inversion result to show that 
    \be
    \label{eqn:HermitianConjugateCharge}
        Q_{\xi_i}^{\dagger} = - Q_{I \xi_i I}.
    \ee 
\ecl 

\bq 
    \textcolor{red}{I am not actually sure how to do this. Simmons-Duffin says to consider acting on the stress-tensor. I tried doing that but couldn't quite get the correct result. This is a note to self to remember to try do this later and come back and include it here.} % Note to Paul: If you know how to do this could I come to your office and see how to do it?
\eq 

For example, recalling \Cref{eqn:KEqualIPI}, we can use \Cref{eqn:HermitianConjugateCharge} to deduce that 
\be
\label{eqn:KDaggerP}
    (\widetilde{K}^{\mu})^{\dagger} = \widetilde{P}^{\mu},
\ee 
in radial quantisation. 

\badr 
    Note that this seems very strange at first. We are used to the momentum charge being Hermitian in a QFT, however we have essentially just shown that its Hermitian conjugation is the special conformal charge. The reason for this mismatch is that the Hilbert space in our radial quantisation picture is different from the Hilbert space of our Lorentzian system. This is actually not a surprising result as we have quantised differently and so we expect the resulting Hilbert space to change. Therefore we should really distinguish between $\dagger_{\text{Cyl}}$ and $\dagger_{\text{flat}}$. We won't make such distinguishes here, but this remark is just included to clear up potential confusion. 
\eadr 

\section{The State-Operator Correspondence}

We are now ready to present the highly non-trivial \textit{state-operator correspondence}, which is also called the state-operator map. 

\mybox{
    \bt[State-Operator Correspondence]
        In a conformal field theory, there is a one-to-one correspondance, i.e. an isomorphism, between states and \textbf{local} operators.
    \et 
}

Note that we have been very careful to say \textit{local} everywhere in these notes. This is because the state-operator map tells us that this correspondence only holds for \textit{local} operators. In particular this tells us that the total number of general operators and states do not agree (indeed this is \textit{never} the case), but simply that the number of local operators and states agree. Even with this reduction the state-operator map is still highly-non trivial.

\badr
    Before we go on to present a proof of this correspondence, let's first see how non-trivial this result really is. The states of our system live in some Hilbert space and so we can represent them as $n$-column matrices. On the other hand, operators act on the Hilbert space and map us from one element to another. We can therefore represent them as $(n\times n)$ matrices. The state-operator map is telling us that there is an isomorphism between these two structures. It's somewhat reasonable to believe we can get a state from a given local operator, however it is really non-trivial that we can recover a local operator from a state.
\eadr

\subsection{Motivating A Proof}

As the name of this section suggests, we do not present a strict proof of the state-operator here but simply motivate why it's actually a reasonable result. A more detailed derivation can be found in Section 4.6 of David Tong's notes or in Polchinski.

\subsubsection{State $\Rightarrow$ Local Operator}

We shall first discuss the more non-trivial step of obtaining a local operator from the state.

First consider states in QM, i.e. 1-particle in 1 spatial dimension. Here states are equivalent to the Schr\"{o}dinger wavefunction $\psi(x,t_0)$ at a particular time. Note that these are defined across an entire spatial slice, we shall use this idea soon. Now in particular let's focus on an initial state, i.e. $t_0 \to -\infty$, this is given by
\bse 
    \psi_I(x) := \lim_{t_0\to-\infty} \psi(x,t_0).
\ese 

In QFTs we don't normally consider wavefunctions, but there is no reason that we can't. Here we don't have wavefunctions but wavefunction\textit{al}s, however this is a minor detail and won't effect much of the intuition. For simplicity, consider just a single scalar particle (the generalisation should be clear), then the $x$ in QM becomes $\varphi(x,t_0)$ in QFT. That is the configuration space becomes the space of fields at a particular time. An initial state $\psi_I(x)$ becomes a limit of a functional
\bse 
    \lim_{t_0\to-\infty} \Psi[\varphi(x,t_0), t_0].
\ese 
What does this mean? Well, let's go to radial quantisation where $t_0\to-\infty$ becomes $r\to 0$, giving us
\bse 
    \lim_{r\to 0} \Psi[\varphi(r,\underline{n}),r],
\ese 
so it is essentially a function of $\varphi(\underline{x})$ at $\underline{x}=0$. That is $\Psi$ is a function of $\varphi(0)$, $\p_{\mu} \varphi(0)$, $\p_{\mu}\p_{\nu}\varphi(0)$ etc, but a function of $\varphi(0)$ etc, is precisely a local operator at $\underline{x}=0$. That is, given some general state we can obtain a local operator by taking an infinite dilatation `backwards', i.e. shrink it down to the origin. 

This is perhaps still confusing, so let's think of it more pictorially.\footnote{This explanation is based off the one given in Shiraz Minwalla's string theory course. See section 4.5 of my notes for more details.} As we said above, a state in a QFT corresponds to specifying it's values on an entire spatial slice. In the cylinder picture this corresponds to `cutting' the cylinder horizontally. When we take our radial map, this tells us that states on the flat space are given by equal radial slices. A general state, then, is just a circle around the origin, with the radius telling us the `time'. This is clearly not local; it's a whole circle! If we have any hope of producing a local operator we need to reduce the radius of this circle to a single point. We know how to do this, though: simply take an infinite dilatation to the origin. This is \textit{exactly} equivalent to considering the initial state that propagated to our general state. 

\begin{center}
    \btik 
        \begin{scope}[xshift=-5cm]
            \draw[thick] (-1.5,-1.5) -- (1.5,-1.5) -- (1.5,1.5) -- (-1.5,1.5) -- (-1.5,-1.5);
            \draw (0,-1.5) -- (0,1.5);
            \draw (-1.5,0) -- (1.5,0);
            \draw[dashed, thick] (0,0) circle [radius=1cm];
            \node at (0,-2) {`Time' $r$};
        \end{scope}
        \begin{scope}
            \draw[thick] (-1.5,-1.5) -- (1.5,-1.5) -- (1.5,1.5) -- (-1.5,1.5) -- (-1.5,-1.5);
            \draw (0,-1.5) -- (0,1.5);
            \draw (-1.5,0) -- (1.5,0);
            \draw[dashed, thick] (0,0) circle [radius=0.5cm];
            \node at (0,-2) {`Time' $r^{\prime}<r$};
        \end{scope}
        \begin{scope}[xshift=5cm]
            \draw[thick] (-1.5,-1.5) -- (1.5,-1.5) -- (1.5,1.5) -- (-1.5,1.5) -- (-1.5,-1.5);
            \draw (0,-1.5) -- (0,1.5);
            \draw (-1.5,0) -- (1.5,0);
            \draw[fill=black] (0,0) circle [radius=0.05cm];
            \node at (0,-2) {`Time' $r=0$};
        \end{scope}
    \etik 
\end{center}

\badr 
    Note that we have assumed that there are no operators inserted between the origin and our chosen general state. This is why would could just take an infinite dilatation to obtain the result at the origin. Of course if there was an operator `in the way' we would have to consider how the operator effects our state. This is not a problem though, as we can always just choose our initial state to lie inside all operator insertions. The only problem being if the operator was inserted at the origin itself. This still isn't a problem as we have translation invariance, so we can just shift the origin to a point where there are no operator insertions and go from there. In fact, as we will see shortly, the idea of having operators inserted within our circle gives us a very powerful result, known as the operator product expansion.
\eadr 

We can therefore obtain a local operator at the origin for a given state. We can therefore label our states by their corresponding operator at the origin,
\bse 
    \ket{\cO} := \cO(0) \ket{0}.
\ese 
We say that the state $\ket{\cO}$ is \textit{dual} to the operator $\cO$.

\bd[Vacuum]
    We can \textit{define} the vacuum to be the state dual to the identity,\footnote{This is actually only true in unitary theories. This is a subtle point but does actually appear, perhaps most notable in string theory: in order to obtain quantum Weyl invariance in string theory we have to include $b,c$ ghost system. These form a non-unitary theory and here the state $\ket{\b1}$ is not the vacuum but is actually an excited state. See Section 11.5.4 of my string theory notes for more details.} 
    \be 
    \label{eqn:VacuumDualToIdentity}
        \ket{0} := \ket{\b1}.
    \ee
\ed 

\br 
    We can now see why our conformal charges annihilate the vacuum, as said in footnote 2 of the last section: the action on the vacuum is given by the commutator with the identity operator, but everything commutes with the identity operator, and so the result vanishes. 
\er 

\subsubsection{State $\Leftarrow$ Local Operator}

So we have seen how to relate an operator at the origin to a state, however not all our operators will be inserted at the origin, so how do we get the states for these operators? That is\footnote{For clarity, $\cO(x)$ does not mean a \textit{function} of $x$, as this would be non-local. Instead it just means $\cO$ inserted at the fixed event $x$.}
\bse 
    \ket{\psi} = \cO(x) \ket{0} = ?
\ese

The answer is we Taylor expand our operator around the origin,
\bse 
    \cO(x) = \cO(0) + x^{\mu}\p_{\mu} \cO(0) + \frac{1}{2}x^{\mu}x^{\nu}\p_{\mu}\p_{\nu}\cO(0) + ...,
\ese 
so that we conclude 
\bse 
    \ket{\psi} = \sum_n \frac{1}{n!} x^{\mu_1} ... x^{\mu_n} \ket{\p_{\mu_1} ... \p_{\mu_n} \cO}.
\ese

The question then obviously becomes "what is $\ket{\p_{\mu}\cO}$?" This is actually quite straight forward: use our state-operator correspondence at the origin to obtain 
\bse 
    \ket{\p_{\mu}\cO} := \p_{\mu}\cO(0)\ket{0} = [\widetilde{P}_{\mu},\cO(0)]\ket{0} = \widetilde{P}_{\mu}\cO(0)\ket{0} = \widetilde{P}_{\mu}\ket{\cO},
\ese 
where we have used $\widetilde{P}_{\mu}\ket{0} = 0$. We can therefore obtain a state from the insertion of a local operator at any $x$.

We can think of the local operator to state map pictorially as well.\footnote{Again see my notes on Shiraz Minwalla's string theory course for more details.} The idea is that the area inside our circles represent the evolution of our states. We start with our initial wavefunctional at the origin and allow it to evolve freely. At some point, $x$, it comes into contact with an operator, and this disturbs the free evolution, resulting in some new state, which then again evolves freely (until it meets another operator). We can then associate the new state with the operator itself, giving us our duality.

\badr 
    Note, as Prof. Tong points out,\footnote{See page 100-101 of his String Theory notes.} this state from local operator is \textit{not} the same idea as using creation/annihilation operators in QFT. They look deceptively similar algebraically, however creation/annihilation operators are given by taking Fourier transforms of local fields, and so are completely \textit{un}localised!
\eadr 

\subsection{Raising \& Lowering Operators}

Recall that normally in a QFT we label the states by their quantum numbers, e.g. their momentum and spin. We now want to do a similar thing for our conformal states. As we said before, in the radial picture the energy corresponds to the dilatation weight, and so we label our states by their $\Delta$. We can check this is a good thing to do by computing the action of the dilatation charge explicitly. Consider a general state (i.e. primary \textit{or} descendant) $\ket{\cO_{\Delta}}$, where $\cO_{\Delta}$ is an operator with weight $\Delta$. Then
\bse 
    \begin{split}
        \widetilde{D}\ket{\cO_{\Delta}} & = \widetilde{D}\cO_{\Delta}(0)\ket{0} \\
        & = \big([\widetilde{D},\cO(0)] + \cO(0) \widetilde{D}\big)\ket{0} \\
        & = \Delta \cO_{\Delta}(0)\ket{0} \\
        \implies D\ket{\cO_{\Delta}} & = \Delta \ket{\cO_{\Delta}},
    \end{split}
\ese 
where we used $\widetilde{D}\ket{0}=0$.\footnote{This makes sense as the vacuum is dual to the identity operator, and the identity doesn't have a dilatation weight.} We can therefore label states by their dilatation weight 
\bse 
    \ket{\cO_{\Delta}} \equiv \ket{\Delta}. 
\ese 

\bbox 
    Show that 
    \bse 
        \widetilde{D}\big( \widetilde{P}_{\mu} \ket{\cO} \big) = (\Delta+1)\widetilde{P}_{\mu} \ket{\cO},
    \ese 
    as well as 
    \bse 
        \widetilde{D}\big( \widetilde{K}_{\mu} \ket{\cO} \big) = (\Delta-1)\widetilde{K}_{\mu} \ket{\cO}.
    \ese
    These results tell us that $\widetilde{P}_{\mu}/\widetilde{K}_{\mu}$ are like raising/lowering operators for dilatation weight. \textit{Hint: Rewrite in terms of commutators, and recall we are at the origin, $x=0$.}
\ebox 

\badr 
     We have essentially already shown the result of the above exericse in \Cref{sec:DescendantFields}. The excerise is included for non-asterisk self contained-ness.
\eadr 

So we have seen that $\widetilde{P}_{\mu}$ and $\widetilde{K}_{\mu}$ act on the states of our Hilbert space as
\bse 
    \ket{\Delta} \overset{\widetilde{P}_{\mu}}{\longrightarrow} \ket{\Delta+1}  \overset{\widetilde{P}_{\mu}}{\longrightarrow} \ket{\Delta+2} ... 
\ese 
and
\bse 
    \ket{\Delta} \overset{\widetilde{K}_{\mu}}{\longrightarrow} \ket{\Delta-1}  \overset{\widetilde{K}_{\mu}}{\longrightarrow} \ket{\Delta-2} ...,
\ese 
which is exactly how raising/lowering operators act. Note this is also consistent with \Cref{eqn:KDaggerP}, i.e. the dagger of the lowering operator is the raising operator. Note from here we see that we can define primary \textit{states} as follows. 

\bd[Primary State]
    A state $\ket{\Delta} \equiv \ket{\cO_{\Delta}}$ is a \textit{primary state} if it is annihilated by $\widetilde{K}_{\mu}$.
\ed 

\noindent We then define descendent \textit{states} similarly.

\section{2-Point Function Again}

Recall that our two point functions of primary operators were given by 
\bse
    \bra{0}\cO_1(x) \cO_2(y) \ket{0} = \begin{cases}
        \frac{C_{12}}{|x_1-x_2|^{2\Delta}} & \text{if } \Delta_1=\Delta_2=\Delta \\
        0 & \text{otherwise},
    \end{cases}
\ese 
where we have written the bra-kets explicitly for comparison of what follows. We can use this to show that $C_{12}$ is just given by the overlap of our operators inserted at the origin, i.e. the overlap of the dual initial wavefunctionals. The first thing we note is that the overlap vanishes unless $\Delta_1=\Delta_2$, as states with different dilatation weight are given by raising/lowering operators. Now it turns out we can diagonalise our Hilbert space such that the inner product of two operators (i.e. the 2-point function) is orthogonal. That is if $\cO_1\neq \cO_2$ are two operators of weight $\Delta$,\footnote{Note this is \textit{not} the same thing as saying we only have one operator of weight $\Delta$, but just that the 2-point function vanishes.}
\bse 
    \la \cO_1(x) \cO_2(y) \ra = 0.
\ese 
It follows from this that the "otherwise" condition is already met.

The initial overlap is given by 
\bse 
    \braket{\cO_{\Delta_1}}{\cO_{\Delta_2}} = \lim_{r\to0}\lim_{u\to \infty} u^{2\Delta_1} \bra{0} \cO_{\Delta_1}(u,\underline{n}_1) \cO_{\Delta_2}(r,\underline{n}_2)\ket{0}
\ese 
where we have used the definition of the Hermitian and the substitution $u=1/r_1$ in the first term. We then plug in the two-point function result above to get  
\bse 
    \braket{\cO_{\Delta_1}}{\cO_{\Delta_2}} = \lim_{r\to0}\lim_{u\to \infty} u^{2\Delta_1} \begin{cases}
        \frac{C_{12}}{|x_1-x_2|^{2\Delta}} & \text{if } \Delta_1=\Delta_2=\Delta \\
        0 & \text{otherwise}.
    \end{cases}
\ese 
Here we have 
\bse 
    (x_1-x_2)^{\mu} = u\underline{n}_1^{\mu} - r\underline{n}_2^{\mu} \overset{r\to0}{\longrightarrow} u \underline{n}_1^{\mu} 
\ese 
and therefore 
\bse 
    (x_1-x_2)^2 \overset{r\to 0}{\longrightarrow} u^2
\ese
so 
\bse 
    \braket{\cO_{\Delta}}{\cO_{\Delta}} = \lim_{u\to\infty} u^{2\Delta} \frac{C_{12}}{u^{2\Delta}} = C_{12}. 
\ese 

We can use this result to cross check our result with what we would have got for the overlap purely algebraically. This calculation is left as the following exercise.

\bbox 
    Given the standard evolution relation
    \bse 
        \cO(x) = e^{x\cdot\widetilde{P}} \cO(0) e^{-x\cdot \widetilde{P}}
    \ese 
    show that 
    \bse 
        \braket{\psi_1}{\psi_2} = u^{2\Delta} \bra{\cO} \exp\big( u^2 x_1 \cdot \widetilde{K}\big) \exp\big(x_2\cdot \widetilde{P}\big) \ket{\cO},
    \ese 
    where $\ket{\psi_i} := \cO(x_i)\ket{0}$, and $u := 1/x_1$.
    \textit{Hint: Recall \Cref{eqn:KDaggerP} and the fact that the conformal charges annihilate the vacuum.}
\ebox 

We can use the result of this exercise to check the two agree. We do this by Taylor expanding the exponentials and using the fact that our states are primary, i.e. $\widetilde{K}^{\mu}\ket{\cO} = 0 = \bra{\cO}\widetilde{P}^{\mu}$ with the second following from taking a Hermitian conjugation. The first few terms in the expansion are then 
\bse 
    \begin{split}
        \braket{\psi_1}{\psi_2} & \approx u^{2\Delta}\Big( \braket{\cO}{\cO} + u^2 x_1^{\mu} x_2^{\nu} \bra{\cO} \widetilde{K}_{\mu}\widetilde{P}_{\nu} \ket{\cO}\Big) \\
        & = u^{2\Delta}\Big( C_{12} + u^2 x_1^{\mu} x_2^{\nu} \bra{\cO} [\widetilde{K}_{\mu},\widetilde{P}_{\nu}]\ket{\cO} \Big) \\
        & = u^{2\Delta}\Big( C_{12} + u^2 x_1^{\mu} x_2^{\nu} \bra{\cO} \big( 2\eta_{\mu\nu} \widetilde{D} -2\widetilde{L}_{\mu\nu} \big)\ket{\cO} \Big) \\
        & = u^{2\Delta}\big( C_{12} + 2\Delta u^2 x_1\cdot x_2  \bra{\cO}\ket{\cO} \big) \\
        & = x_1^{-2\Delta} C_{12} \bigg( 1 + 2\Delta\frac{x_1\cdot x_2}{x_1^2} \bigg),
    \end{split}
\ese 
where we have used the fact that we are considering a scalar primary to drop the $\widetilde{L}_{\mu\nu}$.\footnote{This is because $\widetilde{L}_{\mu\nu}\ket{\cO} = S_{\mu\nu}\ket{0}$.} This result agrees exactly with the first few terms in the expansion of 
\bse 
    \braket{\cO(x_1)}{\cO(x_2)} = \frac{C_{12}}{|x_1-x_2|^{2\Delta}} = x_1^{-2\Delta} \frac{C_{12}}{|1-\frac{x_2}{x_1}|^{2\Delta}}.
\ese 

\bbox 
    Prove that the 3-point function coefficient $C_{123}$ of three scalar operators of dimensions $\Delta_1$, $\Delta_2$ and $\Delta_3$ is equal to sandwiching $\cO_{\Delta_2}(\hat{1})$ in between $\bra{\cO_{\Delta_1}}$ and $\ket{\cO_{\Delta_3}}$ where $\hat{1} = (1,0,0,...)$. That is show 
    \bse 
        \bra{\cO_{\Delta_1}} \cO_{\Delta_2}(\hat{1}) \ket{\cO_{\Delta_3}}.
    \ese 
    \textit{Hint: Start the same as above by taking the limits on the arguments of $\cO_{\Delta_1}$ and $\cO_{\Delta_2}$ and use the known 3-point function result.}
\ebox 

\section{Operator Product Expansion (OPE)}

There is a very powerful consequence of the state-operator correspondance, it is known as the \textit{operator product expansion} (OPE). 

\bl 
    In a CFT, we can write the product of two local operators at different points as a sum of local operators (both primaries and descendants) at one point. In particular, labelling our given local operators with $\Phi$s,
    \be 
    \label{eqn:OPE}
        \Phi_{\Delta_1}(x) \Phi_{\Delta_2}(0) = \sum_{\text{primary  ops.}} C_{\Phi_1\Phi_2\cO} C_{\cO}(x,\p_y) \cO(y) \big|_{y=0}
    \ee 
    where $C_{\cO}(x,\p_y)$ is a power series in $\p_y$ which generates our descendants, and $C_{\Phi_1\Phi_2\cO}$ is a theory dependent prefactor, known as \textit{OPE coefficients}.
\el

The important point to note is that the series in \Cref{eqn:OPE} is convergent (provided there is no other operators between $0$ and $x$).

\br 
    The power series, $C_{\cO}(x,\p_y)\cO(y)$, are a purely kinematical thing and only depends on the representation of the theory. The OPE coefficients, $C_{\Phi_1\Phi_2\cO}$, however are theory dependent and so allow us to distinguish the different CFTs. 
\er 

We can understand the idea of an OPE using our pictorial approach to the state-operator correspondence. Let's imagine the two local operators within one of our state circles. We can view this state as being produced by some initial state being altered by the two operators in turn. However we could `forget' about these operators and instead assume that our state was produced by a single operator at the origin, by the state $\Rightarrow$ operator map. Now of course its highly unlikely that a single primary operator at the origin will give us the state we want, however we know that we can always express a general state in our Hilbert space as a sum over a basis, which is what the right-hand side of \Cref{eqn:OPE} corresponds to. 

\begin{center}
    \btik 
        \begin{scope}[xshift=-5.5cm]
            \draw[thick] (-2,-2) -- (2,-2) -- (2,2) -- (-2,2) -- (-2,-2);
            \draw (-2,0) -- (2,0);
            \draw (0,-2) -- (0,2);
            \node at (-0.5,0.6) {$\cross$};
            \node[above] at (-0.5,0.6) {$\Phi_1$};
            \node at (0.2,-1) {$\cross$};
            \node[above] at (0.2,-1) {$\Phi_2$};
            \draw[ultra thick, ->] (2.25,0) -- (3.25,0); 
        \end{scope}
        \begin{scope}
            \draw[thick] (-2,-2) -- (2,-2) -- (2,2) -- (-2,2) -- (-2,-2);
            \draw (-2,0) -- (2,0);
            \draw (0,-2) -- (0,2);
            \draw[thick, dashed] (0,0) circle [radius=1.5cm];
            \node at (1.3,1.3) {$\ket{\psi}$};
            \draw[ultra thick, ->] (2.25,0) -- (3.25,0);
        \end{scope}
        \begin{scope}[xshift=5.5cm]
            \draw[thick] (-2,-2) -- (2,-2) -- (2,2) -- (-2,2) -- (-2,-2);
            \draw (-2,0) -- (2,0);
            \draw (0,-2) -- (0,2);
            \node at (0,0) {$\cross$};
            \node at (0.4,0.3) {$\cO_{\ket{\psi}}$};
        \end{scope}
    \etik 
\end{center}

Note that in the above picture we have taken the two $\Phi$s to be away from the origin whereas \Cref{eqn:OPE} has $\Phi_2$ at the origin. This is obviously not a problem as we can go between the two by a translation.

\badr 
    The idea of an OPE is actually useful in other, non-CFT, QFTs. However there it is only an approximate statement and requires the separation of the two operators to be smaller than the distance between these operators and other operators. That is, if $\Phi_1$ is at $x_1$ and $\Phi_2$ is at $x_2$, the other operator insertions in our correlator have to be more than $|x_1-x_2|$ away from $\Phi_1$/$\Phi_2$. For CFTs, though, our OPE is \textit{exact} and we just require that the other operator insertions are outside our state circle. 
\eadr 

Why is the OPE so useful? Well recall that we said that $n$-point correlators for $n\geq 4$ aren't fully constrained by our conformal symmetries, but the $2$-point and $3$-point functions were. The OPE allows us to reduce any $n$-point function to an $(n-1)$-point function as 
\bse 
    \la \cO_1 \cO_2 .... \cO_{n-1}\cO_n\ra = \sum_{\cO'} C_{\cO_{n-1}\cO_n\cO'} C_{\cO'}(x,\p_y) \la \cO_1\cO_2...\cO_{n-2}\cO'\ra. 
\ese 
The idea is to reduce all higher point functions down to 3-point functions weighted by the OPE coefficients. 

So how can we compute $C_{\cO}(x,\p_y)$? Well firstly we simplify our lives by assuming that all the operators in our OPE are scalars. Of course this need not be the case as both the $\Phi$s and $\cO$s could be in non-trivial Lorentz representations, i.e. \Cref{eqn:OPE} could have indices everywhere. We will also work with one of the operators inserted at the origin, as with \Cref{eqn:OPE}. 

Ok, now consider the 3-point function
\bse 
    \la \Phi_1(x) \Phi_2(0) \Phi_3(z) \ra.
\ese 
Plugging in the OPE for $\Phi_1(x)\Phi_2(0)$, we get
\bse
    \sum C_{\Phi_1\Phi_2\cO} C_{\cO}(x,\p_y) \la \cO(y)|_{y=0} \Phi_3(x)\ra
\ese 
but we know what both these expressions are (up to a constant). That is we know the general solution to a 2-point function and 3-point function. 

Now recall that the two-point function is only non-zero when the two operators have the same dilatation weight. Also recall that we can pick a basis where any two operators with the same dilatation weight are orthogonal. We therefore conclude that the two-point function is only non-vanishing if the two operators are exactly the same. Finally, assume we normalised the operators so that 
\bse 
    \la \cO(y) \Phi_3(z) \ra = \begin{cases}
        \frac{1}{|y-z|^{2\Delta_3}} & \Phi_3 = \cO \\
        0 & \text{otherwise},
    \end{cases}
\ese 
that is we have set $C_{12}=1$. We shall now assume that $\Phi_x=\cO$ and assume the "otherwise" case is understood implicitly.  

Then recalling the 3-point function result, \Cref{eqn:ThreePointCorrelator}, we have
\be 
\label{eqn:Exercise}
    \frac{C_{\Phi_1\Phi_2\Phi_3}}{|x|^{\Delta_1+\Delta_2-\Delta_3} |z|^{\Delta_2+\Delta_3-\Delta_1} |x-z|^{\Delta_1+\Delta_3-\Delta_2} } = C_{\Phi_1\Phi_3\Phi_3} C_{\Phi_3}(x,\p_y) \frac{1}{|y-z|^{2\Delta_3}}\bigg|_{y=0}. 
\ee 
Now identifying $C_{\Phi_1\Phi_2\Phi_3}$ with the 3-point coefficient,\footnote{That is, the $C_{\Phi_1\Phi_2\Phi_3}$ on the right-hand side is the OPE coefficient, $C_{\Phi_1\Phi_2\cO}$. We have used our $\cO=\Phi_3$ condition to get the right-hand side above.} this factors out. Expanding both sides in $x$ allows us to determine $C_{\Phi_3}(x,\p_y)$ term by term. The left-hand side is
\bse 
    \frac{1}{|x|^{\Delta_1+\Delta_2-\Delta_3} |z|^{\Delta_2+\Delta_3-\Delta_1} |z|^{\Delta_1+\Delta_3-\Delta_2}} = \frac{1}{|x|^{\Delta_1+\Delta_2-\Delta_3} |z|^{2\Delta_3}}
\ese 
which, upon comparing with right-hand side, we conclude that the leading term (as $x\to 0$) is 
\bse 
    C_{\Delta_3}(x,\p_y) = \frac{1}{|x|^{\Delta_1+\Delta_2-\Delta_3}}. 
\ese 

\bbox 
    Use \Cref{eqn:Exercise} to show that, for $\Delta_1=\Delta_2=\Delta$,
    \bse 
        C_{\Phi}(x,\p_y) = \frac{1}{|x|^{2\Delta-\Delta_3}} \bigg[ 1 + \frac{1}{2}x^{\mu}\p_{\mu} + \a x^{\mu}x^{\nu}\p_{\mu}\p_{\nu} + \beta x^2 \p^2 + \cO(x^3)\bigg]
    \ese 
    where 
    \bse 
        \a = \frac{\Delta_3+2}{8(\Delta_3+1)} \qand \beta = - \frac{\Delta_3}{16(\Delta_3 - \frac{D}{2} + 1)(\Delta_3+1)},
    \ese 
    where $D$ is the spacetime dimension. 
\ebox 
\chapter{Conformal Blocks \& Bootstrap Programme}

\section{CFT Data (Summary)}

As we have tried to emphasise, any CFT data is characterised by the set local primary operators $\{\cO_{\Delta,I}\}$, where $I$ denotes the general Lorentz indices as before, and their correlation functions. We refer to the set $\{\cO_{\Delta,I}\}$ as the \textit{spectrum} of local primary operators.

Let's review what we know so far about the correlation functions.
\ben[label=(\roman*)] 
    \item Two-point functions: 
        \ben 
            \item Are fixed up to a normalisation, e.g. for scalars we had $\la \cO^{(1)}_{\Delta}(x)\cO^{(2)}_{\Delta}(y)\ra = \frac{C_{12}}{|x-y|^{2\Delta}}$. 
            \item We can normalise the operators, i.e. $\cO^{(i)}_{\Delta}(x) \to \cO^{(i)}_{\Delta}(x)/\sqrt{A}$, such that the 2-point functions are completely fixed.
            \item We can diagonalise our representations by the dilatation weight, i.e. pick a basis such that two different operators of the same weight are orthogonal. 
        \een 
    \item Three-point functions:
        \ben 
            \item Are fixed up to a constant $C_{123}$, e.g. for scalars we have 
            \bse
                \la \cO_{\Delta_1}(x_1) \cO_{\Delta_2}(x_2)\cO_{\Delta_3}(x_3)\ra = \frac{C_{123}}{|x_{12}|^{\Delta_1+\Delta_2-\Delta_3} |x_{23}|^{\Delta_2+\Delta_3-\Delta_1} |x_{13}|^{\Delta_1+\Delta_3-\Delta_1}}.
            \ese
            \item The constant $C_{123}$ is set by how we choose to normalise our states, in other words it is fixed once we renormalise our operators as (i)(b) above.
        \een 
    \item Higher point functions: 
        \ben 
            \item Can be reduced to three point functions using the OPE.
            \item The terms appearing in the OPE are completely fixed in terms of the three-point functions. 
        \een 
\een 

So we have seen that we can essentially completely categorise a CFT via the above steps. We therefore make the following definition. 

\bd[CFT Data]
    We call the set of the of the spectrum and the OPE coefficients the \textit{CFT data}. That is 
    \be 
    \label{eqn:CFTData}
        (\text{CFT Data}) := \{ \cO_{\Delta,I} , C_{\cO_1\cO_2\cO_3} \}.
    \ee 
\ed 

This name is fitting as given this data we can compute \textit{any} correlation function, and therefore we have solved the CFT. 

\section{Consistency Conditions On CFT Data}

The next question is what kind of CFT data gives a well defined CFT? That is can we just give \textit{any} random set and be OK? The answer is, of course, "no" and we have some consistency conditions.

The first thing we note is that we have completely glazed over unitarity so far. It turns out that unitarity imposes bounds on the allowed values of $\Delta$ for a given Lorentz rep, for example for a scalar field we have
\be 
\label{eqn:UnitarityBoundOnDelta}
    \Delta \geq \frac{D}{2}-1,
\ee 
apart from identity operator which has $\Delta=0$. The bounds for other operator types can be found in Section 7.3 of Simmons-Duffin.

\br 
    Note that a scalar field saturates this bound. This is easily seen from computing its mass dimension from the Lagrangian. 
\er 

Secondly (and more interestingly), let's consider the 4-point functions
\bse 
    \la \Phi_1(x_1)\Phi_2(x_2) \Phi_3(x_3)\Phi_4(x_4)\ra. 
\ese 
We can take an OPE between $\Phi_1(x_1)\Phi_2(x_2)$ and another OPE between $\Phi_3(x_3)\Phi_4(x_4)$ to obtain 
\bse 
    \begin{split}
        \la \Phi_1(x_1)\Phi_2(x_2) \Phi_3(x_3)\Phi_4(x_4)\ra & = \sum_{\cO,\cO^{\prime}} C_{12\cO}C_{34\cO^{\prime}} \big[ C_{\cO}(x_{12},\p_y)  C_{\cO^{\prime}}(x_{34},\p_z) \la \cO(y) \cO^{\prime}(z)\ra \big] \\
        & = \sum_{\cO} C_{12\cO} C_{34\cO} \big[ C_{\cO}(x_{12},\p_y)  C_{\cO}(x_{34},\p_z) \la \cO(y) \cO(z)\ra \big]
    \end{split}
\ese 
where the second line follow from our diagonalisation procedure, i.e. the two point function is only non-vanishing when $\cO=\cO^{\prime}$. Now we recall that the bit in square brackets is completely theory independent, as the $C_{\cO}(x,\p_y)$ are purely dynamical and only depend on the representations (i.e. weights) and similarly the two-point function only depends on $|y-z|$ and weights. We call this square bracketed terms a \textit{conformal partial wave} or \textit{conformal block}, and we shall denote it as 
\bse 
    G_{\cO}(x_1,x_2,x_3,x_4) :=  C_{\cO}(x_{12},\p_y)  C_{\cO}(x_{34},\p_z) \la \cO(y) \cO(z)\ra 
\ese
so that  
\bse 
    \la \Phi_1 ... \Phi_4 \ra = \sum_{\cO} C_{12\cO}C_{34\cO} G_{\cO}(x_1,..., x_4).
\ese 

This is good, however (as we will indeed shortly do) we could have done the OPE contractions as $\Phi_2\Phi_3$ and $\Phi_1\Phi_4$. This would change our definition of our conformal block to be\footnote{Bonus exercise, verify that this is true.} 
\bse 
    G_{\cO}(x_1,x_2,x_3,x_4) :=  C_{\cO}(x_{23},\p_y)  C_{\cO}(x_{14},\p_z) \la \cO(y) \cO(z)\ra. 
\ese 
We can therefore tweak the definition of our conformal blocks to be 
\mybox{
    \be 
    \label{eqn:ConformalBlock}
        G_{\cO}(x_1,x_2,x_3,x_4) := |x_{ij}|^{\Delta_i+\Delta_j} |x_{k\ell}|^{\Delta_k+\Delta_{\ell}} C_{\cO}(x_{ij},\p_y) C_{\cO}(x_{k\ell}, \p_z) \la \cO(y)\cO(z)\ra 
    \ee 
}
\noindent with our 4-point function then given by 
\be 
\label{eqn:FourPointFunctionConformalBlocks}
    \la \Phi_1... \Phi_4 \ra = P_{ijk\ell} \sum_{\cO} C_{ij\cO} C_{k\ell\cO} G_{\cO}(x_1,x_2,x_3x_4),
\ee 
with 
\be 
\label{eqn:Pijkl}
    P_{ijk\ell} := |x_{ij}|^{-(\Delta_i+\Delta_j)} |x_{k\ell}|^{-(\Delta_k+\Delta_{\ell})}
\ee 
This might seem like overkill, however this will come in handy shortly. 

We can depict the conformal block pictorially so that our 4-point function is given by
\begin{center}
    \btik 
        \node at (-3,0) {\Large{$\sum C_{12\cO}C_{34\cO} $}};
        \node at (-4.1,-0.5) {\small{$\cO$}};
        \draw[thick, rotate around={-45:(0,0)}] (0,0) -- (-1,0) node [left] {$1$};
        \draw[thick, rotate around={45:(0,0)}] (0,0) -- (-1,0) node [left] {$2$};
        \draw[thick, rotate around={45:(1,0)}] (1,0) -- (2,0) node [right] {$4$};
        \draw[thick, rotate around={-45:(1,0)}] (1,0) -- (2,0) node [right] {$3$};
        \draw[thick] (0,0) -- (1,0) node [midway, above] {$\cO$};
    \etik 
\end{center}

Now we did the above OPEs in $\Phi_1\Phi_2$ and $\Phi_3\Phi_4$, but we could just have easily done them for $\Phi_2\Phi_3$ and $\Phi_1\Phi_4$. If we do that we of course get a different conformal block corresponding to the following diagram
\begin{center}
    \btik 
        \node at (-3,0) {\Large{$\sum C_{23\cO}C_{14\cO} $}};
        \node at (-4.1,-0.5) {\small{$\cO$}};
        \draw[thick, rotate around={-30:(0,0.5)}] (0,0.5) -- (-1,0.5) node [left] {$1$};
        \draw[thick,rotate around={30:(0,-0.5)}] (0,-0.5) -- (-1,-0.5) node[left] {$2$};
        \draw[thick, rotate around={-30:(0,-0.5)}] (0,-0.5) -- (1,-0.5) node [right] {$3$};
        \draw[thick, rotate around={30:(0,0.5)}] (0,0.5) -- (1,0.5) node [right] {$4$};
        \draw[thick] (0,0.5) -- (0,-0.5) node [midway, right] {$\cO$};
    \etik 
\end{center}

Of course these two things have to agree (as they are both the 4-point function) and so we get a constraint between the allowed OPE coefficients $\{C_{ij\cO}\}$ where $i,j=1,2,3,4$. This relation is known as a \textit{crossing symmetry}, and it is a very powerful constraint on the allowed data by placing a constraint on the OPE associativity. 

\section{Conformal Bootstrap}

\bcl 
    Once we impose the $4$-point crossing symmetry, there are \textit{no new constraints} at higher point functions.
\ecl 

This is a non-trivial result and \textit{massively} simplifies the study of CFTs. Let's prove that it works for the case of a 5-point function. In order to lighten notation, we shall absorb the OPE coefficients into the diagrams.\footnote{They can be read off by looking at the vertices.} We shall use a $\bullet$ to indicate a point in the diagram that still needs to OPEd.

\subsubsection{Step 1: Do $\Phi_1\Phi_2$ OPE}
\begin{center}
    \btik 
        \begin{scope}
            \draw[thick, rotate around={-45:(0,0)}] (0,0) -- (-1,0) node [left] {$1$};
            \draw[thick, rotate around={45:(0,0)}] (0,0) -- (-1,0) node [left] {$2$};
            \draw[thick, rotate around={-60:(0,0)}] (0,0) -- (1,0) node [right] {$3$};
            \draw[thick] (0,0) -- (1,0) node [right] {$4$};
            \draw[thick, rotate around={60:(0,0)}] (0,0) -- (1,0) node [right] {$5$};
            \draw[fill=black] (0,0) circle [radius=0.1cm];
            \node[right] at (1.75,0) {\Large{$=$}};
        \end{scope}
        \begin{scope}[xshift=4.75cm]
            \node at (-1.75,0) {\Large{$\sum$}};
            \node at (-1.75,-0.5) {\small{$\cO$}};
            \draw[thick, rotate around={-45:(0,0)}] (0,0) -- (-1,0) node [left] {$1$};
            \draw[thick, rotate around={45:(0,0)}] (0,0) -- (-1,0) node [left] {$2$};
            \draw[thick] (0,0) -- (1,0) node [midway, above] {$\cO$};
            \draw[thick, rotate around={-60:(1,0)}] (1,0) -- (2,0) node [right] {$3$};
            \draw[thick] (1,0) -- (2,0) node [right] {$4$};
            \draw[thick, rotate around={60:(1,0)}] (1,0) -- (2,0) node [right] {$5$};
            \draw[fill=black] (1,0) circle [radius=0.1cm];
        \end{scope}
    \etik 
\end{center}

\subsubsection{Step 2: Do $\Phi_3\Phi_4$ OPE}

\begin{center}
    \btik 
        \begin{scope}
            \node at (-1.75,0) {\Large{$\sum$}};
            \node at (-1.75,-0.5) {\small{$\cO$}};
            \draw[thick, rotate around={-45:(0,0)}] (0,0) -- (-1,0) node [left] {$1$};
            \draw[thick, rotate around={45:(0,0)}] (0,0) -- (-1,0) node [left] {$2$};
            \draw[thick] (0,0) -- (1,0) node [midway, above] {$\cO$};
            \draw[thick, rotate around={-60:(1,0)}] (1,0) -- (2,0) node [right] {$3$};
            \draw[thick] (1,0) -- (2,0) node [right] {$4$};
            \draw[thick, rotate around={60:(1,0)}] (1,0) -- (2,0) node [right] {$5$};
            \draw[fill=black] (1,0) circle [radius=0.1cm];
            \node[right] at (2.75,0) {\Large{$=$}};
        \end{scope}
        \begin{scope}[xshift=5.75cm]
            \node at (-1.75,0) {\Large{$\sum$}};
            \node at (-1.75,-0.5) {\small{$\cO, \cO^{\prime}$}};
            \draw[thick, rotate around={-45:(0,0)}] (0,0) -- (-1,0) node [left] {$1$};
            \draw[thick, rotate around={45:(0,0)}] (0,0) -- (-1,0) node [left] {$2$};
            \draw[thick] (0,0) -- (1,0) node [midway, above] {$\cO$};
            \draw[thick, rotate around={60:(1,0)}] (1,0) -- (2,0) node [right] {$5$};
            \draw[thick] (1,0) -- (1.7,-0.7);
            \node at (1.6,-0.2) {$\cO^{\prime}$};
            \draw[thick, rotate around={25:(1.7,-0.7)}] (1.7,-0.7) -- (2.7,-0.7) node [right] {$4$};
            \draw[thick, rotate around={-65:(1.7,-0.7)}] (1.7,-0.7) -- (2.7,-0.7) node [right] {$3$};
        \end{scope}
    \etik 
\end{center}

\subsubsection{Step 3: Use $4$-point crossing symmetry on $\Phi_1\Phi_2\cO\Phi_5\cO^{\prime}$}

\begin{center}
    \btik 
        \begin{scope}
            \node at (-1.75,0) {\Large{$\sum$}};
            \node at (-1.75,-0.5) {\small{$\cO, \cO^{\prime}$}};
            \draw[thick, rotate around={-45:(0,0)}] (0,0) -- (-1,0) node [left] {$1$};
            \draw[thick, rotate around={45:(0,0)}] (0,0) -- (-1,0) node [left] {$2$};
            \draw[thick] (0,0) -- (1,0) node [midway, above] {$\cO$};
            \draw[thick, rotate around={60:(1,0)}] (1,0) -- (2,0) node [right] {$5$};
            \draw[thick] (1,0) -- (1.7,-0.7);
            \node at (1.6,-0.2) {$\cO^{\prime}$};
            \draw[thick, rotate around={25:(1.7,-0.7)}] (1.7,-0.7) -- (2.7,-0.7) node [right] {$4$};
            \draw[thick, rotate around={-65:(1.7,-0.7)}] (1.7,-0.7) -- (2.7,-0.7) node [right] {$3$};
            \node[right] at (3.5,0) {\Large{$=$}};
        \end{scope}
        \begin{scope}[xshift=6.75cm]
            \node at (-1.75,0) {\Large{$\sum$}};
            \node at (-1.75,-0.5) {\small{$\cO^{\prime}, \cO^{\prime\prime}$}};
            \draw[thick, rotate around={-30:(0,0.5)}] (0,0.5) -- (-1,0.5) node [left] {$1$};
            \draw[thick,rotate around={30:(0,-0.5)}] (0,-0.5) -- (-1,-0.5) node[left] {$2$};
            \draw[thick, rotate around={30:(0,0.5)}] (0,0.5) -- (1,0.5) node [right] {$5$};
            \draw[thick] (0,-0.5) -- (0.8,-1);
            \node at (0.6,-0.5) {$\cO^{\prime}$};
            \draw[thick, rotate around={25:(0.8,-1)}] (0.8,-1) -- (1.8,-1) node [right] {$4$};
            \draw[thick, rotate around={-65:(0.8,-1)}] (0.8,-1) -- (1.8,-1) node [right] {$3$};
            \draw[thick] (0,0.5) -- (0,-0.5) node [midway, left] {$\cO^{\prime\prime}$};
        \end{scope}
    \etik 
\end{center}

\subsubsection{Step 4: Use OPE on $\Phi_2\cO^{\prime}\Phi_3\Phi_4$ in reverse}

\begin{center}
    \btik 
        \begin{scope}
            \node at (-1.75,0) {\Large{$\sum$}};
            \node at (-1.75,-0.5) {\small{$\cO^{\prime}, \cO^{\prime\prime}$}};
            \draw[thick, rotate around={-30:(0,0.5)}] (0,0.5) -- (-1,0.5) node [left] {$1$};
            \draw[thick,rotate around={30:(0,-0.5)}] (0,-0.5) -- (-1,-0.5) node[left] {$2$};
            \draw[thick, rotate around={30:(0,0.5)}] (0,0.5) -- (1,0.5) node [right] {$5$};
            \draw[thick] (0,-0.5) -- (0.8,-1);
            \node at (0.6,-0.5) {$\cO^{\prime}$};
            \draw[thick, rotate around={25:(0.8,-1)}] (0.8,-1) -- (1.8,-1) node [right] {$4$};
            \draw[thick, rotate around={-65:(0.8,-1)}] (0.8,-1) -- (1.8,-1) node [right] {$3$};
            \draw[thick] (0,0.5) -- (0,-0.5) node [midway, left] {$\cO^{\prime\prime}$};
            \node[right] at (2.5,0) {\Large{$=$}};
        \end{scope}
        \begin{scope}[xshift=5.75cm]
            \node at (-1.75,0) {\Large{$\sum$}};
            \node at (-1.75,-0.5) {\small{$\cO^{\prime\prime}$}};
            \draw[thick, rotate around={-30:(0,0.5)}] (0,0.5) -- (-1,0.5) node [left] {$1$};
            \draw[thick, rotate around={30:(0,0.5)}] (0,0.5) -- (1,0.5) node [right] {$5$};
            \draw[thick,rotate around={30:(0,-0.5)}] (0,-0.5) -- (-1,-0.5) node[left] {$2$};
            \draw[thick, rotate around={10:(0,-0.5)}] (0,-0.5) -- (1,-0.5) node [right] {$4$};
            \draw[thick, rotate around={-60:(0,-0.5)}] (0,-0.5) -- (1,-0.5) node [right] {$3$};
            \draw[fill=black] (0,-0.5) circle [radius=0.1cm];
            \draw[thick] (0,0.5) -- (0,-0.5) node [midway, left] {$\cO^{\prime\prime}$};
        \end{scope}
    \etik 
\end{center}

If we then relabel $\cO^{\prime\prime}\to \cO$ and compare the right-hand side of step 1, we see that crossing at 5-points follows automatically from crossing at 4-points. 

\br 
    We did this whole thing under our assumptions that everything was a scalar. Of course it is possible to extend the idea to operators with spin. We focused on the scalar case because it is easier in practice.
\er 

This above calculation allows us to introduce the following definition for a CFT.\footnote{Multiple people have come to this conclusion, including: Ferrara, Grillo, Gatto, '73; Polyakov '74; Mack '77.}

\bd[CFT] 
   We \textit{define} a CFT as the set of data satisfying this OPE associativity. 
\ed 

This leads to what is known as the conformal \textit{bootstrap programme}, which was used to completely solve some 2d theories analytically in the 1980s (e.g. so-called \textit{minimum models}, i.e. CFTs with a finite number of primaries). In 2008 there was a programme started by Rychkov known as \textit{numerical bootstrap programme} in higher dimensions. Here the idea is to put numerical constraints on the space of theories by examining this crossing equation. 

Let's sketch the idea now.\footnote{For more details see Section 10 (in particular Section 10.4) of Simmons-Duffin.} First consider the crossing equation for 4 identical scalars of weight $\Delta$. 

\begin{center}
    \btik 
        \begin{scope}[xshift=-4cm]
            \node at (-3,0) {\Large{$\sum C_{\Phi\Phi\cO}^2 $}};
            \node at (-3.65,-0.5) {\small{$\cO$}};
            \draw[thick, rotate around={-45:(0,0)}] (0,0) -- (-1,0) node [left] {$\Phi$};
            \draw[thick, rotate around={45:(0,0)}] (0,0) -- (-1,0) node [left] {$\Phi$};
            \draw[thick, rotate around={45:(1,0)}] (1,0) -- (2,0) node [right] {$\Phi$};
            \draw[thick, rotate around={-45:(1,0)}] (1,0) -- (2,0) node [right] {$\Phi$};
            \draw[thick] (0,0) -- (1,0) node [midway, above] {$\cO$};
            \node[right] at (2.75,0) {\Large{$=$}};
        \end{scope}
        \begin{scope}[xshift=4cm]
            \node at (-3,0) {\Large{$\sum C_{\Phi\Phi\cO}^2$}};
            \node at (-3.65,-0.5) {\small{$\cO$}};
            \draw[thick, rotate around={-30:(0,0.5)}] (0,0.5) -- (-1,0.5) node [left] {$\Phi$};
            \draw[thick,rotate around={30:(0,-0.5)}] (0,-0.5) -- (-1,-0.5) node[left] {$\Phi$};
            \draw[thick, rotate around={-30:(0,-0.5)}] (0,-0.5) -- (1,-0.5) node [right] {$\Phi$};
            \draw[thick, rotate around={30:(0,0.5)}] (0,0.5) -- (1,0.5) node [right] {$\Phi$};
            \draw[thick] (0,0.5) -- (0,-0.5) node [midway, right] {$\cO$};
        \end{scope}
    \etik 
\end{center}

It is now that our redefinition \Crefrange{eqn:ConformalBlock}{eqn:Pijkl} comes in handy. As we are considering 4 identical fields our $P_{ijk\ell}$ factor simply becomes 
\bse
    P_{ijk\ell} = |x_{ij}|^{\Delta} |x_{k\ell}|^{\Delta}.
\ese
The reason this is useful is because we note that the two diagrams above correspond to 
\bse 
    P_{1234} = |x_{12}|^{\Delta}|x_{34}|^{\Delta} \qand P_{1423} = |x_{14}|^{\Delta}|x_{23}|^{\Delta},
\ese 
which we can relate to the 4-point cross-ratios, \Cref{eqn:CrossRatios}, as
\bse 
    P_{1234} = |x_{13}|^{\Delta} |x_{24}|^{\Delta} u^{\Delta} \qand P_{1423} = |x_{13}|^{\Delta} |x_{24}|^{\Delta} v^{\Delta}.
\ese 

Next note that crossing symmetry is basically the statement that the result is invariant under $1\leftrightarrow 3$ (or $2\leftrightarrow 4$), which corresponds exactly to $u\leftrightarrow v$, which is how our two $P$s are related. If we then plug all this into \Cref{eqn:FourPointFunctionConformalBlocks}, we see the above equality of diagrams is equivalent to saying 
\bse 
    \sum_{\cO} C^2_{\Phi\Phi\cO} \big[ u^{2\Delta} G_{\cO}(x_1,x_2,x_3,x_4) - v^{2\Delta} G_{\cO}(x_3,x_2,x_1,x_4)\big] = 0
\ese 
for any $u,v$. 

This is the sort of equation people play around with numerically in order to solve $D\geq 3$ CFTs. The obvious question is how does this help in that goal? Well first we notice that $C^2_{\Phi\Phi\cO} >0$. This obviously places a constraint on the allowed set of primaries, i.e. the allowed values of $C_{\Phi\Phi\cO}$. The idea is to then consider the function in square brackets as a vector, so we're asking for a sum of vectors with non-negative coefficients to vanish:
\bse 
    \sum c \underline{v} = 0
\ese 
with $c\geq 0$. The case $c=0$ just means every correlator vanishes and so is boring, so we focus on $c > 0$. Clearly this has no solutions if all the vectors live on one side of a plane as we need to cancel stuff. That is, if we want to cancel a sum of the blue and green arrows in the diagram below we need something like the red arrow. 

\begin{center}
    \btik[scale=0.8]
        \draw[ultra thick, ->, red] (4,0.5) -- (4,-1.5);
        \draw[thick, rotate around={-25:(0,0)}, xscale=1.5, fill = gray!40, opacity = 0.8] (0,0) .. controls (0.8,0.5) and (1.2,0.5) .. (3.5,1) .. controls (4,1.5) and (4,3) .. (4.5,4.5) .. controls (3.2,4) and (3.7,4) .. (1,3.5) .. controls (0.5,3) and (0.5,1.5) .. (0,0);
        \draw[thick, rotate around={-25:(0,0)}, xscale=1.5] (0,0) .. controls (0.8,0.5) and (1.2,0.5) .. (3.5,1) .. controls (4,1.5) and (4,3) .. (4.5,4.5) .. controls (3.2,4) and (3.7,4) .. (1,3.5) .. controls (0.5,3) and (0.5,1.5) .. (0,0);
        \draw[ultra thick, ->, blue, rotate around={10:(4,0.5)}] (4,0.5) -- (4,2.5);
        \draw[ultra thick, ->, green, rotate around={-40:(4,0.5)}] (4,0.5) -- (4,2.5);
        \draw[fill=black] (4,0.5) circle [radius=0.05cm];
    \etik
\end{center}

\noindent These vectors themselves are the data you are given in order to try solve the CFT. 

% For Paul: you made some comment about AdS/CFT here about how bootstrap is used to do strong coupled CFT calculations. I didn't catch everything you said in when I typed it in the lecture and I'm not confident enough with AdS/CFT to expand on this remark myself. 
\part{CFT in $2D$}
\chapter{CFT in 2D}

We now move on to the discussion of CFTs in 2D. As we have said multiple times throughout these notes, CFTs in 2-dimensions are conceptually very different from higher dimensional CFTs. One of the biggest differences, which we will continue to emphasise in this part of the course, is that in 2D we have an \textit{infinite} number symmetries. The fact that this is true can be seen going all the way back to the first exercise, which asked for a derivation of $\xi^{\mu}$: during one step of the calculation you should have got something of the form 
\be 
\label{eqn:D-2Kappa}
    (D-2) \p_{\mu} \p_{\nu} \kappa = 0,
\ee 
which was used to say that $\kappa$ was linear in $x$. However if we set $D=2$ then we can have \textit{any} function of $x$. This clearly gives us conceptually different results, which this part of the course will explain. 

Before moving on, it is worth clarifying a fact that might seem backwards at first.\footnote{At least it took me a while to get my head around this.} The fact that we have more symmetries puts \textit{more} constraints on the CFT (i.e. more symmetries need to be obeyed) and so actually gives us a \textit{better} grasp on the problem. 

\section{Introduction To String Theory}

We start this part of the notes with a brief introduction to string theory. We will not go into much depth of the physics of string theory\footnote{This will be done on the string theory course, of course.} but just use it as a base to study the mathematics of 2D CFTs. 

\subsection{The Relativistic Particle}

The standard way to introduce string theory is to begin by considering the relativistic particle. In this case we think of the coordinates $X^{\mu}$ as being maps 
\bse 
    \begin{split}
        X^{\mu} : \R & \to \mathbb{M}^{1,D-1} \\
        \tau & \mapsto X^{\mu}(\tau),
    \end{split}
\ese 
where $\mathbb{M}^{1,D-1}$ is $D$-dimensional Minkowski spacetime. This is just the statement that for a given proper time, $\tau$, we have a given event in Minkowski spacetime. This just gives us the familiar \textit{worldline} of the particle. 
\begin{center}
    \btik
        \draw[thick, ->] (-8,2.5) -- (-4,2.5);
        \draw[thick] (-7.5,2.7) -- (-7.5,2.3);
        \draw[thick] (-7,2.7) -- (-7,2.3);
        \draw[thick] (-6.5,2.7) -- (-6.5,2.3);
        \draw[thick] (-6,2.7) -- (-6,2.3);
        \draw[thick] (-5.5,2.7) -- (-5.5,2.3);
        \draw[thick] (-5,2.7) -- (-5,2.3);
        \draw[thick] (-4.5,2.7) -- (-4.5,2.3);
        \node at (-3.8, 2.3) {$\R$};
        \draw[->] (-5,3) .. controls (-3,4) and (-1,4) .. (1,3) node [midway, above] {$X^{\mu}$};
        \draw[ultra thick, ->] (-0.5,0) -- (3,0);
        \node at (2.8,-0.3) {$\Vec{x}$};
        \draw[ultra thick, ->] (0,-0.5) -- (0,5);
        \node at (-0.3,4.8) {$x^0$};
        \draw[thick] (0.5,0.5) .. controls (3,3) and (1,3) .. (2,4.5);
        \draw[rotate around={35:(0.6,0.6)}] (0.6,0.4) -- (0.6,0.8);
        \draw[rotate around={38:(1,1)}] (1,0.8) -- (1,1.2);
        \draw[rotate around={40:(1.4,1.5)}] (1.4,1.3) -- (1.4,1.7);
        \draw[rotate around={45:(1.65,2)}] (1.65,1.8) -- (1.65,2.2);
        \draw[rotate around={70:(1.8,2.6)}] (1.8,2.4) -- (1.8,2.8);
        \draw[rotate around={100:(1.75,3.3)}] (1.75,3.1) -- (1.75,3.5);
        \draw[rotate around={70:(1.75,4)}] (1.75,3.8) -- (1.75,4.2);
    \etik
\end{center}

The action for the relativistic particle in terms of the proper time $\tau$ is
\be 
\label{eqn:ParticleAction}
    S = -m \int d\tau \, \sqrt{\dot{X}_{\mu} \dot{X}^{\mu}}, \qquad \text{with} \qquad \dot{X}^{\mu} := \frac{\p X^{\mu}}{\p \tau}.
\ee    

We now stress an important, at first confusing, point. The $\mu = 0, ... , D-1$ index comes from the fact our Minkowski spacetime is $D$ dimensional and so we need to specify how the worldline "lies" in there.\footnote{More technically, we need to specify how our embedding is defined.} However our action is just a 1-dimensional QFT and so the $\mu$ there is just viewed as labelling $D$ different scalar fields on our 1D space. That is $X^1(\tau)$ is just a scalar field as far as \Cref{eqn:ParticleAction} is concerned. 

\badr 
\label{rem:ProperTimeGauge}
    There is actually a subtle point relating to going from integrating over coordinate time $t$ to integrating over proper time $\tau$. When we go to proper time, we "promote" the coordinate time $t$ to a degree of freedom, $\tau$, in order to give a manifestly Lorentz invariant theory. However $t$ is \textit{not} a real degree of freedom (you cannot travel in time as you like), and so really $\tau$ is a \textit{gauge} degree of freedom. This discrimination between a real degree of freedom and an gauge degree of freedom is tied up in the fact that our action is reparameterisation invariant $\tau \to \widetilde{\tau}(\tau)$. This basically tells us that $\tau$ has no physical meaning. This point is made here because the same idea applies to the coordinates on the string, where it is easy to get confused about what's what. For a bit more detailed discussion of this point, see the start of Prof. Tong's string theory notes.\footnote{Or, shameless plug, the start of my string theory notes on Shiraz Minwalla's course.}
\eadr 

\subsection{String Actions}

We now want to extend this logic for a free relativistic particle to a free propagating string. The idea is that instead of having a particle tracing out a worldline, we have a string tracing out a world\textit{sheet}. This string can either be closed (which traces out a cylinder shape) or open (which just gives a sheet, like a piece of paper). In our discussion of string theory we shall focus on the closed string. 

Our worldsheet is now parameterised using two coordinates; the proper time, $\tau$, and another which tells us about moving \textit{around} the sheet. It is standard to label the latter as $\sig$. Our $X^{\mu}$ map then becomes 
\bse 
    \begin{split}
        X^{\mu} : \R^{1,1} & \to \mathbb{M}^{1,D-1} \\
        (\tau,\sig) & \mapsto X^{\mu}(\tau,\sig).
    \end{split}
\ese 
We call the domain and target of this map the "\textit{worldsheet}" and "\textit{target space}", respectively. That is 
\bse 
    \text{Worldsheet} = \R^{1,1}, \qand \text{Target Space} = \mathbb{M}^{1,D-1}.
\ese

Pictorially, this map acts as follows. 
\begin{center}
    \btik
        \draw[thick, ->] (-8.5,0.5) -- (-4,0.5);
        \node at (-3.8,0.2) {$\sig$};
        \draw[thick, ->] (-8,0) -- (-8,4.5);
        \node at (-8.3,4.2) {$\tau$};
        %
        \draw[->] (-5.5,3) .. controls (-3.5,4) and (-1.5,4) .. (0.5,3) node [midway, above] {$X^{\mu}$};
        %
        \draw[ultra thick] (-0.5,0) -- (4,0);
        \node at (3.8,-0.3) {$\Vec{x}$};
        \draw[ultra thick] (0,-0.5) -- (0,5);
        \node at (-0.3,4.8) {$x^0$};
        \draw[thick] (0.8,0.5) .. controls (2.3,3) and (0.3,3) .. (0.8,4.5);
        \draw[thick] (2.3,0.5) .. controls (4.8,2.5) and (0.8,3) .. (2.3,4.5);
        \draw[thick] (1.55,4.5) ellipse (0.74 and 0.18);
        \draw[thick] (0.8,0.5) arc (180:360:0.74 and 0.18);
        \draw[thick, dashed] (0.8,0.5) arc (180:360:0.74 and -0.18);
        \draw[thick, dashed] (2.22,2) ellipse (0.87 and 0.2);
        \draw[->, thick] (3.3,2) .. controls (3.1,2.5) .. (2.6,3);
        \node at (3.3,2.5) {$\tau$};
        \draw[->, thick] (1.5,1.7) .. controls (2.25,1.5) .. (2.9,1.7);
        \node at (2.25, 1.4) {$\sig$};
    \etik
\end{center}

\badr 
    It is at this point that \Cref{rem:ProperTimeGauge} becomes important. As we have done in the diagram above, we often talk about about $\tau$ and $\sig$ in our Minkowski spacetime as "going up the string" and "going around the string". However these are the analogue of the proper time for the worldline and have \textit{no} physical meaning. What $\tau$ and $\sig$ label are the coordinates on the worldsheet (i.e. the left-hand side of the diagram above). However it is often more intuitive to think of how $(\tau,\sig)$ live in target space as it gives us an idea of the worldsheet being the thing on the right-hand side of our figure, but we just \textit{defined} the worldsheet to be our $\R^{1,1}$. A bit more technically the thing on the right-hand diagram is an \textit{embedding} of the worldsheet into $\mathbb{M}^{1,D-1}$. This is completely analogous to differentiating between the abstract $2$-sphere $S^2$ and the embedding into $\R^3$, which is our intuitive notion of a ball. 
\eadr 

We now want to define a QFT around this, and in order to do that we need an action. We again want to use the particle action, \Cref{eqn:ParticleAction}, as a guiding light. It should be familiar from a GR course that this action just gives the length of the curve. This is an absolute property of the worldline, i.e. it doesn't depend on how we choose to do the embedding. We therefore try to extend this idea to say that our string action should give us the \textit{area} of the worldsheet. 

\subsection{Nambu-Goto Action}

The question is "how do we do this?" Well recall from GR that we measure lengths and angles using the metric. We therefore need some kind of metric on the worldsheet. We do this by introducing what is known as the \textit{induced metric}. This is given by
\be 
\label{eqn:InducedMetric}
    h_{ab} = \p_a X^{\mu} \p_b X_{\mu}, \qquad a, b \in \{1,2\},
\ee 
where we have used the fact that our target space is Minkowski and so has the flat metric $\eta_{\mu\nu}$. You obtain this result by considering the pullback of the target space metric onto the embedding of the worldsheet. We do not go through the details here but refer intrested readers to Section 1.2 of Prof. Tong's notes.\footnote{Again or 1.3.1 of my notes on Prof. Minwalla's course.} 

We can use this result to get the following action, known as the \textit{Nambu-Goto} action.
\mybox{
    \be 
    \label{eqn:NambuGotoAction}
        S_{NG} = -\frac{1}{2\a'} \int d^2 x \sqrt{-\det(h_{ab})},
    \ee 
}
\noindent where the brackets around $(h_{ab})$ tell us to consider the matrix whose entries are given by $h_{ab}$. The fact that this gives the area can be seen by considering the area of `infintesimal tiles' in a Euclidean space.\footnote{Again for more detials see either Prof. Tong's notes or my notes.} The $\a^{\prime}$ appearing in \Cref{eqn:NambuGotoAction} is known as the \textit{Regge slope}, and simple dimensional analysis tells us $[\a^{\prime}] = [L^2]$. We can  relate it to the, more physical, tension of the string via
\bse 
    \frac{1}{2\pi\a'} = T.
\ese

\br 
    In order to understand how this tension fits in to our pictorial intuition, we should really think of the strings as a rubber band trying to shrink down to zero radius. The tension then corresponds to its resistance to being stretched out again.
\er 

\bbox 
    Vary the Nambu-Goto action w.r.t. $X^{\mu}$ to arrive at the equations of motion 
    \be 
    \label{eqn:NambuGotoEOM}
        \p_a \big( \sqrt{-h} h^{ab}\p_b X^{\mu}\big) = 0,
    \ee 
    where $h:=\det(h_{ab})$.
\ebox 

Let's see what symmetries our Nambu-Goto action has.
\ben[label=(\roman*)]
    \item The spacetime Poincar\'{e} generators, ${\Lambda^{\mu}}_{\nu}$ and $a^{\mu}$, are independent of the worldsheet coordinates,\footnote{This is what we mean about $\mu$ just labelling the different fields from the worldsheet's perspective.} so we have spacetime Poincar\'{e} symmetry from the fact that the indices are contracted. This is genuine \textit{global} symmetry from the worldsheet's perspective. 
    \item Just as we had the gauge reparameterisation symmetry $\tau\to\widetilde{\tau}(\tau)$ for the worldline, we have the gauge reparameterisation symmetries $\tau\to\widetilde{\tau}(\tau)$ and $\sig\to\widetilde{\sig}(\sig)$ for the worldsheet. Again it is important to note that these are conceptually different from the Poincar\'{e} symmetries; they are \textit{local} gauge symmetries.
\een 

\badr 
    An important repercussion of the above statements is that the Poincar\'{e} symmetries, being global symmetries, will give rise to Noether currents. However the gauge symmetries will not. 
\eadr 

\subsection{Polyakov Action}

This is great, but this is meant to be a CFT course and our action is clearly not a CFT; it doesn't have the needed Weyl invariance! Why are we talking about it, then? Well we note the following "theorem"
\bt 
    Square roots nasty things to have in actions, and make quantisation very hard.\footnote{This is bascially because in the path integral formulation we get the propagators by rewritting the action in the form $\int \phi A \phi$, and then take the inverse of this operator to find the propagator (see, e.g., my notes on Dr. Nabil Iqbal's QFT II course for more details). This is hard to do for multiple reasons if we have a pesky square root.}
\et 

We would, therefore, like to remove this square root. In order to do that, we invoke a neat trick that finds use when studying the relativistic particle.\footnote{This presentation is taken, almost directly, from Prof. Tong's string theory notes.} The idea is that we can rewrite \Cref{eqn:ParticleAction} in the following form 
\bse 
    S = \frac{1}{2}\int d\tau \big( e^{-1} \dot{X}^{\mu}\dot{X}_{\mu} - em^2\big),
\ese
where $e=e(\tau)$ is some new field we introduce. We can make this look more familiar by switching to notation 
\bse 
    e = \sqrt{-\g_{\tau\tau}},
\ese 
so that our action becomes 
\bse 
    S = \frac{1}{2}\int d\tau \, \sqrt{-g_{\tau\tau}} \big( g^{\tau\tau} \dot{X}^{\mu}\dot{X}_{\mu} + m^2\big) \qquad g^{\tau\tau} := (g_{\tau\tau})^{-1}.
\ese 
This now looks like we have coupled our scalar fields $X^{\mu}(\tau)$ to a 1D gravity, which is where the name "einbein" comes from.\footnote{In GR we have something called a \textit{vierbein}.} 

\br 
    It would be fair to raise concern about $e(\tau)$ having a drastic impact on the theory. However, as a quick calculation will show, it is completely fixed by the equations of motion and so we don't need to panic. 
\er 

Ok so we want to take this argument and try remove the square root from our Nambu-Goto action. This gives us what is known as the \textit{Polyakov} or \textit{Brink-DiVecchia-Howe-Deser-Zumino} action.\footnote{It was the latter that discovered the action. It was Polyakov was responsible with understanding how to manipulate it in a path integral, and so it is normally given by his name.} 
\mybox{
    \be 
    \label{eqn:PolyakovAction}
        S_P = -\frac{1}{4\pi\a'} \int d^2 x \sqrt{-\g} \g^{ab} \p_a X^{\mu} \p_b X_{\mu}. 
    \ee 
}
\noindent Our new field here is $\g_{ab}$, with inverse $\g^{ab}$. We have also used the standard notation $\g := \det(\g_{ab})$. First let's check this is actually (classically) equivalent to the Nambu-Goto action. Computing the Euler Lagrangian equations for $\g$, we find 
\bse
    0 = \p_a X^{\mu} \p_b X_{\mu} - \frac{1}{2}\g_{ab}\g^{cd} \p_c X^{\mu} \p_d X_{\mu} = h_{ab} - \frac{1}{2}\g_{ab}\g^{cd}h_{cd}
\ese 
which can be manipulated into the form
\bse 
    \sqrt{-\g} \g^{ab} \p_a X^{\mu} \p_b X_{\mu} = \sqrt{- h}.
\ese 

\bbox 
    Fill in the missing lines in the above calculation. \textit{Hint: Recall that $\det(aM) = a^2\det M$ for any scalar $a$ and matrix $M$.}
\ebox 

Ok that's good, but what is this new field $\g_{ab}$? Well if we compute the equations of motion w.r.t $X^{\mu}$ we simply get 
\bse 
    \p_a \big( \sqrt{-\g} \g^{ab}\p_a X^{\mu}\big) = 0,
\ese 
which is exactly the same as the Nambu-Goto equations of motion, \Cref{eqn:NambuGotoEOM}, but with $h\to \g$. This tells us that we can \textit{think} of $\g_{ab}$ as the induced metric in the Polyakov case.

\badr 
    The emphasis on "think" above is for the following reason. The induced metric $h_{ab}$ is a genuine metric, given by the pullback metric from the spacetime metric $\eta_{\mu\nu}$. However $\g_{ab}$ is a \textit{field}, and is therefore dynamical. As with the einbein $e$, its dynamics are completely fixed by the equations of motion, however we should not be so quick to say it is a legitimate metric. We therefore refer to $\g_{ab}$ as the \textit{dynamical metric} on the worldsheet. We will, however, just think of it as a `normal' metric from now on.
\eadr 

The Polyakov action is the action of free scalar fields $X^{\mu}(\sig,\tau)$ (from the 2-dimensional point of view) coupled to a background metric. It is a free field as the $X^{\mu}$ equations of motion are 
\bse 
    \nabla^2 X^{\mu} = 0 
\ese 
where $\nabla$ is the covariant derivative in 2D w.r.t. $\g_{ab}$. 

\subsection{Symmetries Of Polyakov Action}

The important thing about the Polyakov action for us are it's symmetries. It still has the spacetime Poincar\'{e} and worldsheet reparameterisation symmetries (which we will call worldsheet diffeos from now on) present in the Nambu-Goto action. However it also possess a Weyl symmetry given by 
\bse 
    \g_{ab} \to \Omega^2(x) \g_{ab}.
\ese 
This is easily seen by the fact that
\bse 
    \det\g \to \Omega^4(x)\det\g , \qand \g^{ab} \to \Omega^{-2}(x)\g^{ab},
\ese 
so that 
\bse 
    \sqrt{-\g} \g^{ab} \to \sqrt{-\g} \g^{ab}.
\ese 
These are \textit{local} symmetries, and so we can gauge fix them. 

Now $\g_{ab}$ is a $2\times 2$ symmetric matrix and so has $3$ independent components. However we have just said that we have 2 worldsheet reparameterisation gauge symmetries, and so we can use a gauge transformation to fix $\g_{ab}$ so that it is of diagonal form. We can then use a Weyl transformation to give us a flat worldsheet metric:
\bse 
    \g_{ab} \, \overset{\text{Diffeos.}}{\longrightarrow} \, \Omega^2(x)\eta_{ab} \, \overset{\text{Weyl}}{\longrightarrow} \, \eta_{ab}.
\ese 

\br 
    We can also see this result by considering the Riemann tensor $R_{\mu\nu\rho\sig}$. In 2D it reduces down to $\frac{R}{2}(g_{\mu\rho} g_{\nu\sig} - g_{\mu\sig}g_{\nu\rho})$. This tells us that essentially the Riemann tensor goes like $R$, which is just a single scalar degree of freedom, which we can make flat using the Weyl transformations. 
\er 

Now comes the important result for us: there is still a left over residual gauge symmetry after we fix $\g_{ab}=\eta_{ab}$. That is, even after picking this gauge, there is still some left over gauge invariance. It is not too hard to see that this gauge invariance is given \textit{precisely} by conformal transformations. That is any worldsheet diffomorphism that only changes $\g_{ab}$ by a Weyl factor is a left over residual gauge. This allows us to make the all important conclusion

\mybox{
    \begin{center}
        String theory is a 2D CFT!
    \end{center}
}

\badr 
    The fact that the Polyakov action possess the additional Weyl invariance tells us that if we want to continue studying the Nambu-Goto action through the Polyakov one, we actually have to take an equivalence class on our solutions, i.e. we mod out by Weyl transformations. In other words, two solutions that differ purely by a Weyl transformation are considered the same solution. The Weyl invariance is obviously true at the classical level, and it places some significant restraints on the allowed terms in the action. For example we couldn't include a cosmological constant term as then there wouldn't be a $\g^{ab}$ factor to cancel the transformation from the determinant. Things become a bit more subtle when it comes to the quantum theory, and more details about this can be found in my notes on Prof. Shiraz Minwalla's string theory course. 
\eadr 

\br 
    Note that in higher dimensions we were making the \textit{choice} to consider a flat space for our CFT. However for string theory we can actually always do this. 
\er 

This concludes our discussion of string theory itself, and we will now study abstract 2D CFTs. A lot of the work that follows will be very similar to the work we did before for higher dimensional CFTs, but there are some important differences, that we shall point out along the way. Of course the stuff we talk about for the rest of the course will be applicable to string theory (as it is a 2D CFT), but it is important to note that it is not string theory specific. 

\bnn 
    As we are now moving on from string theory our $\mu$ indices that follow take two values. That is they label the worldsheet coordinates $(\tau,\sig)$ not some higher dimensional spacetime. 
\enn 

\section{Conformal Transformations In 2D}

As we said at the start of this part of the notes, the space of conformal transformations is \textit{much} bigger, which came from the fact that the $D-2$ on the left-hand side of \Cref{eqn:D-2Kappa} identically vanishes. This tells us that $\kappa$ can be an arbitrary polynomical in $x$, rather than quadratic as for the $D>2$ case.

We can see this result in another useful way by considering so-called \textit{light-cone coordinates}
\be 
\label{eqn:LightconeCorrdinates}
    x^{\pm} := \frac{1}{\sqrt{2}} (\sig \pm \tau). 
\ee 
The line element in these coordinates becomes
\bse 
    ds^2 = d\sig^2 - d\tau^2 = 2dx^+ dx^- = \begin{pmatrix}
        dx^+ & dx^-
    \end{pmatrix} \begin{pmatrix}
        0 & 1 \\
        1 & 0
    \end{pmatrix} \begin{pmatrix}
        dx^+ \\
        dx^-
    \end{pmatrix}
\ese 
from which we can read off the worldsheet metric (which we denote by $g$ from now on) as 
\bse 
    (g_{ab}) = \begin{pmatrix}
        0 & 1 \\
        1 & 0
    \end{pmatrix}.
\ese 
Now consider the general transformation
\bse 
    x^- \mapsto f(x^-) \qand x^+ \mapsto g(x^+) 
\ese
where $f$ and $g$ can \textit{any} functions of only their respective arguments, i.e. $f$ doesn't depend on $x^+$ at all and similarly $g$ doesn't depend on $x^-$. Then our metric transforms as
\bse 
    2dx^+ dx^- \to  2\frac{df(x^-)}{dx^-}\frac{dg(x^+)}{dx^+} dx^-dx^+,
\ese 
but this is just a Weyl transformation, as the metric is just scaling under this transformation. This is therefore a conformal transformation for \textit{any} functions $f,g$. This gives us an infinte class of symmetries for the $D=2$ case. In higher dimensions things are a bit more complicated,\footnote{Note in higher dimensions going to light-cone coordinates would break Lorentz invariance as we would have to single out two of our directions in order to define $x^{\pm}$. However for $2D$ we're alright because there are only two coordinates.} and, as we have seen already, we get at most quadratic terms.

We now abandon light-cone coordinates and look at the much more useful complex coordinates, i.e. we define
\bse 
    z := x^+ \qand \Bar{z} := x^-. 
\ese 
There are three cases:
\ben[label=(\roman*)] 
    \item Real Minkowski space: $\sig,\tau\in\R$ and therefore $z,\Bar{z}\in\R$ and independent. 
    \item Complexify to get Complex Minkowski: $\sig,\tau\in\C$ and $z,\Bar{z}\in\C$ and independent. 
    \item Real Euclidean, $\tau\to i\tau$: $\sig,\tau\in\R$ but $z,\Bar{z}\in \C$ with $\Bar{z} = z^*$.
\een 

Case (iii)  is often most useful as it gives us a direct relation between $z$ and $\bar{z}$, namely complex conjugation.\footnote{Note it's for this reason that we call the second variable $\bar{z}$. That is for cases (i) and (ii) the bar does not indicate complex conjugation, but here it does.} This essentially allows us to derive all the results by just considering the $z$ case and then obtain the $\bar{z}$ results by complex conjugation. For this reason, this is the case we will use in this course.

Our conformal transformations are then just simply meromorphic\footnote{It is likely I will use mermorphic and holomorphic interchangeabley. These are obviously technically different, but close enough that hopefully the ideas are understood.} transformations
\bse 
    z^{\prime} = f(z) \qand \bar{z}^{\prime} = g(\bar{z}).
\ese 

\badr 
    Note although we still have an infinite space of conformal transformations, it is still \textit{drastically} smaller than the original diffeomorphisms. This is because our $f$ and $g$ are only functions of a single variable. The easiest way to see this is to consider taking a Taylor expansion of functions that depends on both variables $F(\tau,\sig)$ and $G(\tau,\sig)$. If we expand $F(\tau,\sig)$ about $\sig$, only \textit{the first term} in this expansion is $\sig$ independent, and similarly only the first term in the expansion of $G(\tau,\sig)$ about $\tau$ will be $\tau$ independent. Therefore essentially we have truncated our space of diffeomorphisms down to a first order.
\eadr 

\section{Generators: Witt Algebra}

We can write our conformal transformations infinitesimally as 
\bse 
    z^{\prime} = z + \xi(z)  
\ese 
where $\xi(z)$ is any meromorphic function of $z$, which we can expand as
\bse 
    \xi(z) = \sum_n a_n z^{n+1}.
\ese 
Our conformal transformation is then generated by 
\bse 
    \xi(z) \p_z = \sum_n a_n z^{n+1}\p_z = -\sum_a a_n L_n
\ese 
where 
\mybox{
    \be 
    \label{eqn:Ln}
        L_n := -z^{n+1}\p_z
    \ee
}
\noindent are the generators. If these are to be generators, they are meant to be elements of the Lie algebra, and so we better check that they close under the Lie bracket, which is just the commutator. By direct calculation we have
\bse 
    [L_n,L_m] = [z^{n+1}\p_z, z^{m+1}\p_z] = ... = (m-n) z^{n+m+1}\p_z,
\ese 
which we can rewrite as
\bse 
    [L_n, L_m] = (n-m) L_{n+m},
\ese
where we have used the minus sign to flip $(m-n)$. As we mentioned above, we can easily obtain the barred version by taking the complex conjugate (i.e. put bars everywhere). That is repeating the above for the $\bar{z}$s will give us
\bse 
    [\bar{L}_n, \bar{L}_m] = (n-m) \bar{L}_{n+m}.
\ese 

This is an important result, so we write it again in a nice box. 
\mybox{
    \be 
    \label{eqn:WittAlgebra}
        [L_n,L_m] = (n-m) L_{n+m} \qand [\bar{L}_n,\bar{L}_m] = (n-m) \bar{L}_{n+m}.
    \ee 
}
\noindent This is known as the \textit{Witt algebra}. 

\subsection{M\"{o}bius Subgroup}

The Witt algebra at first seems completely new, however after staring at it for a moment or two, we see that it has a subalgebra spanned by
\bse 
    L_0, L_1, L_{-1} \qquad \text{or} \qquad \bar{L}_0, \bar{L}_1, \bar{L}_{-1}
\ese 
depending on which part of the algebra we're looking at. Explicitly the commutators are 
\bse 
    [L_{-1}, L_1] = -2L_0 \qquad [L_0,L_1] = -L_1 \qand [L_0,L_{-1}] = L_{-1},
\ese
and similarly for the barred expressions. We can then construct a Lie algebra isomorphism between this subalgebra and the $\mathfrak{sl}(2)$ algebra.\footnote{For an explicit statement see Section 6.1.1 of my notes on Prof. Shiraz Minwalla's course.}

We can translate this Lie algebra isomorphism to a Lie group isomorphism. That is there is a subgroup of our $2D$ conformal group that is isomorphic to $SL(2,\C)$. These are the subset of transformations which are globally well defined on the whole Riemann sphere. In other words, when we have $n<-1$ our generators $L_n/\bar{L}_n$ go with negative powers of $z$, and so are ill-defined at $z=0$. Similarly for $n>1$ they are ill-defined at $z=\infty$. The latter condition might seem a bit strange, but we can see it by considering an inversion, then 
\bse 
    z^n \mapsto \frac{1}{z^n}, \qand \infty \mapsto 0,
\ese 
so they are ill-defined.

So what are the subgroup generators?
\ben[label=(\roman*)]
    \item $L_{-1}, \bar{L}_{-1}$ are translations, as they are just $\p_z, \p_{\bar{z}}$. In other words they correspond to $P_{\mu}$. 
    \item $L_0, \bar{L}_0$ are dilatations. This just just because $D = z\p_z + \bar{z}\p_{\bar{z}}$, and so can be generated by $L_0 = z\p_z$ and $\bar{L}_0 = \bar{z}\p_{\bar{z}}$. We can also get the Lorentz transformations\footnote{Note that the indices only take values $\{0,1\}$ as we are in 2D.} $L_{01} = -L_{10} = z\p_{\bar{z}} - \bar{z}\p_z$.
    \item $L_1, \bar{L}_1$ give us the special conformal transformations $K_{\mu}$. We can see this by a similar argument to (ii) above. 
\een 

This global subgroup is called the \textit{M\"{o}bius group}. For a finite transformation it is given by\footnote{Hopefully it's reasonably clear why this is isomorphic to $SL(2,\C)$ from here.}
\bse 
    z \mapsto \frac{az+b}{cz+d} \qquad \text{with} \qquad ad-bc = 1.
\ese 
We can get the infinitesimal version by setting
\bse 
    a = 1 + \a, \qquad b = 1+\beta, \qquad d = 1 + \del, \qand c = 1+\g, 
\ese 
where the $ad-bc=1$ condition gives us $\a = - \del$. This will reproduce the 3 generators $\{L_0,L_1,L_{-1}\}$, as it must as the Lie algebra is given by working infinitesimally around the identity. 

\section{Transformations Of Fields}

We now want to introduce fields into our CFT. A general field will depend on both variables, $\varphi(z,\bar{z})$, however there are also fields which only depend on one. We call a field that only depends on $z$, $\varphi(z)$, \textit{Chiral} and similarly a field $\varphi(\bar{z})$ \textit{antichiral}.

The first thing we note is that we can consider the scaling of $z$ and $\bar{z}$ separately, i.e. $z\to \l z$ and $\bar{z}\to \bar{\l}\bar{z}$. Our fields will obviously scale by both these, and so we write the weights of the fields as a doublet $(h,\bar{h})$. That is a general field of weight $(h,\bar{h})$ transforms as
\bse 
    \varphi \to \varphi^{\prime}, \qquad \varphi^{\prime}(\l z, \bar{\l}\bar{z}) = \l^{-h}\bar{\l}^{-\bar{h}} \varphi(z,\bar{z}). 
\ese 

Now we note that this combines both dilatations and rotations. We can see this by letting $\l = re^{i\theta}$ and $\bar{\l} = r e^{-i\theta}$, then we have
\bse 
    \varphi^{\prime}(\l z, \bar{\l} \bar{z} ) = r^{-(h+\bar{h})} e^{-i\theta(h-\bar{h})} \varphi(z,\bar{z}).
\ese
The $r$ factor is just our dilatation scaling, while the $\theta$ term is a rotation in the $\C$ plane. We therefore make the following definitions/conclusions. 
\mybox{
    \be 
        \Delta = h + \bar{h} \qand S = h-\bar{h},
    \ee 
}
\noindent where $\Delta$ is the dilatation weight and $S$ is the \textit{spin} of the field. 

Ok so how does a field transform under a \textit{general} conformal transformation? Again we answer this by \textit{defining} what we mean by a primary field and then move on to descendants. 
\bd[Primary Field (2D)] 
    A \textit{primary field} of weight $(h,\bar{h})$ transforms as $\varphi\to\varphi^{\prime}$ where 
    \be 
    \label{eqn:PrimaryField2D}
        \varphi^{\prime}\big(f(z),\bar{f}(\bar{z})\big) = \big(\p_z f\big)^{-h} \big(\p_{\bar{z}}\bar{f}\big)^{-\bar{h}} \varphi(z,\bar{z}).
    \ee 
\ed 
\noindent If we then restrict to the M\"{o}bius subgroup we can find how a field transforms in the same way as $D>2$ case before. 

An infinitesimal transformation
\bse 
    z \mapsto z + \xi(z), \qand \bar{z} \mapsto \bar{z} + \bar{\xi}(\bar{z}),
\ese 
of primary of weight $(h,\bar{h})$ is then 
\be 
\label{eqn:InfintesimalPrimaryTransformation2D}
    \del \varphi = -h \frac{\p \xi}{\p z} \varphi - \xi \frac{\p\varphi}{\p z} -\bar{h} \frac{\p \bar{\xi}}{\p \bar{z}} \varphi - \bar{\xi} \frac{\p\varphi}{\p \bar{z}}
\ee 

\br 
    Note, as we tried to emphasise at the start of this chapter, this is a much \textit{stronger} requirement than in higher-dimensions. That is, we have more constraints and so it gives us a smaller set of primary fields. 
\er 

With the above remark in mind, we can define

\bd[Quasi-primary Field]
    A \textit{quasi-primary field} of weight $(h,\bar{h})$ transforms under an infinitesimal transformation as \Cref{eqn:InfintesimalPrimaryTransformation2D} \textit{but} it need only hold for the global M\"{o}bius subgroup, i.e. for transformations were $\xi(z)$ is quadratic in $z$.
\ed 

To clarify, a quasi-primary field \textit{must} transform as \Cref{eqn:InfintesimalPrimaryTransformation2D} when $\xi(z)$ is quadratic in $z$, but it can transform arbitrarily for higher polynomials. It follows from this that all primary operators are quasi-primary, but the reverse is not true. This is exactly what we mean in the remark above; there are \textit{less} primary operators then there are quasi-primary ones.

\br 
    Note in $D>2$ CFTs quasi-primaries and primaries become indistinguishable. This is purely because we only have quadratic transformations for $D>2$. 
\er 

\section{Stress-Energy Tensor}

As we are still considering a CFT the stress energy tensor still obeys
\bse 
    \p_{\mu}T^{\mu\nu} = 0 \qand {T^{\mu}}_{\mu} = 0.
\ese
The stress-tensor is symmetric, and so we have 3 independent components,
\bse 
    T^{zz} \qquad T^{\bar{z}\bar{z}} \qand T^{z\bar{z}} = T^{\bar{z}z}
\ese 
to start with. The question now is "what do our constraints above tell us about these components?" First consider the traceless condition. Recall that our metric is off-diagonal form, and so it essentially trades $z\leftrightarrow \bar{z}$ when we raise/lower indices. We therefore have 
\bse 
    T^{\mu\nu}g_{\mu\nu} = 2 T^{z\bar{z}} = 0. 
\ese 
This removes one of our components, leaving us with just 2. Now let's look at current conservation
\bse 
    \p_{\mu}T^{\mu\nu} = 0.
\ese 
This is actually two equations, given by $\nu=z,\bar{z}$. Let's consider these separately 
\ben[label=(\roman*)]
    \item $\nu=z$: 
    \bse 
        \p_z T^{zz} +\p_{\bar{z}} T^{\bar{z}z} = 0 \qquad \implies \qquad \p_z T^{zz} = 0,
    \ese 
    where we have used the traceless condition. 
    \item $\nu=\bar{z}$: Similarly we get
    \bse 
        \p_{\bar{z}} T^{\bar{z}\bar{z}} = 0.
    \ese 
\een 
We therefore see that the $T^{zz} = T^{zz}(\bar{z})$ and $T^{\bar{z}\bar{z}} = T^{\bar{z}\bar{z}}(z)$. This seems a little unappealing; it would be nicer if $T$ with $z$ indices was only $z$ dependent, and similarly $T$ with $\bar{z}$ indices. Well if we again use the fact that the metric is off-diagonal we can lower the indices while swapping $z\leftrightarrow \bar{z}$, to define
\be 
\label{eqn:StressTensor2D}
    \begin{split}
        T(z) & := T_{zz} = T^{\bar{z}\bar{z}} \\
        \bar{T}(\bar{z}) & := T_{\bar{z}\bar{z}} = T^{zz}.
    \end{split}
\ee 
With this notation in mind we shall also introduce the notation 
\be 
\label{eqn:PartialsNotation}
    \p := \p_z \qand \bar{\p} := \p_{\bar{z}}.
\ee 
In this notation we have 
\mybox{
    \be 
    \label{eqn:PartialsStressTensor2D}
        \bar{\p} T(z) = 0 = \p \bar{T}(\bar{z}).
    \ee 
}

\section{Radial Quantisation}

Next on the list: radial quantisation. Again this is just the same as before, but our map is really from a cylinder, $\R\times S^1$, to a plane, $\C$. Our radius is now just given by $r=|z|$, and we parameterise the other coordinate using the angle $\arg(z) = \theta$. In other words we replace our $(r,\underline{n}^{\mu})$ from the $D>2$ case with $(r,\theta)$. Of course this is just a standard coordinatisation of the complex plane with $z=re^{i\theta}$ and so we use $(z,\bar{z})$ as coordinates, as we have been up to now. 

\subsection{Hermitian Conjugation In 2D Radial Quantisation}

We now get to something new coming from the fact that we are considering a 2D CFT. Recall that the Hermitian conjugation was given by
\bse 
    \cO_{\text{flat}}(r,\underline{n})^{\dagger} = \frac{1}{r^{2\Delta}} \cO_{\text{flat}}\bigg(\frac{1}{r},\underline{n}\bigg),
\ese
If we therefore consider a scalar field --- i.e. no spin, so $h=\bar{h}$ --- we have $\Delta=2h$, and so our complex conjugation becomes 
\bse 
    \cO(z,\bar{z})^{\dagger} = \frac{1}{(z\bar{z})^{2h}} \cO\bigg(\frac{1}{\bar{z}}, \frac{1}{z}\bigg)
\ese 
where we have used $r^2=|z|^2 = z\bar{z}$ along with
\bse 
    \frac{1}{\bar{z}} = \frac{1}{re^{-i\theta}} = \frac{1}{r}e^{i\theta}.
\ese 
The natural generlisation for fields with spin is just
\be 
\label{eqn:HermitianConjugatePhi2D}
    \Phi(z,\bar{z})^{\dagger} = \frac{1}{z^{2\bar{h}}} \frac{1}{\bar{z}^{2h}} \Phi\bigg(\frac{1}{\bar{z}}, \frac{1}{z}\bigg).
\ee 

Ok so why is this different to the higher dimensional case? Well in our $\C$ plane we can take a Laurent expansions for our fields as
\bse 
    \Phi(z,\bar{z}) = \sum_{m,n\in \Z} \frac{\varphi_{nm}}{z^n \bar{z}^m}.
\ese
As will become clear soon, it is actually convenient to shift\footnote{This is a shift in the sense of we are just shifting where $m\to m^{\prime} = m+h$ and $n\to n^{\prime} = n+\bar{h}$. As the sum is over all integers, we can take the sum over $m,n$ again instead of $m^{\prime}, n^{\prime}$ without ruining anything.} this Laurent expansion as 
\bse 
    \Phi(z,\bar{z}) = \sum_{m,n\in \Z} \frac{\varphi_{nm}}{z^{n+h} \bar{z}^{m+\bar{h}}}.
\ese

Now the right-hand side of \Cref{eqn:HermitianConjugatePhi2D} is then
\bse 
    \Phi^{\dagger}(z,\bar{z}) = \frac{1}{z^{2\bar{h}}} \frac{1}{\bar{z}^{2h}} \sum_{n,m} \varphi_{n,m} \bar{z}^{n+h} z^{m+\bar{h}} = \sum_{m,n} \varphi_{nm} z^{m-\bar{h}}\bar{z}^{n-h},
\ese 
but the left-hand side is just 
\bse 
    \sum_{n,m} \frac{\varphi_{nm}^{\dagger}}{\bar{z}^{n+h} z^{m+\bar{h}} } = \sum_{m,n} \varphi^{\dagger}_{-n,-m} \bar{z}^{n-h} z^{m-\bar{h}}, 
\ese 
where the equality follows from the relabelling $m\to -m$ and $n\to -n$. So comparing these two see see 
\mybox{
    \be 
    \label{eqn:VarphiNMDagger}
        \varphi^{\dagger}_{n,m} = \varphi_{-n,-m}.
    \ee 
}
\noindent The reason we did the shift of the sum above was so that this result turned out nicely. In other words, if we hadn't do the shift we would have factors of $h/\bar{h}$ appearing in \Cref{eqn:VarphiNMDagger}.

\section{Operator State Correspondance}

We now move on to discuss the state-operator correspondance. Recall that the state-operator map allowed us to define the state $\ket{\Phi}$ as
\bse 
    \ket{\Phi} := \Phi(0,0)\ket{0}.
\ese 
If we then use the mode expansion above,
\bse 
    \ket{\Phi} = \sum \frac{\varphi_{nm}}{z^{n+h}\bar{z}^{m+\bar{h}}}\bigg|_{z,\bar{z}=0} \ket{0},
\ese
we see that we need 
\bse 
    \varphi_{nm}\ket{0} = 0 \qquad \forall (n+h) > 0 \qand (m+\bar{h}) >0,
\ese
and 
\bse 
    \ket{\Phi} = \varphi_{-h,-\bar{h}}\ket{0}.
\ese 
We then define the $\bar{\Phi}$ state by Hermitian conjugation
\bse 
    \bra{\Phi} := \bra{0}\Phi^{\dagger}(0,0) = \bra{0} \big(\varphi_{-h,-\bar{h}}\big)^{\dagger} = \bra{0}\varphi_{h,\bar{h}}.
\ese 

\section{OPE \& Stress-Tensor}

Finally we move on to the OPE. Of course this works in exactly the same way as before, and we have
\bse 
    \Phi_1(z) \Phi_2(0) = \sum_{\text{primaries}} C_{123} C(z,\p_y) \Phi_3(y)\big|_{y=0}.
\ese 
As above, we now consider putting in the mode expansions. In order to simplify things, we shall just consider a Chiral field (i.e one that just depends on $z$), and obtain the general field relation by simply adding the barred bit. Our mode expansion is then just
\bse 
    \Phi(z) = \sum \frac{\varphi_n}{z^{n+h}}.
\ese 

Now, for a reason that will become clear in a moment, let's consider what happens when we take a contour integral around our local insertion $\Phi(z)$. We will weight this contour integral by $z^{n+h-1}$, giving us
\bse
    \frac{1}{2\pi i}\oint \Phi(z) z^{n+h-1} dz = \frac{1}{2\pi i} \oint \sum_{n^{\prime}} \frac{\varphi_{n^{\prime}}}{z^{n^{\prime}+h}} z^{n+h-1} dz = \frac{1}{2\pi i} \oint \sum_{n^{\prime}} \frac{\varphi_{n^{\prime}}}{z^{n^{\prime}-n}} \frac{1}{z} dz = \varphi_n
\ese 
where the integral is taken over a contour enclosing $z$, and where  the last line follows from the residue theorem (i.e. we get $\del_{n,n^{\prime}}$ and then use the sum). We therefore see that we can extract the modes via a contour integral. We shall return to this shortly.

We now want to go back to talking about Noether currents and their associated charge. Recall in higher dimensions we had a conformal Killing vector with a corresponding Noether current, which in turn has a corresponding conserved charge\footnote{Written here in time quantisation. This is why we have a $0$ index and are integrating over space.}
\bse 
    \xi^{\nu}\p_{\nu} \longrightarrow J^{\mu} = T^{\mu\nu}\xi_{\nu} \longrightarrow Q_{\xi} = \int T^{0\nu} \xi_{\nu} d^3 x.
\ese 
How does this translate to our 2D picture? Well of course it is all the same, but the idea is to try use our complex coordinates to simplify stuff. 

In the complex picture we have two conserved currents given by the transformations $z\to z + \xi(z) $ and $\Bar{z} \to \Bar{z}+\Bar{\xi}(\bar{z})$. The currents are simply 
\bse 
    J(z) = \xi(z)T(z) \qand \bar{J}(z) = \bar{\xi}(\bar{z})\bar{T}(\bar{z}).
\ese
Our charges are then simply given by
\bse 
    Q_{\xi} := \frac{1}{2\pi i} \oint \xi(z) T(z) dz, \qand Q_{\bar{\xi}} := \frac{1}{2\pi i} \oint \bar{\xi}(\bar{z}) \bar{T}(\bar{z}) d\bar{z}
\ese
where again our contours are done over the region containing $z/\bar{z}$. The reason for this form is because here we are working in radial quantisation and so our equal time slices (i.e. integrating over $d^3x$) become equal radial slices (i.e. closed contour integrals).

Next also recall that in QFT charges become the generators of the symmetry via commutators. That is, for a Chiral primary of weight $h$,\footnote{Note we don't need to say $(h,0)$ as Chiral already tells us $\bar{h}=0$ as there is no $\bar{z}$ dependence.}
\be 
\label{eqn:QPhiCommutator}
    [Q_{\xi},\Phi(z)] = - \del_{\xi} \Phi = h (\p \xi) \Phi + \xi \p \Phi,
\ee 
where the last line comes from \Cref{eqn:InfintesimalPrimaryTransformation2D} (where we have ignored the barred terms). 

This is our CFT result, but we also know that the commutator is just given by
\bse 
    [Q_{\xi}, \Phi(w)] = Q_{\xi}\Phi(w) - \Phi(w) Q_{\xi} 
\ese 
The right-hand side is an operator, and so it is meant to be understood inside a correlator. In `normal' QFT this corresponds to time ordering, however here we are working in radial quantisation, and so time ordering is replaced by radial ordering. The first term in the above result then corresponds to the $Q_{\xi}$ contour enclosing $\Phi(w)$ (i.e. $r_{\xi} > r_w$) and the second term corresponds to the contour not enclosing $\Phi(w)$ (i.e. $r_{\xi} < r_w$). Pictorially,
\begin{center}
    \btik 
        \begin{scope}[xshift=-3cm]
            \draw[thick] (-2,-2) -- (2,-2) -- (2,2) -- (-2,2) -- (-2,-2);
            \draw (-2,0) -- (2,0);
            \draw (0,-2) -- (0,2);
            \node at (1,1) {$\cross$};
            \node[below] at (1,1) {$\Phi$};
            \draw[thick, blue, decoration={markings, mark=at position 0.15 with {\arrow{>}}}, postaction={decorate}] (0,0) circle [radius=1.75cm];
        \end{scope}
        \node at (0,0) {\Huge{$-$}};
        \begin{scope}[xshift=3cm]
            \draw[thick] (-2,-2) -- (2,-2) -- (2,2) -- (-2,2) -- (-2,-2);
            \draw (-2,0) -- (2,0);
            \draw (0,-2) -- (0,2);
            \node at (1,1) {$\cross$};
            \node[below] at (1,1) {$\Phi$};
            \draw[thick, blue, decoration={markings, mark=at position 0.15 with {\arrow{>}}}, postaction={decorate}] (0,0) circle [radius=0.75cm];
        \end{scope}
        \node at (6,0) {\Huge{$=$}};
        \begin{scope}[xshift=9cm]
            \draw[thick] (-2,-2) -- (2,-2) -- (2,2) -- (-2,2) -- (-2,-2);
            \draw (-2,0) -- (2,0);
            \draw (0,-2) -- (0,2);
            \node at (1,1) {$\cross$};
            \node[below] at (1,1) {$\Phi$};
            \draw[thick, blue, decoration={markings, mark=at position 0.15 with {\arrow{<}}}, postaction={decorate}] (0,0) circle [radius=0.75cm];
            \draw[thick, blue, decoration={markings, mark=at position 0.15 with {\arrow{>}}}, postaction={decorate}] (0,0) circle [radius=1.75cm];
        \end{scope}
    \etik 
\end{center}
where we note that the arrow of the inner circle switches in the last diagram. This is just the absorption of the minus sign to switch the contour direction. 

We now use some complex analysis: we use the fact that our charges are holomorphic to deform our contours. Assume that there are no poles apart from $z$ and $w$ (this is equivalent to saying there are no other fields near by), we then just get 
\begin{center}
    \btik 
        \begin{scope}[xshift=-3cm]
            \draw[thick] (-2,-2) -- (2,-2) -- (2,2) -- (-2,2) -- (-2,-2);
            \draw (-2,0) -- (2,0);
            \draw (0,-2) -- (0,2);
            \node at (1,1) {$\cross$};
            \draw[thick, blue, decoration={markings, mark=at position 0.15 with {\arrow{<}}}, postaction={decorate}] (0,0) circle [radius=0.75cm];
            \draw[thick, blue, decoration={markings, mark=at position 0.15 with {\arrow{>}}}, postaction={decorate}] (0,0) circle [radius=1.75cm];
        \end{scope}
        \node at (0,0) {\Huge{$=$}};
        \begin{scope}[xshift=3cm]
            \draw[thick] (-2,-2) -- (2,-2) -- (2,2) -- (-2,2) -- (-2,-2);
            \draw (-2,0) -- (2,0);
            \draw (0,-2) -- (0,2);
            \node at (1,1) {$\cross$};
            \draw[thick, blue, decoration={markings, mark=at position 0.2 with {\arrow{>}}}, postaction={decorate}] (1,1) circle [radius=0.25cm];
        \end{scope}
    \etik  
\end{center}
This can be seen by either "pinching" the two contours together and seeing that they cancel apart from around the insertion, or by considering first contracting the inner circle to nothing and then contracting the outer circle around the insertion. 

The conclusion is that the commutator 
\bse 
    [Q_{\xi},\Phi] = \frac{1}{2\pi i} \oint \xi(z) T(z) \Phi(w) dz,
\ese 
where the contour is taken around $w$. Now this is supposed to be equal to \Cref{eqn:QPhiCommutator}. This tells us something about the $T(z)\Phi(w)$ OPE. This might not be immediately clear so let's air out any confusion. If we are going to get the terms on the right-hand side of \Cref{eqn:QPhiCommutator} we need our contour integral to pick up two poles. The easier one to see is the $\xi(w)\p_w \Phi(w)$:\footnote{Note it is $w$ not $z$, as $z$ is the integration variable.} if the OPE contained the term 
\bse 
    T(z)\Phi(w) = ... + \frac{1}{z-w} \p_w \Phi(w) + ...
\ese 
then we would just pick up this pole and set $z=w$ giving us exactly what we want. The other term takes a bit more work. We see that there is a derivative acting on the $\xi$, this suggests we want to some kind of integration by parts within our integral. Then noting 
\bse 
    \p_z\bigg(\frac{1}{z-w}\bigg) = -\frac{1}{(z-w)^2},
\ese 
we see that the term we need is 
\bse 
    T(z)\Phi(w) = ... + \frac{h}{(z-w)^2} \Phi(w) + ...
\ese 
In other words we have
\bse 
    \begin{split}
        \frac{1}{2\pi i}\oint \xi(z) \p_z\bigg(-\frac{1}{z-w}\bigg) h\Phi(w) dz & = \frac{1}{2\pi i}\oint \big(\p_z\xi(z)\big) \frac{1}{z-w} h\Phi(w) dz \\
        & = \big(\p_z\xi(z)\big)h\Phi(w)\big|_{z=w} \\
        & = h\big(\p_w\xi(w)\big)\Phi(w),
    \end{split}
\ese
where the minus sign from the derivative is cancelled by the minus sign from the integration by parts. Finally, by extension of this idea, it's hopefully clear that we don't want any higher poles. That is we don't want any $\frac{1}{(z-w)^3}$ poles etc. We therefore conclude that the OPE between a Chiral primary of weight $h$ and the stress-energy tensor is
\bse 
    T(z)\Phi(w) = h\frac{\Phi(w)}{(z-w)^2} + \frac{\p_w \Phi(w)}{z-w} + \text{non-singular}
\ese
We can then extend this argument to the OPE of a general primary operator (i.e. depending on both $z$ and $\bar{z}$) with $T/\bar{T}$ as follows.\footnote{We drop the subscripts on the derivatives and assume that the variable is understood by its action.}
\mybox{
    \be 
    \label{eqn:TPhiOPE}
        \begin{split}
            T(z)\Phi(w) & = h\frac{\Phi(w,\bar{w})}{(z-w)^2} + \frac{\p \Phi(w,\bar{w})}{z-w} + \text{non-singular} \\
            \bar{T}(\bar{z})\Phi(w) & = h\frac{\Phi(w,\bar{w})}{(\bar{z}-\bar{w})^2} + \frac{\bar{\p} \Phi(w,\bar{w})}{\bar{z}-\bar{w}} + \text{non-singular}
        \end{split}
    \ee 
}

What we have essentially shown is that the OPE of a primary operator and the stress-energy tensor is equivalent to the transformation of such an operator. In fact some people\footnote{For example: Polchinski, Prof. Tong and Prof Minwalla all do this.} give \Cref{eqn:TPhiOPE} as the \textit{definition} of a primary operator. You can then work back from this definition and obtain our definition of a primary operator in terms of its transformation property. In fact we could have done something similar in higher dimensions but we were more focused on getting toward conformal bootstrap and so just wanted the Ward identities there. 

So far everything is written with $\xi(z)$, but what about the generators, $L_n$? Well recall that we got them by expanding $\xi(z)\p_z$ as a Laurent series, giving us \Cref{eqn:Ln}. For the quantum theory we do a similar thing and expand $\xi(z)$ in $Q_{\xi}$ to give us 
\mybox{
    \be 
    \label{eqn:LnOint}
        L_n = \frac{1}{2\pi i} \oint dz \, z^{n+1} T(z).
    \ee
}
We now do something that we \textit{can't} do in the higher dimensional case by recalling the idea that we can obtain the modes by doing a contour integral weighted by $z^{n+h-1}$ to see that we can reconstruct the stress-energy tensor from the modes
\mybox{
    \be 
    \label{eqn:TLaurent}
        T(z) = \sum_n \frac{L_n}{z^{n+2}}.
    \ee
}
\noindent The $L_n$ are so-called \textit{Virasoro modes}. We obtain a version of the $T(z)\Phi(w)$ OPE on the modes,
\be 
\label{eqn:LmVarphinCommutator}
    [L_m, \varphi_n] = \big((h-1)m-n\big)\varphi_{m+n}
\ee 

\bbox 
    Use a contour integral argument to show \Cref{eqn:LmVarphinCommutator}.
\ebox 

\chapter{Example Of 2D CFT: Free Boson}

Let's now study an actual example of a $2D$ CFT. We will study the easiest case, that of a free scalar field. Even though it is the simplest model, it is a very important thing to study as string theory is essentially the study of $D$ free scalar fields in a 2D CFT. As a CFT is, in particular, a QFT essentially what we're doing is just the quantisation of a free scalar field. The only unusual feature compared to what we know from `normal' QFTs is that here we put our theory on a cylinder rather then some flat Minkowski spacetime. The idea of us putting the QFT on a cylinder corresponds to studying the closed string in string theory. 

\section{Mode Expansion}

Our action is simply the standard massless free action
\bse 
    S = -\frac{1}{4\pi \a^{\prime}} \int d^2 x \, \p_a X \p^a X,
\ese
where $x=(t,\sig)$ and $a=0,1$. Again to be clear, the $X$ here is just a scalar field from the CFT's perspective. That is we could replace it by the symbol $\phi$ an no confusion would arise. We stick to the $X$ notation as it is allows the relation to string theory to be made easier. 

We know that the equations of motion for such an action is just a 1D wave equation
\bse 
    -\p_t^2 X + \p_{\sig}^2 X = 0.
\ese 
Now comes the fact that we are working on the cylinder, so we impose periodic boundary conditions
\be 
\label{eqn:XPeriodicBoundaryConditions}
    X(t,\sig) = X(t,\sig + 2\pi).
\ee 
The general solution to such a problem is
\be 
\label{eqn:XGeneralSolution}
    X = x_0 - \sqrt{2\a^{\prime}} \a_0 t - i \sqrt{\frac{\a^{\prime}}{2}} \sum_{n\neq 0} \bigg( \frac{\a_n}{n} e^{-in(t+\sig)} + \frac{\bar{\a}_n}{n} e^{-in(t-\sig)}\bigg),
\ee 
where $x_0\in \R$\footnote{As we consider a real scalar field.} is just some constant, and where the funny $\a^{\prime}$ factors are included to make the commutators that come later nicer. Now we have just pulled this solution out of nowhere so let's check it makes sense 
\begin{itemize}
    \item The $x_0$ constant term is obviously a solution as its derivatives vanish. 
    \item A linear term in $t$ again vanishes because we have $\p_t^2$ in our equations of motion. 
    \item Taking the two derivatives on the exponential terms will cancel as there is a relative minus sign in our equations of motion. These terms obviously obey the periodic boundary condition by $e^{in(x+2\pi)} = e^{inx}$.
    \item The only other term we could imagine is something linear in $\sig$. However this is \textit{not} allowed as it wouldn't obey our periodic boundary condition, \Cref{eqn:XPeriodicBoundaryConditions}.
\end{itemize}

Now that we're happy this is the most general solution, let's make some comments/conclusions. As $X$ is a real field, we require $X^*=X$. It follows from this that 
\be 
\label{eqn:AlphaComplexConjugate}
    \a^*_n = \a_{-n} \qand \bar{\a}^*_n = \bar{\a}_{-n}. 
\ee 
It is important to note that the bar here does \textit{not} indicate complex conjugation. In order to avoid confusion we could replace it was a tilde, $\bar{\a}_n \to \widetilde{\a}_n$, however hopefully now that this has been pointed out we can stick with a bar without any confusion. $\a_n$ and $\bar{\a}_n$ play a very important role in string theory and are called \textit{oscillator modes}, as they correspond to the exponential oscillating part of \Cref{eqn:XGeneralSolution}. Next we define  
\be 
\label{eqn:Palpha0}
    p := \sqrt{\frac{2}{\a^{\prime}}} \a_0,
\ee
which at this point is just a label, however if we consider our light-cone coordinates we can see that it is actually related to the conjugate momenta $\pi$,\footnote{See Remark 2.1.1. of my notes on Prof. Minwalla's string theory course.} and so we conclude that $p$ is the centre of mass momentum of the string.

This was done in terms of $(t,\sig)$, but we now want to proceed through as before to obtain the result in terms of $(z,\bar{z})$. There we,
\ben[label=(\roman*)]
    \item Wick rotate $\tau=it$, and
    \item Define $r=e^{\tau}$ so that $e^{-i(t+\sig)}$ becomes $e^{-\tau}e^{-i\sig} = 1/(re^{i\sig}) = 1/z$. 
\een 
so that
\bse 
    X = x_0 + i\a^{\prime} p\ln|z| - i\sqrt{\frac{\a^{\prime}}{2}} \sum_{n\neq0} \bigg( \frac{\a_n}{n} z^{-n} + \frac{\bar{\a}_n}{n} \bar{z}^{-n}\bigg).
\ese 
This object is not a nice object from a CFT point of view mainly because of the $\ln|z|$ term, which is not meromorphic. However we now notice that if we take a derivative, $\p$ or $\bar{\p}$, we will get a much nicer expression
\bse 
    \p X = i \sqrt{\frac{\a^{\prime}}{2}} \sum_{n\neq 0} \frac{\a_n}{z^{n+1}} \qand \bar{\p} X = -i \sqrt{\frac{\a^{\prime}}{2}} \sum_{n\neq 0} \frac{\bar{\a}_n}{\bar{z}^{n+1}}.
\ese 
We can show\footnote{See, for example, section 4.3.3 of Prof Tong's String Theory notes.} that these are primary fields of weights $(1,0)$ and $(0,1)$, respectively. However the field $X$ itself is \textit{not} a primary field, and so we much prefer to work with $\p X$ and $\bar{\p}X$.

\br 
    We can get a rough feeling for the weights above being correct by considering the dimension of $X$. From our action we see $[X]=0$, and so from $[\p]=[\bar{\p}]=1$, we expect dilatation weights $\Delta=1$ for $\p X$ and $\bar{\p}X$. We then make arguments about $z/\bar{z}$ independence to get $(1,0)$ and $(0,1)$.
\er 

This has all been a classical discussion, let's not move on and quantise it.

\section{Quantisation of the 2d Free Scalar}

For the quantisation procedure we will go back to the $(t,\sig)$ coordinates, as we are more familiar with quantising things in this way. We know from our canonical QFT courses that the commutation relations for a free Bosonic field are given by the field and the conjugate momenta. We take these to be equal time commutation relations. So let's compute the conjugate momenta
\bse 
    \pi = \frac{\p \cL}{\p \dot{X}} = \frac{1}{2\pi \a^{\prime}} \dot{X}.
\ese
Putting this into our expected equal time commutation relations we get
\bse 
    [X(t,\sig), \dot{X}(t, \sig^{\prime})] = 2\pi \a^{\prime} i \del(\sig-\sig^{\prime}),
\ese 
where the $2\pi\a^{\prime}$ factor on the right comes from us putting $\dot{X}$ in the commutator not $\pi$. We can then plug our mode expansion for $X$ in to obtain commutation relations between the modes themselves. The results are
\mybox{
    \be 
    \label{eqn:ModeCommutations}
        [x_0,p] = i \qquad [\a_m,\a_n] = m\del_{m+n,0} \qand [\bar{\a}_m,\bar{\a}_n] = m\del_{m+n,0}
    \ee 
}
\noindent We see that all the funny $\a^{\prime}$ factors have disappeared in our commutation relations. This is why we introduced them in the first place.

\badr 
    Note that the $[x_0,p]=i$ commutation relation further supports the idea that $p$ is the centre of mass momentum and that $x_0$ labels the centre of mass. That is it agrees with the usual $[q,p]=i$ relation. 
\eadr 

\bbox 
    Prove \Cref{eqn:ModeCommutations}. \textit{Hint: Use}
    \bse 
        \sum_n e^{-2in(\sig-\sig^{\prime})} = n\del(\sig-\sig^{\prime}).
    \ese 
\ebox 

Next we want to compute the Hamiltonian. The standard QFT derivation gives us 
\be 
\label{eqn:HamiltonianBosonicField}
    H = \int d\sig \, \big( \pi \dot{X} - \cL\big) = \frac{\a^{\prime}p^2}{2} + \sum_{n>0} \big(\a_{-n}\a_n + \bar{\a}_{-n}\bar{\a}_n\big)  
\ee 

\bbox 
    Fill in the missing steps to arrive at \Cref{eqn:HamiltonianBosonicField}.
\ebox 

From here we can compute the commutator
\bse 
    [H,\a_{-n}] = n \a_{-n} \qand [H,\bar{\a}_{-n}] = n \bar{\a}_{-n}
\ese 
which is easily verified. Now we note that these look just like creation/annihilation operator commutators with the energy spacing being given by $n$. That is $\a_n/\bar{\a}_n$ are raising operators when $n<0$ and lowering operators when to $n>0$. It is important to note that we have \textit{two}, independent, sets of creations/annihilation operators. Therefore our Hilbert space is built out of \textit{two} towers of states given by acting with the creation operators. It turns out that there is a relation, known as the \textit{level matching condition}, that tells us the number of $\a_n$ excitations must equal the number of $\bar{\a}_n$ excitations, and so we only need one label in our bra-ket vector. As we will see shortly our Hilbert space also has another sector given by the momentum. So in total we denote an element of the Hilbert space by $\ket{n;k}$, where $n$ indicates the number of harmonic oscillator states and $k$ the momentum. As usual, we construct our Hilbert space by defining the vacuum $\ket{0;k}$ such that 
\bse 
    \a_m \ket{0;k} = 0 = \bar{\a}_m \ket{0,k} \qquad \forall m >0.
\ese 

This seems a bit strange, because it appears we don't have a unique vacuum, i.e. we haven't constrained $k$. So what's going on? Well now is a point where the fact that our CFT is on a cylinder becomes important.  Normally we think of the fields as dying off at infinity, however on the cylinder the fields can go round and round. This corresponds exactly to the ground state not being unique, and the different ground states being labelled by the momentum. So how we do extract the momentum? Well in the above we only considered the $\a_n/\bar{\a}_n$s for $n\neq 0$, but what about $\a_0$? Well we have already seen/said that it is related to the centre of mass momentum. We can therefore use this to extract $k$:
\bse 
    p\ket{0;k} = k\ket{0;k}.
\ese 

Now that we have defined our vacuua, we can build up the Hilbert space by acting with the raising operators. That is any state in the Hilbert space is given in the form
\bse 
    (\a_{-1})^{n_1} (\a_{-2})^{n_2} ... (\bar{\a}_{-1})^{\bar{n}_1} (\bar{\a}_{-2})^{\bar{n}_2}\ket{0;k}.
\ese 

There is an obvious question to ask "what about when $k=0$?" This is indeed an important case, and we define $\ket{0}:= \ket{0;0}$ which is often called the \textit{absolute} vacuum. As we will see shortly, this is invariant under the whole global conformal subgroup, i.e. the M\"{o}bius subgroup. We can then construct $\ket{0;k}$ from $\ket{0}$ via the following claim (we are going back to $(z,\bar{z})$ coordinates\footnote{We only really used the $(t,\sig)$ coordinates to do the quantisation because we were familiar with how to quantise in equal time slices. That's all done now, so it makes sense to go back to the more useful complex coordinates.}).
\bcl 
    We can produce the vacuum state $\ket{0;k}$ from the absolute vacuum in the following way.
    \bse 
        \ket{0;k} = e^{ikX(z=0)}\ket{0}.
    \ese
\ecl 

Before proving this, let's just make a couple comments. This is an important thing in string theory and it is called a \textit{vertex operator}. Note that we \textit{cannot} write this down in higher dimensions, as then $[X]\neq 0$ so its dilatation weight is non-vanishing. That is, if we Taylor expanded the exponential we wouldn't have a well-defined dilatation transformation property. 
\bq 
    Essentially what we need to do is prove it has momentum $k$, and that it is annihilated by the lowering operators $\a_n/\bar{\a}_n$ for $n >0$.  
    \bse 
        \begin{split}
            p \ket{0;k} & = p e^{ikX(0)}\ket{0} \\
            & = \Big[p,e^{ikX(0)}\Big]\ket{0} \\
            & = e^{ikX(0)}[p,ikX(0)]\ket{0} \\
            & = ik e^{ikX(0)} [p,x_0]\ket{0} \\
            & = k e^{ikX(0)} \ket{0}\\
            & = k\ket{0;k},
        \end{split}
    \ese 
    where we have used the mode expansion for $X$ and then only kept the commutators that are non-vanishing. Similarly, we have, for $n>0$, 
    \bse 
        \begin{split}
            \a_n\ket{0;k} & = \Big[\a_n, e^{ikX(0)}\Big] \ket{0} \\
            & = e^{ikX(0)} [\a_n, ik X(0)] \ket{0} \\
            & = e^{ikX(0)} ik \sum_{m} \bigg[\a_n, \frac{\a_{m}}{m} z^{-m}\bigg]\bigg|_{z=0} \ket{0} \\
            & = e^{ikX(0)} \frac{ikn}{m}  \sum_m \del_{n+m,0} z^{-m}\big|_{z=0} \ket{0} \\
            & =  e^{ikX(0)} ik z^n\big|_{z=0} \ket{0}\\
            & = 0,
        \end{split}
    \ese 
    where again we have used the mode expansion and only kept the non-vanishing commutators. A similar calculation works for the $\bar{\a}_n$ lowering operators.
\eq 

Now that we know how to construct our Hilbert space we can use our state-operator correspondance to see what operators these states are dual to. We just give a couple examples here.

\begin{center}
	\begin{tabular}{@{} C{5cm} C{5cm} @{}}
		\toprule
		State & Operator  \\
		\midrule 
		$\ket{0;0}$ & $\b1$ \\
		$\a_{-m}\ket{0;0}$ & $\p^m X$ \\
		$\a_{-m^{\prime}}\a_{-m}\ket{0;0}$ & $\cl \p^{m^{\prime}}X\p^m X \cl$ \\
		$\ket{0;k}$ & $\cl e^{ikX} \cl $ \\
		\bottomrule
	\end{tabular}
\end{center}
where $\cl ... \cl$ means normal ordering.\footnote{It is actually something known as \textit{conformal normal ordering}. This is defined as $\cl A B \cl := \lim_{z\to w} [A(z)B(w)- \la A(z) B(w)\ra ]$. The idea is that a conformal normal ordering gives vanishing vev, $\la \cl A(z)B(w) \cl \ra$. In these notes we will ignore this small difference and just say normal ordering, but feel free to insert "conformal" in front when reading.}

\badr 
    It turns out that the vertex operator $\cl e^{ikX} \cl$ is a primary operator with weight 
    \bse 
        h = \frac{\a^{\prime} k^2}{4}.
    \ese 
    You can show this by considering the OPE with the stress-energy tensor (which we will obtain in a moment). Details of this calculation can be found in section 8.6 of my notes on Prof. Minwalla's string theory course. 
\eadr 

\subsection{Propagator}

Recall that in QFT the propagator is one of the most important objects, and it is given by the two-point function
\bse 
    \la X(z) X(w)\ra = k(z,w).
\ese 
Also recall that we require the propagator to be a Green's function for the quadratic operator, i.e.
\bse 
    \p \bar{\p} k(z,w) = -\del^{(2)}(z-w).
\ese 
The solution to this is a log expression:
\bse 
    \la X(z) X(w)\ra = -\frac{1}{2\pi} \log |z-w|^2,
\ese
where 
\bse 
    |z-w|^2 = (z-w)(\bar{z}-\bar{w}).
\ese 

This tells us that $X$ is not a primary operator, as we know that two point functions of scalar primary operators go as $1/(z-w)^{2\Delta}$, but we have a log. However we note that if we take two derivatives we get 
\bse 
    \la \p X(z) \p X(w) \ra = -\frac{1}{2\pi} \frac{1}{(z-w)^2},
\ese 
where the variable of the derivative is understood by the argument of the field it acts on. Note also we don't have a modulus $|z-w|$ in the denominator. This is because we have only differentiated w.r.t. $z/w$ (i.e. not barred) and 
\bse 
    \log|z-w|^2 = \log\big[(z-w)(\bar{z}-\bar{w})\big] = \log(z-w) + \log(\bar{z}-\bar{w}).
\ese 

We therefore see that $\p X(z)$ \textit{is} a scalar primary operator with weight $h=1$. By our "everything follows by sticking a bar on it" mantra, we also have 
\bse 
    \la \bar{\p} X(\bar{z}) \bar{\p} X(\bar{w}) \ra = -\frac{1}{2\pi} \frac{1}{(\bar{z}-\bar{w})^2}.
\ese 
We summarise this result in a nice box. 
\mybox{
    \begin{center}
        $\p X(z)$ is a scalar primary operator of weight $(1,0)$. \\
        $\bar{\p}X(\bar{z})$ is a scalar primary of weight $(0,1)$.
    \end{center}
}

\subsection{Stress-Energy Tensor}

Ok so what exactly is the stress-energy tensor for our free boson theory? Well classically we can show\footnote{See Section 4.1.3 of Prof. Tong's notes.} that it just comes out to be
\bse 
    T = - \frac{1}{\a^{\prime}} \p X \p X, \qand \bar{T} = - \frac{1}{\a^{\prime}} \bar{\p} X \bar{\p} X.
\ese
When we quantise this result, we have to include our normal ordering\footnote{This is one of the big motivations for defining conformal normal ordering the way we do. Essentially we want the vev of $T$ to vanish, and so we construct it such that this is the case.}
\mybox{
    \be
    \label{eqn:TBoson}
        T = -\frac{1}{\a^{\prime}} \cl \p X \p X \cl.
    \ee 
}

Now recall that we wrote the stress energy tensor as a Laurent expansion 
\bse 
    T = \sum_n \frac{L_n}{z^{n+2}},
\ese 
we can then insert our mode expansions for $X$ into \Cref{eqn:TBoson} and obtain 
\bse 
    T = \frac{1}{2} \cl \sum_m \frac{\a_m}{z^{m+1}} \sum_n \frac{\a_n}{z^{n+1}} \cl ,
\ese 
If we then pick off the $z^{-(m+2)}$ part of this to get the Virasoro modes 
\bse 
    L_m = \cl \frac{1}{2}\sum_n \a_{m+n}\a_{-n} \cl = \frac{1}{2}\sum_n \a_{m+n}\a_{-n} \qquad m \neq 0.
\ese 
where the second line follows from the fact that when $m\neq0$ we always get a vanishing commutator between $\a_{m+n}\a_{-n}$ so no problems arise from our normal ordering. 

What about when $m=0$? Well there we have 
\bse 
    L_0 = \frac{1}{2} \cl \sum_n \a_n \a_{-n} \cl = \frac{1}{2}\a_0^2 + \sum_{n>0} \a_{-n}\a_n,
\ese 
where we have split the sum into $n>0$ and $n<0$ and then relabelled to get rid of the 1/2. We can then finally rewrite this as 
\bse 
    L_0 = \frac{\a^{\prime}}{4}p^2 + \sum_{n>0} \a_{-n}\a_n.
\ese 

Of course again we get a similar result by putting bars everywhere. 

\br 
    In string theory $L_0=0$ gives us a condition on the masses of our particles. This is known as the \textit{level matching} condition and is precisely the condition we used above to say that we always require the number of $\a$ and $\bar{\a}$ excitations match. 
\er 

\bbox 
    Using the relations above prove that 
    \be 
    \label{eqn:LalphaCommutator}
        [L_m, \a_n] = -n\a_{m+n}, \qand [\bar{L}_m, \bar{\a}_n] = -n\bar{\a}_{m+n}
    \ee 
\ebox 

Comparing this with the transformation of the modes of a primary,
\bse 
    [L_m,\varphi_n] = \big((h-1)m-n\big) \varphi_{m+n},
\ese 
we see this is consistent with $\a_n$ being the modes of $\p X$, which is a primary of weight $(1,0)$. That is putting $h=1$ in the above expression will give exactly \Cref{eqn:LalphaCommutator}.

Now we come to the interesting bit, the commutator of $L_m$ with $L_n$. Plugging through the calculations we obtain
\bse 
    [L_m,L_n] = (m-n)L_{m+n} + \frac{1}{12} m(m^2-1) \del_{m+n,0},
\ese
where we recognise the first term from the Witt algebra, but the second term is new. This extra term is called a \textit{central extension} of the algebra,\footnote{This name comes from the fact that the centre of an algebra is something that commutes with every element, which this last term does as it's just a number.} i.e. it breaks our closure of the algebra.

\br 
    Note that the extra term vanishes for our subalgebra, i.e. $\{L_0,L_1,L_{-1}\}$. From this we see that we never see an equivalent in higher dimensions. 
\er 

Let's derive this. We will cheat slightly and then explain why we cheated at the end. We will focus on the case when $m=-n$ so that we can get this extra piece as this is the harder case.\footnote{We can also obtain these results using contour type arguments, once we know the $T(z)T(w)$ OPE. For details of this calculation see, e.g., Section 4.5.2 of Prof. Tong's notes.}

\bse 
    \begin{split}
        [L_m,L_{-m}] & = \frac{1}{4} \sum_n \sum_{n^{\prime}} [ \a_{m-n}\a_n, \a_{-m-n^{\prime}}\a_{n^{\prime}}] \\
        & = \frac{1}{4}\sum_{n,n^{\prime}} \big( \a_{m-n}[\a_n,\a_{m-n^{\prime}}]\a_{n^{\prime}} + \a_{-m-n^{\prime}}[\a_{m-n},\a_{n^{\prime}}]\a_n \\
        & \qquad \qquad  + \a_{m-n}[\a_n,\a_{n^{\prime}}]\a_{-m-n^{\prime}} + \a_{n^{\prime}} [\a_{m-n}, \a_{-m-n^{\prime}}] \a_n \big) \\
        & = \frac{1}{4} \sum_n 2 \big( n\a_{m-n}\a_{n-m} + (m-n)\a_{-n}\a_n\big),
    \end{split}
\ese 
where we have used the fact that the commutators give us bunch of delta functions. Now in order to relate it to the $L$s we need to normal order. In order to do that, notice that when $m>n$, we have 
\bse 
    \a_{m-n}\a_{n-m} = \a_{n-m}\a_{m-n} + (m-n).
\ese
Also notice that when $n<0$, we have
\bse 
    \a_{-n}\a_n = \a_n\a_{-n} + (-n).
\ese 
Now consider the case when both these inequalities are satisfied. Then the extra terms from normal ordering above cancel. That is our commutator will contain the terms
\bse 
    n(m-n) + (m-n)(-n) = 0.
\ese 
Therefore the only additional terms we pick up from normal ordering occur in the inequalities $0<n<m$. 
\begin{center}
    \btik 
        \draw[->] (-4,0) -- (4,0);
        \node at (3.8,-0.2) {$n$};
        \draw[ultra thick, red] (-2,0) -- (2,0);
        \draw[fill=black] (-2,0) circle [radius=0.05cm] node [below] {$0$} ; 
        \draw[fill=black] (2,0) circle [radius=0.05cm] node [below] {$m$};
    \etik 
\end{center}
We can then perform the sum over these extra terms and obtain 
\bse 
    \frac{2}{4}\sum_{n=0}^m n(m-n) = \frac{1}{12}m(m^2-1).
\ese 

So how did we cheat? Well we are really cancelling two infinite sums. Of course we showed they cancel term by term, but the `better' way to do is it introduce a normal ordering constant and show that this is the only thing that satisfies the Jacobi identity. 


\chapter{Quantum Stress Tensor \& Virassoro Reps}

Everything we did in the last chapter was for our specific example of a free Bosonic field. Let's go back now to a general CFT in $2D$. Not much is too different. For example the Witt algebra is simply modified to
\bse 
    [L_m,L_n] = (m-n)L_{m+n} + \frac{c}{12} m(m^2-1)\del_{m+n,0}.
\ese 
which is known as the \textit{Virasoro algebra}. Again technically this is not really an algebra but is a central extension of one, where $c$ is the central extension. $c$ is known as the \textit{central charge}. As we just saw, for the single free Boson we have $c=1$. It's not too hard to see that if we extend our theory to $D$ free Bosons, we get $c=D$. From this we can lean on intuition and arrive at the result that $c$ essentially measures the number of degrees of freedom on the system. As we will see shortly, the central charge plays a massive role in string theory and basically gives us the dimension of the spacetime. 

Ok, now recall that for primaries\footnote{In $2D$ one often only writes the singular terms and leaves the "$+$ non-singular" implicit. This is why we don't write it in this formula here.} 
\bse 
    T(z) \Phi(w) = h\frac{\Phi(w) }{(z-w)^2} + \frac{\p_w \Phi(w)}{(z-w)} \qquad \iff \qquad [L_m,\varphi_n] = \big((h-1)m-n)\big)\varphi_{m+n}.
\ese 
The Virasoro algebra tells us that $L_m$ no longer corresponds to a primary of weight $2$. That is if we  put $\varphi_n=L_n$ we get the central charge term which breaks our primary operator relation. Recalling that $L_m$ are the mode expansion of the stress-energy tensor, what we're saying is that $T$ is \textit{not} primary. By the $\iff$ above, this is equivalent to saying that the $T(z)T(w)$ OPE contains a higher pole singularity. It turns out that the result is in fact 
\mybox{
    \be 
    \label{eqn:TTOPE}
        T(z)T(w) = \frac{c/2}{(z-w)^4} + 2\frac{T(w)}{(z-w)^2} + \frac{\p T(w)}{z-w}.
    \ee
}
A derivation of this, for the case of a free Boson, can be found on page 82 of Prof. Tong's notes. However let's note why it makes sense. $T$ is still a dimension $(2,0)$ operator, it is just not primary. We see from this that it has $\Delta=2$, from which it follows that every term in the OPE must have $\Delta=2+2=4$. Now $(z-w)^{-1}$ has $\Delta=1$, and so it's possible for us to have everything up to $(z-w)^{-4}$ poles.\footnote{You might ask about higher poles and having some operator in the numerator with $\Delta_{\cO}<0$. These turn out to be excluded from unitary theories, so we drop these cases.} It's then a simple game of asking "what operators can I put in the numerator to get the weights right?", and we arrive at something of the form \Cref{eqn:TTOPE}. 

\badr 
    It's a fair question to ask "why don't we have a $(z-w)^{-3}$ term in \Cref{eqn:TTOPE}?" The answer to this involves a bit of work but essentially it comes down to the fact that we require the symmetry $T(z)T(w)=T(w)T(z)$. The cubic term would violate this. The $(z-w)^{-1}$ term \textit{appears} to also violate this, but it can be justified by a Taylor expansion argument: $T(z) = T(w) + (z-w)\p T(w) + ...$. Again for more details see Prof. Tong's notes.
\eadr 

\br 
    The fact that $T$ is not a primary is not actually a problem for our CFT, just means we're using this extension of the algebra. Note, however, that $T$ is indeed quasi-primary. 
\er 

Right, let's see how $T$ transforms under the conformal group, then. That is it doesn't transform as a primary, so how \textit{does} it transform? We do this by considering the contour arguments we made before and then use the OPE \Cref{eqn:TTOPE}. That is we compute\footnote{Again the minus sign on the right-hand side comes from the difference between the $[Q_{\xi},\cO] = -\del_{\xi}\cO$.}
\bse 
    \del_{\xi} T(w) = - \frac{1}{2\pi i} \oint dz \xi(z) T(z) T(w),
\ese
by plugging in the OPE:
\bse 
    \del_{\xi}T(w) = -\frac{c}{12} \p_{w}^3 \xi(w) - 2T(w) \p_w \xi(w) - \xi(w) \p_w T(w)
\ese 
The important thing to note is that for the global M\"{o}bius subgroup we had that $\xi(w)$ was quadratic in $w$, and so the first term in the above vanishes. So we see that for higher dimensional CFTs the stress-energy tensor is like a primary operator if weight $\Delta=D$.\footnote{The easiest way to see that we have $\Delta = D$ is to recall that we get the energy by integrating the stress-energy tensor over space. Each spatial integral measure carries weight $[dx]=-1$, and then using $[E]=1$, we conclude $[T]=(D-1)+1=D$. Then simply use that the dilatation weight and mass dimension agree.} 

\section{Virasoro Representations}

We now want to organise our Hilbert space into the irreps of the Virasoro algebra. We label our states by the weight $\ket{h}$ such that $L_0\ket{h} = h\ket{h}$. We saw before that $\widetilde{P}/\widetilde{K}$ were creating/annihilation operators, we now want to see what these translate to. Let's consider the Virasoro descendant $L_n\ket{h}$ and check its weight: 
\bse 
    L_0 L_n \ket{h} = [L_0,L_n] \ket{h} + h L_n\ket{h} = (h-n)L_n\ket{h},
\ese 
where we have used $L_0\ket{h} = h\ket{h}$ and $[L_0,L_n] = -nL_n$. So we see that $L_n$ lowers the weight by $n$ when $n>0$ and raises the weight when $n<0$. Note that the difference to before is now we can lower a state by $n$ in \textit{one go}, whereas before we had to do $n$ steps. This is important because it tells us that there is more than one way to produce a descendent of a given weight. For example we can product a descendant of weight $h+2$ via 
\bse 
    L_{-2}\ket{h} \qquad \text{or} \qquad L_{-1}^2 \ket{h}.
\ese    

Let's have a look at the identity operator. This had no descendants in the higher dimensional case because it was annihilated by the raising operator $\widetilde{P}_{\mu}$. However in $2D$ it \textit{does} have descendants when $c\neq0$. Of course consistency with the higher dimensional cases tell us that the state dual to the identity, the absolute vacuum, is annihilated by the M\"{o}bius subgroup. So we have\footnote{Interestingly, it can be shown that the \textit{only} state that is annihilated by both $L_{-1}$ and $\bar{L}_{-1}$ is the state dual to the identity. This is not actually hard to see: simply recall that $L_{-1}=\p$ and $\bar{L}_{-1}=\bar{\p}$, and so essentially they shift the local operator away from the origin. Therefore if $L_{-1}$ and $\bar{L}_{-1}$ is to annihilate the state, it must be position independent. The only position independent local operator is the identity. This footnote is included to make the point that the only state that is annihilated by the full M\"{o}bius subgroup is the absolute vacuum.} 
\bse 
    L_n \ket{0} \qquad \forall n \geq -1.
\ese 
However now consider the action of $L_{-2}$. In particular consider acting with $L_{+2}$ on $L_{-2}\ket{0}$:
\bse 
    L_{+2}L_{-2}\ket{0} = [L_2,L_{-2}]\ket{0} = \bigg(4L_0 + \frac{c}{2}\bigg)\ket{0} = \frac{c}{2}\ket{0},
\ese
so we clearly need $L_{-2}\ket{0}\neq0$ when $c\neq 0$. Now let's see what unitarity can tell us. Using $L_n^{\dagger} = L_{-n}$, we have 
\bse 
    \bra{0} L_2 L_{-2}\ket{0} = \big|L_{-2}\ket{0}\big|^2 = \frac{c}{2}\braket{0}{0},
\ese 
so if we have a unitary theory (i.e. inner products are non-negative definite) we have 
\mybox{
    \be 
    \label{eqn:cPositive}
        c > 0. 
    \ee 
}

\badr 
    It is a fair point to say "actually all we can conclude $c$ is non-negative so \Cref{eqn:cPositive} should be $c \geq 0$". Well it turns out that when $c=0$ the only state in the whole theory is the absolute vacuum, which is boring. So \Cref{eqn:cPositive} is meant to be understood for any non-trivial theory. 
\eadr 

We can now look at building up our irreps. We start by defining the lowest weight state $\ket{h}$ such that 
\bse 
    L_n \ket{h} = 0 \qquad \forall n >0.
\ese 
We then obtain all other states in the representation by acting on this with raising operators $L_{-n}$. We list the first few descedents below 

\bse 
    \textcolor{blue}{\ket{h}}
\ese 
\bse 
    \textcolor{blue}{L_{-1}\ket{h}}
\ese 
\bse 
    \textcolor{blue}{L_{-1}^2\ket{h}} \quad \textcolor{red}{L_{-2}\ket{h}}
\ese 
\bse 
    \textcolor{blue}{L_{-1}^3\ket{h}} \quad \textcolor{red}{L_{-1}L_{-2}\ket{h}} \quad \textcolor{green}{L_{-3}\ket{h}}.
\ese 
\bse 
    \textcolor{blue}{L_{-1}^4\ket{h}} \quad \textcolor{red}{L_{-1}^2L_{-2}\ket{h}} \quad \textcolor{green}{L_{-1}L_{-3}\ket{h}} \quad \textcolor{purple}{L_{-2}^2\ket{h}} \quad L_{-4}\ket{h}
\ese 
The complete set of such states is known as the \textit{Verma module}. They form an irreducible representation of the Virasoro algebra. 

What is the colour coordination meant to mean? Well on a closer look we see that within each colour, the only thing that changes as we move down the rows is the action of a $L_{-1}$. These are therefore representations of the global M\"{o}bius subgroup. In other words the first state that appears in a given colour is a primary state through the lenses of the M\"{o}bius group, and then the states of the same colour are descendants. Note that $L_{-2}^2\ket{h}$ is \textit{not} a descendant of $L_{-2}\ket{h}$ from the M\"{o}bius subgroups lenses.

\br 
    Note that we don't need to include the state $L_{-2}L_{-1}\ket{h}$ on the third row (and similarly for lower rows). Why? Well because we can use the Virasoro algebra to relate this to $L_{-1}L_{-2}\ket{h}$ and $L_{-3}\ket{h}$, and so it is \textit{not} a new state. 
\er 

\badr 
    As an extension of the above remark, there is actually a subtly between the states of the Verma module, and it turns out that they are not all necessarily independent. We do not give an example here but the interested reader is directed to the bottom of page 97 of Prof. Tong's notes. 
\eadr 

\section{Another example of a CFT: bc Ghost Theory}

We end the course with a very quick overview of another example of a $2D$ CFT. This system plays a crucial role in one approach to string theory and it is known as the \textit{$bc$ ghost theory.} We do not present any proofs of the statements made here, but details will be flushed out in the string theory course.\footnote{Of course more details can also be found in either Prof. Tong's notes or my notes on Prof. Minwalla's course.}

We introduced the central charge $c$ above with absolutely no problems, and indeed it doesn't pose any problem for \textit{flat} $2D$ CFTs. However, as the emphasis indicates, this is not true for CFTs on \textit{curved} backgrounds. The reason is related to the comment we made about Weyl anomolies all the way back in \Cref{rem:WeylAnom}: the trace of the stress-energy tensor in $2D$ turns out to be given by
\be 
\label{eqn:WeylAnom}
    {T^a}_{a} = -\frac{c}{12} R,
\ee 
where $R$ is the Ricci scalar. As we said in the remark, this is known as a Weyl anomaly. 


\badr 
    It turns out that we can trace the presence of the central charge, $c$, back to the Weyl transformations in our 2D CFT. Essentially the idea is that the $(z-w)^{-4}$ term appearing in the $T(z)T(w)$ OPE comes purely from the Weyl part of a conformal transformation, and therefore $c$ comes from the Weyl transformations. Putting this together with the fact that \textit{classically} conformal invariance gave us the traceless condition, we see why \Cref{eqn:WeylAnom} is called a Weyl anomaly; we loose the exact classical Weyl invariance in the quantum theory. For more details flushing out this idea, see Lecture 7 of my notes on Prof. Shiraz Minwalla's string theory course.
\eadr 

Why is this a problem? Well recall that in our discussion of string theory we needed Weyl invariance in order to make the Polyakov and Nambu-Goto actions equivalent. That is if we don't impose Weyl invariance on the Polyakov action then our set of solutions is (infinitely) bigger. This Weyl symmetry also allowed us to restrict the theory to a flat metric and gave us our conformal symmetry. \Cref{eqn:WeylAnom} breaks our Weyl invariance and so ruins everything. There is a way we can save this, though: make $c=0$. 

However we already said that if $c=0$ then the only state in our theory is the absolute vacuum, so it appears we are doomed. However we now note that important word used above: in a \textit{unitary} theory we require $c>0$. That is a non-unitary theory can have negative central charge. From previous courses, we know of a field which we can add to our theory which is non-unitary: ghosts! The idea is to consider a theory which consists of both free Bosons, $X^{\mu}$, and ghosts, which we standardly denote $b$ and $c$, such that their \textit{collective} central charge vanishes. That is we want 
\bse 
    c_{\text{tot}} = c_{\text{Bosons}} + c_{\text{ghosts}} = 0.
\ese

\subsection{Ghost Theory}
The action for our ghost theory on our 2D space turns out to be
\bse 
    S_g = \frac{1}{2\pi} \int d^2 x \big( b\bar{\p} c + \bar{b}\p \bar{c}\big)
\ese 
where $b,c$ are our ghosts. A quick calculation shows that the equations of motion are
\bse 
    \bar{\p} b = \bar{\p} c = \p \bar{c} = \p \bar{b} = 0
\ese
which tells us that $b,c$ are holororphic and $\bar{b}, \bar{c}$ are antiholomorphic. We can show that 
\bse 
    b(z) = \sum_m \frac{b_m}{z^{m+2}} \qand c(z) = \sum_m \frac{c_m}{z^{m-1}} 
\ese 
so $b$ has weights $(2,0)$ and $c$ has weights $(-1,0)$. The anticommutators between the two ghost types turn out to be
\bse 
    \{ b_m, c_n\} = \del_{m+n,0} 
\ese
and we can show\footnote{We stress again here that, as the $bc$ system is non-unitary, the state dual to the identity need \textit{not} be the vacuum. Indeed it turns out that $\ket{\b1} = b_{-1}\ket{\downarrow}$, where $\ket{\downarrow}$ is one of the two ground states of the $bc$ system. A proof of all this can be found in lecture 11 of my notes on Prof. Shiraz Minwalla's string theory course.} 
\bse 
    \begin{split}
        b_m \ket{\b1} = 0 \qquad m \geq -1 \\
        c_n \ket{\b1} = 0 \qquad m \geq 2.
    \end{split}
\ese

As the ghost system has its own action, it also has it's own stress-energy tensor, which is easily verified to be 
\bse 
    T = 2(\p c) b + c \p b.
\ese 
This in turn tells us that 
\bse 
    L_n = \sum_m (2n-m)b_m c_{n-m}.
\ese 

Ok that was a \textit{lot} of plucked-out-of-the-air stuff, so why do we want it? Well we can now compute the inner product $\bra{0}L_2L_{-2}\ket{0}$ again. If we plug through all the algebra we obtain
\bse 
    \bra{\b1}L_2L_{-2} \ket{\b1} = -13\braket{\b1}{\b1},
\ese
which comparing to 
\bse 
    \bra{\b1}L_2L_{-2}\ket{\b1} = \frac{c}{2}\braket{\b1}{\b1}
\ese 
tells us
\bse 
    c_{\text{ghost}} = -26.
\ese 
We then need to compensate for this with $26$ Bosonic fields $X^{\mu}$, as each Bosonic field contributes $c_{\text{Boson}}=+1$. But the number of Bosonic fields is given by the dimension of the spacetime, and so we conlclude that 
\mybox{
    \bse 
        D=26.
    \ese 
}   
\noindent This is the famous result that string theory is a 26 dimensional theory. 

\br 
    There is actually another way to obtain the result $D=26$ without having to introduce ghosts. This involves breaking manifest Lorentz invariant and going to light-cone gauge. You then quantise the system there and look at the resulting irreducible representations and \textit{insist} that these form a valid decomposition of the Lorentz group. Details of this can be found in Prof Tong's notes (again or my notes on Prof Minwalla's notes).
\er 

\badr 
    For completeness, we just make one comment about the $D=10$ result that people also say. This corresponds to \textit{super}string theory. This corresponds to considering free Fermions, and so is sometimes also called \textit{Fermionic} string theory. In this case we end up getting $c_{\text{ghost}}=-15$ and each Fermion contributes $c_{\text{Fermion}} = 3/2$, so in total we need $10$ spacetime dimensions. Much more details on this can be found in my notes on Prof Minwalla's course.\footnote{Unfortunately, Prof. Tong's notes don't discuss the Fermionic string in analytic detail, so unless you wish to give Polchinski Volume II a bash, I'm afraid all I can reference are my notes.}
\eadr 

% -------------------------------------------------------------------
% Bibliography/Further Readings
% -------------------------------------------------------------------

\chapter*{Useful Texts \& Further Readings}

\section*{Part I: CFT in $D>2$}

\begin{itemize}
    \item Simon Ross's notes (on duo). Focused on $2d$ CFT, but used them if applicable.
    \item Simmons-Duffin, [1602.07982].
    \item Rychkov lectures [1601.05000], he is leading numerical conformal bootstrap research.
    \item Qualls lectures [1511.04074].
\end{itemize}

\section*{Second half}
\begin{itemize}
    \item Tong string notes (first half is CFT).
    \item Simmons-Duffin, [1602.07982].
    \item Osborn (Cambridge) lecture notes on CFT.
    \item Polchinski "string theory" vol. 1
    \item Big Yellow Book, CFT (Di Francesco) (focusing on $2D$)
\end{itemize}






%\bibliographystyle{agsm} 
%\bibliography{mybibliography} 
%\printbibliography[heading=bibintoc]


% -------------------------------------------------------------------
% Appendices
% -------------------------------------------------------------------

%\begin{appendices}
%\chapter{An Appendix of Some Kind}

\lipsum[1-3]  % Replace with your text

%\end{appendices}

\end{document}
